\pandocbounded{
\includegraphics[keepaspectratio]{https://raw.githubusercontent.com/JantN1966/MasterManual/main/Images/logo_MiXBLUP.jpg}}

\hfill\break

\section{Manual MiXBLUP}\label{manual-mixblup}

\hfill\break

\textbf{Version 3.2.1 -- 2026 -- 06} \newpage

MiXBLUP User's Guide\\
MiXBLUP, best linear unbiased prediction (BLUP) software for PCs for
large genetic evaluation systems

This manual is for MiXBLUP version 3.2, released in May 2026\\
MiXBLUP is developed jointly by LUKE National Resources Institute
Finland and Animal Breeding \& Genomics, Wageningen University \&
Research.

Authors:\\
J. ten Napel, J. Vandenplas, M. Schrauf, M.P.L. Calus, R.F. Veerkamp\\
M. Lidauer, I. Stranden, M. Taskinen, E. Mäntysaari

Animal Breeding \& Genomics, Wageningen University \& Research P.O. Box
338\\
6700 AH Wageningen\\
The Netherlands

LUKE National Resources Institute Finland\\
FI-31600\\
Jokioinen\\
Finland

More information on the \href{http://www.mixblup.eu/}{MiXBLUP website}\\
\strut \\

\newpage

\subsection{Table of Contents}\label{Tabl01}

\hyperref[Intr01]{1. Introduction}\\
\hyperref[Intr02]{1.1. Overview}\\
\hyperref[Intr03]{1.2. Manual}\\
\hyperref[Intr04]{1.3. System requirements}\\
\hyperref[HowT01]{2. How to start}\\
\hyperref[HowT02]{2.1. Installing MiXBLUP software}\\
\hyperref[HowT03]{2.2. MiXBLUP Licenses}\\
\hyperref[HowT04]{2.2.1. Trial License}\\
\hyperref[HowT05]{2.2.2. Commercial licenses}\\
\hyperref[HowT06]{2.2.3. Generating a license-request file and
installing the license}\\
\hyperref[HowT07]{2.2.4. Alternative license directory}\\
\hyperref[Inst01]{3. Instruction file}\\
\hyperref[Inst02]{3.1. Parts of the instruction file}\\
\hyperref[Inst03]{3.1.1. Title of the analysis}\\
\hyperref[Inst04]{3.1.2. Observations \& systematic effects}\\
\hyperref[Inst05]{3.1.3. Genetic similarity among individuals}\\
\hyperref[Inst06]{3.1.4. Components of variance and covariance among
traits}\\
\hyperref[Inst07]{3.1.5. Statistical models}\\
\hyperref[Inst08]{3.1.6. Control of analysis and output}\\
\hyperref[Inst09]{3.2. General syntax of the instruction file}\\
\hyperref[Obse01]{4. Observations \& systematic effects}\\
\hyperref[Obse02]{4.1. Data file}\\
\hyperref[Obse03]{4.1.1. General}\\
\hyperref[Obse04]{4.1.2. Input file}\\
\hyperref[Obse05]{4.1.3. Syntax}\\
\hyperref[Obse06]{4.1.4. Associated output files}\\
\hyperref[Obse07]{4.2. Covariate table file}\\
\hyperref[Obse08]{4.2.1. General}\\
\hyperref[Obse09]{4.2.2. Input file}\\
\hyperref[Obse10]{4.2.2.1. Syntax using an existing covariate table for
the MiX99 solver}\\
\hyperref[Obse11]{4.2.2.2. Syntax using an existing covariate table for
the hpblup solver}\\
\hyperref[Obse12]{4.2.2.3. Syntax using a newly created covariate table
for the MiX99 solver}\\
\hyperref[Obse13]{4.2.2.4. Syntax using newly created covariate tables
for the hpblup solver}\\
\hyperref[Obse15]{4.2.3. Associated output files}\\
\hyperref[Obse16]{4.3. General covariate files}\\
\hyperref[Obse17]{4.3.1. General}\\
\hyperref[Obse18]{4.3.2. Input file}\\
\hyperref[Obse19]{4.3.3. Syntax}\\
\hyperref[Obse20]{4.3.3.1. Syntax of a general covariate file and
associated variance-covariance file}\\
\hyperref[Obse23]{4.3.3.2. Syntax of fitting a general covariate file in
the model for the MiX99 solver}\\
\hyperref[Obse25]{4.3.3.3. Syntax of fitting a general covariate file in
the model for the hpblup solver}\\
\hyperref[Obse27]{4.3.4. Associated output files}\\
\hyperref[Obse28]{4.4. Random effects with correlated level effects}\\
\hyperref[Obse29]{4.4.1. General}\\
\hyperref[Obse30]{4.4.2. Input file}\\
\hyperref[Obse31]{4.5. Syntax}\\
\hyperref[Obse32]{4.5.1. Associated output files}\\
\hyperref[Gene01]{5. Genetic similarity among individuals}\\
\hyperref[Gene02]{5.1. Preparing pedigree data}\\
\hyperref[Gene03]{5.1.1. General}\\
\hyperref[Gene04]{5.1.2. Recommended file formats}\\
\hyperref[Gene05]{5.1.2.1. Pedigree file}\\
\hyperref[Gene06]{5.1.2.2. Pedigree inbreeding coefficient file}\\
\hyperref[Gene07]{5.1.3. Pedigree inbreeding coefficients}\\
\hyperref[Gene08]{5.1.3.1. General}\\
\hyperref[Gene09]{5.1.3.2. Syntax of calculating inbreeding
coefficients}\\
\hyperref[Gene10]{5.1.3.3. Syntax of using file with inbreeding
coefficients}\\
\hyperref[Gene12]{5.1.4. Unknown parents are from a single large base
population}\\
\hyperref[Gene13]{5.1.5. Unknown parents are from multiple large base
populations}\\
\hyperref[Gene14]{5.1.5.1. Syntax of multiple large base populations
using Westell grouping}\\
\hyperref[Gene15]{5.1.5.2. Associated output files for Westell
grouping}\\
\hyperref[Gene16]{5.1.5.3. Syntax of multiple large base populations
using precalculated genetic group covariates}\\
\hyperref[Gene99]{5.1.5.4. Syntax of multiple large base populations
using calculated genetic group covariates}\\
\hyperref[Gene17]{5.1.5.5. Associated output files for genetic group
covariates}\\
\hyperref[Gene18]{5.1.6. Unknown parents are from multiple related base
populations (metafounders)}\\
\hyperref[Gene19]{5.1.6.1. Syntax of multiple related base populations
using metafounders}\\
\hyperref[Gene20]{5.1.6.2. Associated output files for metafounders}\\
\hyperref[Gene21]{5.2. Preparing genomic data}\\
\hyperref[Gene22]{5.2.1. General}\\
\hyperref[Gene23]{5.2.2. Recommended file formats}\\
\hyperref[Gene24]{5.2.2.1. Genomic data}\\
\hyperref[Gene25]{5.2.2.2. Allele frequencies}\\
\hyperref[Gene26]{5.3. Genetic similarity from pedigree only}\\
\hyperref[Gene27]{5.3.1. General}\\
\hyperref[Gene28]{5.3.2. Syntax of using pedigree BLUP}\\
\hyperref[Gene29]{5.4. Genetic similarity from genomic data}\\
\hyperref[Gene30]{5.4.1. Genetic similarity from genomic data with
unknown pedigree}\\
\hyperref[Gene31]{5.4.1.1. General}\\
\hyperref[Gene32]{5.4.1.2. Syntax of using GBLUP}\\
\hyperref[Gene33]{5.4.1.3. Syntax of using SNPBLUP for the MiX99
solver}\\
\hyperref[Gene34]{5.4.1.4. Syntax of using SNPBLUP for hpblup solver}\\
\hyperref[Gene35]{5.4.2. Genetic similarity from pedigree and genomic
data}\\
\hyperref[Gene36]{5.4.2.1. General}\\
\hyperref[Gene37]{5.4.2.2. Syntax of SSGBLUP: single-step genomic BLUP
with full inverse of a weighted G}\\
\hyperref[Gene46]{5.4.2.3. Syntax of ssGTacBLUP: single-step GBLUP with
component-wise Ta decomposition of a weighted G}\\
\hyperref[Gene47]{5.4.2.4. Syntax of ssSNPBLUP: single-step BLUP using
SNP genotypes as covariates}\\
\hyperref[Gene48]{5.4.2.5. Syntax of DGV-Pedigree BLUP from Tac: using
previously estimated SNP effects as prior information in pedigree
BLUP}\\
\hyperref[Gene49]{5.4.2.6. Syntax of DGV-Pedigree BLUP from ssSNPBLUP:
using previously estimated SNP effects as prior information in pedigree
BLUP}\\
\hyperref[Gene50]{5.4.3. Modelling a genetic difference between
genotyped and non-genotyped individuals with the hpblup solver}\\
\hyperref[Gene51]{5.4.3.1. General}\\
\hyperref[Gene52]{5.4.3.2. Syntax of fitting J factor covariate}\\
\hyperref[Gene53]{5.4.4. Obtaining solutions for genetic marker effects
with the hpblup solver}\\
\hyperref[Gene54]{5.5. External genetic relationship matrix}\\
\hyperref[Gene55]{5.5.1. General}\\
\hyperref[Gene56]{5.5.2. Recommended file format}\\
\hyperref[Gene57]{5.5.3. Syntax of using an external relationship
matrix}\\
\hyperref[Gene58]{5.6. Genetic similarity in case of multiple breeds or
lines and crosses}\\
\hyperref[Gene59]{5.6.1. General}\\
\hyperref[Gene60]{5.6.2. Fixed effect in the model}\\
\hyperref[Gene61]{5.6.3. Base population for each genetic line}\\
\hyperref[Gene62]{5.6.4. Allele frequencies specific for a genetic
line}\\
\hyperref[Gene63]{5.7. Non-additive genetic similarity}\\
\hyperref[Gene64]{5.7.1. General}\\
\hyperref[Gene65]{5.7.2. Expected heterosis and recombination in
crossbreds}\\
\hyperref[Gene66]{5.7.3. Genomic dominance effects (to complete)}\\
\hyperref[Gene67]{5.7.3.1. General}\\
\hyperref[Gene68]{5.7.3.2. Syntax}\\
\hyperref[Gene69]{5.7.3.3. Output files}\\
\hyperref[Gene70]{5.7.4. Genomic epistasis effects (to complete)}\\
\hyperref[Gene71]{5.7.4.1. General}\\
\hyperref[Gene72]{5.7.4.2. Syntax}\\
\hyperref[Gene73]{5.7.4.3. Output files}\\
\hyperref[Gene74]{5.8. Genetic similarity in case of polyploidy or mixed
ploidy}\\
\hyperref[Gene75]{5.8.1. General}\\
\hyperref[Gene76]{5.8.2. Required format of pedigree file}\\
\hyperref[Gene77]{5.8.3. Syntax}\\
\hyperref[Comp01]{6. Components of variance and covariance among
traits}\\
\hyperref[Comp02]{6.1. General parameter file}\\
\hyperref[Comp03]{6.1.1. General}\\
\hyperref[Comp04]{6.1.2. Input file in lower-triangular-matrix format}\\
\hyperref[Comp05]{6.1.3. Input file in sparse-matrix format for the
MiX99 solver}\\
\hyperref[Comp06]{6.1.4. Syntax}\\
\hyperref[Comp07]{6.2. Parameter files for general covariates}\\
\hyperref[Comp08]{6.2.1. General}\\
\hyperref[Comp09]{6.2.2. Input file}\\
\hyperref[Comp10]{6.2.3. Syntax}\\
\hyperref[Comp11]{6.3. Parameters for SNP covariate files}\\
\hyperref[Comp12]{6.3.1. General}\\
\hyperref[Comp13]{6.3.2. Input file}\\
\hyperref[Comp14]{6.3.3. Syntax}\\
\hyperref[Comp15]{6.4. Parameters in case of heterogeneous residual
variances}\\
\hyperref[Comp16]{6.4.1. General}\\
\hyperref[Comp17]{6.4.2. Input file}\\
\hyperref[Comp18]{6.4.3. Syntax}\\
\hyperref[Stat01]{7. Statistical models}\\
\hyperref[Stat02]{7.1. Basic models}\\
\hyperref[Stat03]{7.1.1. General}\\
\hyperref[Stat04]{7.1.2. Syntax}\\
\hyperref[Stat05]{7.1.3. Associated output files}\\
\hyperref[Stat06]{7.2. Repeatability models}\\
\hyperref[Stat07]{7.2.1. General}\\
\hyperref[Stat08]{7.2.2. Syntax}\\
\hyperref[Stat09]{7.2.3. Associated output files}\\
\hyperref[Stat10]{7.3. Maternal genetic models}\\
\hyperref[Stat11]{7.3.1. General}\\
\hyperref[Stat12]{7.3.2. Syntax}\\
\hyperref[Stat13]{7.3.3. Associated output files}\\
\hyperref[Stat14]{7.4. Social interaction models}\\
\hyperref[Stat15]{7.4.1. General}\\
\hyperref[Stat16]{7.4.2. Syntax of the social interaction model with one
group size for all groups for the MiX99 solver}\\
\hyperref[Stat17]{7.4.3. Syntax of the social interaction model with
slightly varying group sizes for the MiX99 solver}\\
\hyperref[Stat18]{7.4.4. Syntax of the social interaction model for the
hpblup solver}\\
\hyperref[Stat19]{7.4.5. Associated output files}\\
\hyperref[Stat20]{7.5. Random regression models}\\
\hyperref[Stat21]{7.5.1. General}\\
\hyperref[Stat22]{7.5.2. Syntax of a simple non-genetic random
regression model for the MiX99 solver}\\
\hyperref[Stat23]{7.5.3. Syntax of a simple genetic random regression
model for the MiX99 solver}\\
\hyperref[Stat24]{7.5.4. Syntax of a polynomial regression model using a
covariate table for the MiX99 solver}\\
\hyperref[Stat25]{7.5.5. Syntax of a simple non-genetic random
regression model for the hpblup solver}\\
\hyperref[Stat26]{7.5.6. Syntax of a simple genetic random regression
model for the hpblup solver}\\
\hyperref[Stat27]{7.5.7. Syntax of a polynomial regression model using a
covariate table for the hpblup solver}\\
\hyperref[Stat28]{7.5.8. Associated output files}\\
\hyperref[Stat29]{7.6. Weighting residuals by record}\\
\hyperref[Stat30]{7.6.1. General}\\
\hyperref[Stat31]{7.6.2. Syntax}\\
\hyperref[Stat32]{7.6.3. Associated output files}\\
\hyperref[Stat33]{7.7. Combining effects across traits}\\
\hyperref[Stat34]{7.7.1. General}\\
\hyperref[Stat35]{7.7.2. Syntax}\\
\hyperref[Stat36]{7.7.3. Associated output files}\\
\hyperref[Stat37]{7.8. Correction of heterogeneous residual variances}\\
\hyperref[Stat38]{7.8.1. General}\\
\hyperref[Stat39]{7.8.2. Syntax}\\
\hyperref[Stat40]{7.8.3. Associated output files}\\
\hyperref[Stat41]{7.9. Using a threshold model for a categorical trait
(MiX99 solver only)}\\
\hyperref[Stat42]{7.9.1. General}\\
\hyperref[Stat43]{7.9.2. Input files}\\
\hyperref[Stat44]{7.9.3. Syntax}\\
\hyperref[Stat45]{7.9.4. Associated files}\\
\hyperref[Cont01]{8. Control of analysis and output}\\
\hyperref[Cont02]{8.1. Control of the analysis}\\
\hyperref[Cont03]{8.1.1. General}\\
\hyperref[Cont04]{8.1.2. Syntax}\\
\hyperref[Cont05]{8.1.2.1. Syntax when using MiX99 solver}\\
\hyperref[Cont06]{8.1.2.2. Syntax when using hpblup solver}\\
\hyperref[Cont07]{8.2. Control of output}\\
\hyperref[Cont08]{8.2.1. General}\\
\hyperref[Cont09]{8.2.2. Syntax}\\
\hyperref[Cont10]{8.2.2.1. Syntax when using MiX99 solver}\\
\hyperref[Cont11]{8.2.2.2. Syntax when using hpblup solver}\\
\hyperref[Reli01]{9. Reliabilities}\\
\hyperref[Reli02]{9.1. General}\\
\hyperref[Reli03]{9.2. Exact reliabilities}\\
\hyperref[Reli05]{9.2.1. Syntax}\\
\hyperref[Reli07]{9.3. Approximate reliabilities}\\
\hyperref[Reli08]{9.3.1. General}\\
\hyperref[Reli10]{9.3.1.1. Block variable}\\
\hyperref[Reli11]{9.3.1.2. Common-block variable}\\
\hyperref[Reli12]{9.3.1.3. Sorting data and pedigree file on block
variable}\\
\hyperref[Reli13]{9.3.1.4. Strategies for block definition}\\
\hyperref[Reli14]{9.3.2. Differences between the syntax of the
instruction file for approximate reliability calculation and breeding
value estimation}\\
\hyperref[Reli15]{9.3.2.1. Data file}\\
\hyperref[Reli16]{9.3.2.2. Genetic similarity between individuals}\\
\hyperref[Reli17]{9.3.2.3. Statistical model}\\
\hyperref[Reli18]{9.3.2.4. Control of analysis}\\
\hyperref[Reli21]{9.4. Approximate genomic reliabilities}\\
\hyperref[Reli26]{9.5. Command-line interface for calculating
reliabilities}\\
\hyperref[MTGW01]{10. Multi-trait genome-wide association studies
(GWAS)}\\
\hyperref[MTGW02]{10.1. General}\\
\hyperref[MTGW03]{10.2. Computation of frequentist p-values for limited
datasets}\\
\hyperref[MTGW04]{10.2.1. General}\\
\hyperref[MTGW05]{10.2.2. Syntax for calculating frequentist p-values}\\
\hyperref[MTGW06]{10.2.3. Associated output files}\\
\hyperref[MTGW07]{10.2.4. Example (move to appendix later)}\\
\hyperref[MTGW08]{10.3. Approximation of frequentist p-values for
large-scale datasets}\\
\hyperref[MTGW09]{10.3.1. General}\\
\hyperref[MTGW10]{10.3.2. Syntax for calculating frequentist p-values}\\
\hyperref[MTGW11]{10.3.3. Associated output files}\\
\hyperref[MTGW12]{10.3.4. Example (move to appendix later)}\\
\hyperref[Desc01]{11. Descriptive analyses}\\
\hyperref[Desc02]{11.1. General}\\
\hyperref[Desc03]{11.2. Descriptive statistics of performance data,
genomic data and pedigree}\\
\hyperref[Desc04]{11.2.1. General}\\
\hyperref[Desc05]{11.2.2. Syntax}\\
\hyperref[Desc06]{11.2.3. Counts (option N)}\\
\hyperref[Desc07]{11.2.4. Means and standard deviations (option D)}\\
\hyperref[Desc08]{11.2.5. Class effect levels grouped by available
information (option H)}\\
\hyperref[Desc09]{11.2.6. Valid data records for each combination of
trait and class effect level (option L)}\\
\hyperref[Desc10]{11.2.7. Associated output files}\\
\hyperref[Desc11]{11.3. Diagnostics of use of base populations}\\
\hyperref[Desc12]{11.3.1. General}\\
\hyperref[Desc13]{11.3.2. Equivalent number of base animals genotyped}\\
\hyperref[Desc14]{11.3.3. Auto-similarity to another base population}\\
\hyperref[Desc15]{11.3.4. Number of generations between pedigree and
genomic base populations}\\
\hyperref[Desc16]{11.3.5. Syntax}\\
\hyperref[Desc17]{11.3.6. Associated output files}\\
\hyperref[Vali01]{12. Validation studies with MiXBLUP}\\
\hyperref[Vali03]{12.1. Validation individuals}\\
\hyperref[Vali04]{12.2. Validation effect}\\
\hyperref[Vali05]{12.3. Creating partial dataset}\\
\hyperref[Vali06]{12.4. Types of validation}\\
\hyperref[Vali07]{12.5. Command-line syntax}\\
\hyperref[Vali08]{12.6. Syntax}\\
\hyperref[Vali09]{12.7. Validation statistics}\\
\hyperref[Vali10]{12.8. Associated output files}\\
\hyperref[Indi01]{13. Indirect prediction}\\
\hyperref[Indi02]{13.1. General}\\
\hyperref[Indi03]{13.2. Indirectly predicting genomic estimated breeding
values}\\
\hyperref[Indi04]{13.2.1. General}\\
\hyperref[Indi05]{13.2.2. Supported evaluations}\\
\hyperref[Indi06]{13.2.3. Syntax}\\
\hyperref[Indi07]{13.2.4. Output files}\\
\hyperref[Runn01]{14. Running MiXBLUP}\\
\hyperref[Runn02]{14.1. Starting a MiXBLUP evaluation}\\
\hyperref[Runn03]{14.2. Choosing a breeding value evaluation or a
reliability calculation}\\
\hyperref[Runn04]{14.3. A breeding value analysis with previous
solutions as starting values}\\
\hyperref[Runn05]{14.4. Monitoring and checking the process}\\
\hyperref[Runn06]{14.5. Interrupting a process of the kernel}\\
\hyperref[Deco01]{15. Decoding any file with coded class effect
labels}\\
\hyperref[Deco02]{15.1. General}\\
\hyperref[Deco03]{15.2. Decoding coded labels}\\
\hyperref[Deco04]{15.2.1. General}\\
\hyperref[Deco05]{15.2.2. Syntax}\\
\hyperref[Deco06]{15.2.3. Decoding individual coded class effect
labels}\\
\hyperref[Deco07]{15.2.4. Decoding a file that contains coded class
effect labels}\\

\newpage

\subsection{1. Introduction}\label{Intr01}

MiXBLUP has been developed for routine breeding value estimation in
commercial genetic programmes and supports modern applications, such as
random regression models, group selection, the use of genetic markers or
haplotypes and the use of genomic information.

\subsubsection{1.1. Overview}\label{Intr02}

The intention of developing MiXBLUP was to utilize efficient computing
strategies for solving mixed model equations. With MiXBLUP it is
possible to use sophisticated models in estimation of breeding values in
animals, like cattle, pigs, poultry, sheep, horses, goats, dogs, and
aquatic species. The software also supports many ways to specify genetic
similarity between individuals, including pedigree, marker information
and genomic information. The statistical method used for genetic
evaluation is best linear unbiased prediction (BLUP), which is currently
the common methodology for genetic evaluation.\\
MiXBLUP supports two solvers. The MiX99 solver has been developed for
efficient use of disk space and memory. Due to iteration on data and a
very fast algorithm in the solver (preconditioned conjugate gradient,
PCG), it is able to solve mixed model equations very fast. It is derived
from MiX99 and was initially developed for classical genetic evaluation
without the use of markers or genes by LUKE National resources Institute
Finland. The adaptation for the use of marker and genomic information
was implemented by Wageningen UR Livestock Research in collaboration
with LUKE.\\
The hpblup solver has been developed specifically for efficient genetic
evaluation using a very large amount of genomic information. It is also
based on a PCG algorithm, but genomic information is stored in memory
during solving and it uses multiple cores whenever beneficial.

\subsubsection{1.2. Manual}\label{Intr03}

This manual will guide the user through the use of MiXBLUP. The examples
provide a way to test the software, to get a feel for the software. A
set of examples is provided as an Appendix to the manual. The number of
the example refers to the corresponding chapter of this manual.\\
A schematic overview of the input files, output files and instruction
file is in Figure 1.\\

\pandocbounded{\includegraphics[keepaspectratio]{https://raw.githubusercontent.com/JantN1966/MasterManual/main/Images/Introd01_MiX.jpg}}

\emph{Figure 1.} Schematic overview of the input and output files of
MiXBLUP.\\

\subsubsection{1.3. System requirements}\label{Intr04}

MiXBLUP is written in standard Fortran 90 language and is
self-contained. The program runs in Windows, Linux and Unix environments
and is available in 64-bit version. In Windows, it runs in the
command-line interpreter, cmd.exe or in the Windows PowerShell. The
Windows release it is routinely tested in a Windows 11 operating
system.\\
MiXBLUP allocates memory depending on the need. Small applications can
be run with a minimum of memory available. Very large applications may
need a substantial amount of memory, especially genomic analyses and the
calculation of reliabilities.\\
MiXBLUP supports the use of multiple cores. The MiX99 solver uses all
available cores for the most common genomic evaluations, only. The
hpblup solver is optimised for 10-15 cores for all available types of
evaluation. Preparation of data for solving and processing its results
are done with a single core.

\hyperref[Tabl01]{Back to Table of Contents}

\newpage

\subsection{2. How to start}\label{HowT01}

MiXBLUP is easy to use and easy to install. This chapter describes how
to install the software and how to obtain and install a license.

\subsubsection{2.1. Installing MiXBLUP software}\label{HowT02}

Download the appropriate zip-file from http://www.MiXBLUP.eu and unzip
the folder with the executables. Copy the executables to a central
folder that can be accessed from other folders. The user needs to create
a file, named `SysDir.inp', which contains the path to the central
folder with executables. This file should be copied to any folder from
which MiXBLUP is run. The path to MiXBLUP.exe should be included in the
command file that starts up the analysis or added to the system path.
MiXBLUP uses SysDir.inp to locate the other executables.

\subsubsection{2.2. MiXBLUP Licenses}\label{HowT03}

To run MiXBLUP software on your computer you need a license. There are
different license types for MiXBLUP. A license can be ordered at
http://www.MiXBLUP.eu.

The license key of the commercial licenses is computer-specific.
Therefore, if executables and the license key `LICENSE.DAT' are moved to
another computer, MiXBLUP will give an error message. Running MiXBLUP
with the run-time option --D l (minus, uppercase D, lowercase L) writes
the host name, license type and expiry date in the license file to the
screen output. So if you want to transfer the MiXBLUP software with an
existing license to a new computer, you have to request a new license
from info@mixblup.eu with the LICREQST.DAT attached (how to generate a
LICREQST.DAT file see below). You will receive a new license for the
remainder of the license period. The license key provides the
information about the MiXBLUP version, the license type and the expiry
date of the license. A trial license can be used for one month and a
trial license key is not computer-specific. The small and full
commercial license can be used for one year. The license key for these
licenses is computer-specific.

\paragraph{2.2.1. Trial License}\label{HowT04}

Order a trial license at http://www.MiXBLUP.eu. After receiving your
order, we send the necessary license key to the e-mail address in the
order.

\paragraph{2.2.2. Commercial licenses}\label{HowT05}

Order a commercial license at http://www.MiXBLUP.eu. While entering the
order you are asked to upload one or more `LICREQST.DAT' files. For each
computer you need to upload a separate `LICREQST.DAT' file. This file is
required to generate a license key for your computer. Also renewing a
license for the next calendar year you need to do by filling in the
MiXBLUP License Order \& Renewal Form on the website.

\paragraph{2.2.3. Generating a license-request file and installing the
license}\label{HowT06}

The name of the license request file is `LICREQST.DAT'. The name of a
license file is `LICENSE.DAT'.

\begin{itemize}
\tightlist
\item
  Run MiXBLUP.exe once without the need for an instruction file. MiXBLUP
  creates the file LICREQST.DAT in the working directory.\\
\item
  After payment of the license one or more `LICENSE.DAT' files will be
  sent back and should be saved in the bin folder of the corresponding
  computer(s).\\
\item
  Store the license key `LICENSE.DAT' in the
  C:\textbackslash MiXBLUP\textbackslash bin-folder for Windows or in
  the /usr/bin-folder for Linux.
\end{itemize}

\paragraph{2.2.4. Alternative license directory}\label{HowT07}

If the license key cannot be stored in the default directory, the user
may create a file, named LicDir.inp, which contains the path to the
license file. If this file exists, MiXBLUP will look for the license
file in the specified folder.

\hyperref[Tabl01]{Back to Table of Contents}

\newpage

\subsection{3. Instruction file}\label{Inst01}

The instruction file contains all information that MiXBLUP needs for the
analysis. This chapter gives an overview of the instruction file. The
various parts of the instruction file are discussed in detail in the
next chapters.

\subsubsection{3.1. Parts of the instruction file}\label{Inst02}

The information in the MiXBLUP instruction file is presented in six
parts. These parts are:

\begin{enumerate}
\def\labelenumi{\arabic{enumi}.}
\tightlist
\item
  Description of the analysis\\
\item
  Observations \& systematic effects (Chapter \hyperref[Obse01]{4.})\\
\item
  Genetic similarity among individuals (Chapter \hyperref[Gene01]{5.})\\
\item
  Components of variance and covariance among traits (Chapter
  \hyperref[Comp01]{6.})\\
\item
  Statistical models (Chapter \hyperref[Stat01]{7.})\\
\item
  Control of analysis and output (Chapter \hyperref[Cont01]{8.})
\end{enumerate}

These parts may be presented in the instruction file in any order.
Sections within a part may also appear in any order. Below the example
instruction file is given for a bivariate animal model for two traits
(phen1 and phen2).

\begin{quote}
TITLE breeding value estimation for phen1 and phen2 using pedigree\\
\textbf{\# Observations \& systematic effects}\\
DATAFILE ExampleDat.txt !MISSING -99\\
animal A\\
fix1 A\\
fix2 I\\
cov R\\
ran A\\
phen1 T\\
phen2 T\\
blk I\\
\textbf{\# Genetic similarity among individuals}\\
PEDFILE ExamplePed.txt\\
animal A\\
sire A\\
dam A\\
blkped I\\
\textbf{\# Components of variance and covariance among traits}\\
PARFILE ExamplePar.dat\\
\textbf{\# Statistical models}\\
MODEL\\
phen1 \textasciitilde{} fix1 cov !RANDOM ran G(animal)\\
phen2 \textasciitilde{} fix2 cov !RANDOM ran G(animal)\\
\textbf{\# Control of analysis and output}\\
SOLVING\\
!MAXIT 1000\\
\end{quote}

\emph{Example}. Parts of the instruction file.

\paragraph{3.1.1. Title of the analysis}\label{Inst03}

The instruction file must start with a specification of the title of the
analysis. The TITLE keyword is optional. If omitted, the first line must
start with a hash (\#). This comment line is then used as the title of
the analysis. This line can be used to describe the analysis and
distinguish it from other analyses.

\paragraph{3.1.2. Observations \& systematic effects}\label{Inst04}

The data observations part of the instruction file contains the name of
the files with data or covariates, their location and their record
layout. The sections that can be used in this part are DATAFILE,
CVRTABLE and REGFILE. The syntax of these sections, more advanced
options and examples are presented in Chapter \hyperref[Obse01]{4.} of
this manual.

\paragraph{3.1.3. Genetic similarity among individuals}\label{Inst05}

Genetic similarity among individuals can be specified in many different
ways. It may be based on pedigree information only, genomic information
only or both sources of information simultaneously. Pedigree information
may contain genetic groups for unknown parents or a single code to
denote an unknown parent. Inbreeding can be taken into account or
ignored. Genomic information may be incorporated through covariances
between individuals or through regression on SNP covariates. Sections
that can be used in this part are PEDFILE, ERMFILE, INBRFILE, SNPFILE,
REGFILE, CORRFILE and CVMATRIX. The syntax of these sections and
examples for the various options are presented in Chapter
\hyperref[Gene01]{5.} of this manual.

\paragraph{3.1.4. Components of variance and covariance among
traits}\label{Inst06}

Genetic and non-genetic random effects have components of variance and
covariance among traits in the model. Residual (co)variance components
may also vary between groups of data records. Section that can be used
in this part are PARFILE, RESFILE, SNPPARFILE and REGPARFILE. The syntax
of these sections is presented in Chapter \hyperref[Comp01]{6.}

\paragraph{3.1.5. Statistical models}\label{Inst07}

Statistical models are specified by trait. Each trait starts on a new
line. The only sections in this part of the instruction file are MODEL,
LINKEDEFFECTS and COMBINE. The syntax of the various statistical models
supported by MiXBLUP are presented in Chapter \hyperref[Stat01]{7.}

\paragraph{3.1.6. Control of analysis and output}\label{Inst08}

The control part of the instruction file can be used to specify (1)
whether to solve the system (i.e.~estimate breeding values) or calculate
approximate reliabilities, (2) whether or not to use starting values,
(3) which resources to use for parts of the process, (4) when to stop
the iterative process and write out the solutions, (5) how to present
the solutions, (6) which additional output files to create after the
solving process has been completed and (7) how to manage temporary
files. The sections that can be used in this part are SOLVING, TRAITEBV,
PRECON and TMPDIR. Syntax is presented in Chapter \hyperref[Cont01]{8.}

\subsubsection{3.2. General syntax of the instruction
file}\label{Inst09}

\begin{itemize}
\tightlist
\item
  The maximum record length of the instruction file is 5,000 characters
\item
  The instruction file may contain empty lines for the convenience of
  the user
\item
  Comments may be inserted on a new line or after instructions on the
  same line, provided that any comment starts with a hash (\#). Any text
  on a line following a hash is ignored by MiXBLUP
\item
  The keyword of any section must be the first word of the line
\end{itemize}

\hyperref[Tabl01]{Back to Table of Contents}

\newpage

\subsection{4. Observations \& systematic effects}\label{Obse01}

\textbf{The data observations part of the instruction file is used to
specify observed traits and any factors or covariates that cause
systematic variation between observations for these traits. This chapter
describes the various ways to present observations and systematic
effects.}

\subsubsection{4.1. Data file}\label{Obse02}

\paragraph{4.1.1. General}\label{Obse03}

Observations and systematic effects are normally presented in the data
file. All traits and effects in the statistical model must have a column
in the data file, except for covariates in a covariate table file and
covariates in an external covariate file. The name of the data file is
specified in the instruction file. The data file is located by default
in the work directory, but it can be in any other folder if this is
specified as part of the name of the file
(e.g.~d:\textbackslash PerformanceTest\textbackslash BreedP.txt). The
order of the fields in the DATAFILE section must be the same as the
order of the fields in the data file.

\paragraph{4.1.2. Input file}\label{Obse04}

The data fields (individuals, systematic effects and trait observations)
each have their own column in the data file. The data file must be
provided in space-separated format, which means that any two columns are
separated by at least one space. Data fields can be integer values or
alphanumeric labels for class effects or real values for covariates and
trait observations. Real values are read with a decimal point.

Details of the layout of the data file:

\begin{itemize}
\item
  The maximum column width in the data file is 25 characters.
\item
  The maximum record length of the data file is 5,000 characters.
\item
  When data is alphanumeric, any of the symbols on the keyboard can be
  used, including a slash (`/').
\item
  An alphanumeric string must not contain spaces or it will be
  interpreted as two strings.
\item
  For the MiX99 solver, a class effect must not be zero or negative if
  it is a number, regardless of whether it is declared as integer or
  alphanumerical. Data records with a class effect in the model that is
  zero are omitted from the analysis by the kernel as invalid data
  points. Therefore, MiXBLUP replaces any classes of zero or a negative
  number with a 1 for an integer class effect or ``NA'' for an
  alphanumerical class effect. This will not affect the results of the
  evaluation if the invalid classes are not associated with a valid
  trait observation. It is left to the user to verify that this is
  indeed the case.
\item
  The default missing-value indicator for traits and covariates is zero.
  Data records with a covariate in the model that is equal to the
  missing-value indicator are omitted from the analysis by the kernel.
  If zero is a valid level for one of the covariates in the model,
  another missing-value indicator should be used. The missing value
  indicator has to be numerical.
\end{itemize}

\pandocbounded{\includegraphics[keepaspectratio]{https://raw.githubusercontent.com/JantN1966/MasterManual/main/Images/Observ01.jpg}}\\

\emph{Example}. Columns in data file: animal ID, mean, herd, sex, dam
ID, haplotype 1, haplotype 2, common environment, pen mate 1, pen mate
2, age 1, age 2, genotype, body weight at age 1, bodyweight at age 2.

\paragraph{4.1.3. Syntax}\label{Obse05}

\begin{quote}
DATAFILE \textless filename\textgreater{} {[}!SKIP \textless n
lines\textgreater{]} {[}!MISSING \textless value\textgreater{]}
{[}!SLASH{]} {[}!STATS {[}N{]}{[}D{]}{[}H{]}{[}L{]}{]} !MINMAX
\textless filename\textgreater{}\\
\textless field 1\textgreater{} \textless field type:
I/R/T/A\textgreater{}\\
\ldots{}\\
{[}\textless field i\textgreater{} I !BLOCK{]}\\
{[}\textless field j\textgreater{} I !RESVARCLASS{]}\\
\ldots{}\\
{[}\textless field n\textgreater{]} {[}I/R/T/A{]}
\end{quote}

Section:

\textbf{DATAFILE}\\
The DATAFILE section contains all the details of the file with trait
observations and systematic effects.

Qualifiers:

\textbf{!MISSING \textless value\textgreater{}}\\
If the value specified for !MISSING is encountered when reading traits
or covariates in the data file, it is interpreted as a missing
observation for that trait or covariate. A missing covariate invalidates
the trait for which the covariate is included in the model.

\textbf{!BLOCK}\\
This field is used as the block variable. If used, the data file and
pedigree file both need to contain this column. It is required for the
calculation of reliabilities, but might be beneficial in some
computationally heavy genetic evaluations. The field must be integer.
The !block qualifier must not be specified in the PEDFILE section, but
the fourth column in the pedigree file must have the same field name as
the block variable in the data file.

\textbf{!RESVARCLASS}\\
This field is used to specify the residual variance class of the data
record, in case the residual variance differs for groups of records. The
field must be integer. The qualifier !RESVARCLASS must be used if the
section RESFILE is specified.

\textbf{!SKIP \textless value\textgreater{}}\\
With this qualifier, one (!SKIP 1) or more (e.g.~!SKIP 2) header lines
in a data file can be ignored when reading the data file.

\textbf{!SLASH}\\
The qualifier !SLASH is optional and is used when any of the input files
contains a forward slash (`/') as a character. A forward slash is also a
control character in certain file formats. If !SLASH is not specified,
but MiXBLUP encounters a record with a forward slash, it will re-start
reading the file as if !SLASH had been specified. Reading of data with
!SLASH specified is slower than normal reading of data.

\textbf{!STATS NDHL}\\
The qualifier !STATS can be used to obtain a summary of descriptive
statistics of files in the evaluation, written to Statistics.log. There
are four types of statistics that can be produced: N for numbers of
records in data, pedigree and genotype file; D for means and standard
deviations of traits and covariates in the data; H for grouping class
effect levels for each trait by the number of records per class and L
for a table by trait with number of records for each class effect level.
For large evaluations, it is recommended to use !STATS NDH, as the
option L might produce a very large output file. Types may be specified
in any order. If D, H or L are specified, N is automatically included.

\textbf{!MINMAX \textless filename\textgreater{}}\\
The qualifier !MINMAX can be used to specify a file with the valid
ranges of traits and covariates. The file contains three fields for each
record: the name of the field in the data file (case-sensitive!), the
minimum and the maximum valid value. Field records may be in any order
and may contain field records of other data files, like the parameter
file.

More details of the syntax of the DATAFILE section:

\begin{itemize}
\item
  The field specification must start on the line following the line
  containing the DATAFILE keyword
\item
  The field type indicates whether a field in the data file should be
  read as an integer value (I), a real value for covariates (R), a real
  value for a trait (T) or a text string (A).
\item
  Maximum length of field names is 8 characters. A field name may be up
  to 19 characters long, but only the first 8 characters are used to
  distinguish fields, so a warning is given to remind the user.
\item
  Field names longer than 19 characters result in an error.
\item
  Field names are case-sensitive.
\item
  The qualifier !BLOCK only affects the MiX99 solver. If it is specified
  for multiple data fields, only the first specification is used.
\item
  Alphanumerical labels of a class effect (fields coded with A) are
  converted into integer values for the analysis. Solutions are decoded
  back to the original alphanumerical labels of the effect.
\item
  Each alphanumerical label in a field in the data file gets a unique
  numerical value. There is no apparent relation between the
  alphanumerical label and numerical value, so the numerical value of a
  string may vary across runs without using old solutions as starting
  values. The numerical value of a string does not change if old
  solutions are used as starting values by specifying !RESTART in the
  SOLVING section.
\item
  When using the hpblup solver, there is effectively no difference
  between field types A and I, as both types will be treated as
  alphanumeric.
\item
  The ID of animal in the data file, and the IDs of animal, its sire and
  its dam in the pedigree file must all be of the same type, so either
  alphanumeric (A) or numeric (I).
\item
  The largest integer number that can be used as level of a class effect
  is approximately 2,100,000,000 (2\textsuperscript{31}). For class
  effects with levels that exceed this number, the field type has to be
  set to alphanumerical (A).
\item
  The version of the data file with alphanumerical labels converted to
  integer values is `data.txt' for the MiX99 solver and hpData.txt for
  the hpblup solver.
\item
  The use of names reserved as section keywords, qualifiers or functions
  as field names is not supported.
\end{itemize}

\paragraph{4.1.4. Associated output files}\label{Obse06}

{\def\LTcaptype{none} % do not increment counter
\begin{longtable}[]{@{}ll@{}}
\toprule\noalign{}
Output file & Description \\
\midrule\noalign{}
\endhead
\bottomrule\noalign{}
\endlastfoot
data.txt & temporary file; data file prepared for analysis by kernel \\
\end{longtable}
}

\subsubsection{4.2. Covariate table file}\label{Obse07}

\paragraph{4.2.1. General}\label{Obse08}

If the relationship between an independent variable and a dependent
trait is modelled as an n\textsuperscript{th} order polynomial, a
covariate table file with all levels of the independent variable between
its minimum and maximum value in the data and (n+1) columns of
covariates may be used for easy presentation of covariates and syntax of
the instruction file. The name of a pre-defined covariate table file is
specified in the instruction file. The name may include the path to the
covariate table file. A covariate table file can also be created in
MiXBLUP. Currently only a Legendre polynomial is supported. A covariate
table is created using the minimum and maximum value of the independent
variable and the required order of the polynomial. The minimum and
maximum value of the independent variable can either be specified by the
user or determined from the data. For the MiX99 solver, only one
covariate table can be used, but its columns may be fitted within
multiple class effects. Additional polynomials using other independent
variables should be added as columns in the data file prior to calling
MiXBLUP. For the hpblup solver, it is possible to use multiple covariate
table files and indices to link covariate records to data records may be
different between covariate tables.

\paragraph{4.2.2. Input file}\label{Obse09}

A covariate table file may be created outside of MiXBLUP, it may have
been created in a previous analysis or it may be created at run-time. It
consists of the original independent variable and the n+1 covariates
derived from it, with n being the order of the polynomial. If the order
is n, the covariate columns in the table are numbered from 0 to n,
giving n+1 covariate columns in addition to the original independent
variable. The independent variable has to have an integer field type.
The covariate table should contain all levels between the minimum and
maximum value with steps of one. It means that an independent variable
with decimals must be converted to integer values before a covariate
table can be used for it. The independent variable links the record in
the data file with the covariate record in the covariate table. The
column in the data file with the independent variable must contain a
valid entry for every record. For the hpblup solver, each covariate
table must have a unique label that starts with TABLE followed by a
number between 01 and 99.

\pandocbounded{\includegraphics[keepaspectratio]{https://raw.githubusercontent.com/JantN1966/MasterManual/main/Images/Observ02.jpg}}\\

\emph{Example}. A covariate table file for an independent variable with
values in the data between 86 and 115. The order of the Legendre
polynomial is 2. The table was created with the line !CVRMAKE LEG
!CVRNUM 2 !CVRMIN 86 !CVRMAX 115 in the CVRTABLE section of the
instruction file.

\subparagraph{4.2.2.1. Syntax using an existing covariate table for the
MiX99 solver}\label{Obse10}

\begin{quote}
DATAFILE \textless filename\textgreater{}\\
\ldots{}\\
\textless field k\textgreater{} I !CVRIND\\
\ldots{}\\
CVRTABLE \textless filename\textgreater{}\\
MODEL\\
\textless trait\textgreater{} \textasciitilde{} \textless fixed
effects\textgreater{} \textless Class1\textgreater*CVR(n1) !RANDOM
\textless Class2\textgreater*CVR(n2) G(Animal*CVR(n3))\\
\ldots{}
\end{quote}

Sections:

\textbf{CVRTABLE}/ The CVRTABLE section contains the details of the
existing or new covariate table.

Qualifiers:

\textbf{!CVRIND} / The field marked with !CVRIND is the independent
variable used in polynomial regression. Any level of the field specified
with !CVRIND must exist in the covariate table file. The field must not
contain a missing value indicator for a valid trait observation. The
qualifier !CVRIND must be used when the section CVRTABLE is specified.
The field must be integer. The qualifier !CVRIND, specified in DATAFILE
section, should not be confused with !CVRindex that is used with hpblup
solver and specified in the CVRTABLE section.

\textbf{CVR(\ldots)} / The CVR function is used in the MODEL section and
is a shorthand for all polynomial terms to be fitted and may be used in
the same way as any individual random regression term. The alternative
way to specify polynomial random regression is to use the individual
columns of the covariate table file. The names of the columns are cvr00,
cvr01, cvr02, \ldots, cvrnn.

\subparagraph{4.2.2.2. Syntax using an existing covariate table for the
hpblup solver}\label{Obse11}

\begin{quote}
\ldots{}\\
CVRTABLE !nCVRTABLES 2\\
TABLE01 \textless filename\textgreater{} !CVRNUM
\textless n\textsuperscript{th} order\textgreater{} !CVRMIN
\textless minimum value\textgreater{} !CVRMAX \textless maximum
value\textgreater{} !CVRSingleCov !CVRIndex \textless index field
name\textgreater{}\\
TABLE04 \textless filename\textgreater{} !CVRNUM
\textless n\textsuperscript{th} order\textgreater{} !CVRMIN
\textless minimum value\textgreater{} !CVRMAX \textless maximum
value\textgreater{} !CVRSingleCov !CVRIndex \textless index field
name\textgreater{}\\
\strut \\
MODEL\\
\textless trait\textgreater{} \textasciitilde{} \textless fixed
effects\textgreater{} \textless Class1\textgreater*TABLE01 !RANDOM
\textless Class2\textgreater*TABLE04 G(Animal*TABLE04)\\
\ldots{}
\end{quote}

Additional qualifiers:

\textbf{!nCVRTABLES}\\
This qualifier specifies the number of covariate tables included in this
section

\textbf{!CVRIndex}\\
This qualifier specifies the field name of the index in the DATAFILE.
Please note that this option is different from !CVRIND, which is used
with the MiX99 solver and specified in the DATAFILE section.

\textbf{!CVRSingleCov}\\
This qualifier is used to create a separate file for each covariate in
table specified. Each covariate in the table is then be fitted as a
separate effect for the hpblup solver.

\textbf{TABLE\emph{tt} in the MODEL section}\\
A covariate table file specified in the CVRTABLE section can be fitted
in the model by fitting its label. It may be used in the same way as any
individual random regression term. The names of its columns in variance
covariance matrix files are cvr\emph{tt}\_00 to cvr\emph{tt}\_\emph{nn},
where \emph{tt} is the number in the label of the covariate table and
\emph{nn} the order of the polynomial specified for the covariate table
\emph{tt}.

\subparagraph{4.2.2.3. Syntax using a newly created covariate table for
the MiX99 solver}\label{Obse12}

\begin{quote}
DATAFILE \textless filename\textgreater{}\\
\ldots{}\\
\textless field k\textgreater{} I !CVRIND\\
\ldots{}\\
CVRTABLE !CVRMAKE LEG !CVRNUM \textless n\textsuperscript{th}
order\textgreater{} !CVRMIN \textless minimum value\textgreater{}
!CVRMAX \textless maximum value\textgreater{}\\
MODEL \textless trait\textgreater{} \textasciitilde{} \textless fixed
effects\textgreater{} \textless Class1\textgreater*CVR(n1) !RANDOM
\textless Class2\textgreater*CVR(n2) G(Animal*CVR(n3))\\
\ldots{}
\end{quote}

Additional qualifiers:

\textbf{!CVRMAKE}\\
If !CVRMAKE is specified, MiXBLUP generates a covariate table file using
the settings specified with the !CVRNUM, !CVRMIN and !CVRMAX qualifiers.
Currently, only a covariate table containing Legendre polynomials can be
created, by specifying LEG as the argument of !CVRMAKE. The name of the
new covariate table file is `cvrtable.txt'.

\textbf{!CVRNUM}\\
The qualifier !CVRNUM must be specified and is used to specify the order
of the polynomial in the covariate table. The expected number of columns
to read is the order + 2, one for the level of the independent variable
and one for the order being 0. It is up to the user to make sure that
the order specified in the MODEL section is equal to or lower than the
order specified with !CVRNUM.

\textbf{!CVRMIN and !CVRMAX}\\

The qualifiers !CVRMIN and !CVRMAX can be used to specify the lowest and
highest value of the independent variable that were used to estimate the
genetic parameters. Legendre polynomials are dependent on the lowest and
highest value of the independent variable and so are the genetic
parameters of Legendre polynomials. If !CVRMIN or !CVRMAX is
nevertheless omitted, the lowest or highest value of the independent
variable in the data is used, instead.

\subparagraph{4.2.2.4. Syntax using newly created covariate tables for
the hpblup solver}\label{Obse13}

\begin{quote}
\ldots{}\\
CVRTABLE !nCVRTables \textless value\textgreater{}\\
TABLE01 !CVRMAKE LEG !CVRSingleCov !CVRNUM
\textless n\textsuperscript{th} order\textgreater{} !CVRMIN
\textless minimum value\textgreater{} !CVRMAX \textless maximum
value\textgreater{} !CVRIndex \textless field name of the
index\textgreater{}\\
TABLE02 !CVRMAKE LEG !CVRNUM \textless n\textsuperscript{th}
order\textgreater{} !CVRMIN \textless minimum value\textgreater{}
!CVRMAX \textless maximum value\textgreater{} !CVRIndex \textless field
name of the index\textgreater{}\\
\ldots{}\\
MODEL \textless trait\textgreater{} \textasciitilde{} \textless fixed
effects\textgreater{} TABLE01 !RANDOM
\textless Class\textgreater*TABLE02 G(TABLE02*animal)
\end{quote}

Additional qualifiers:

\textbf{!CVRMAKE}\\

If !CVRMAKE is specified, MiXBLUP generates a covariate table file using
the settings specified with the !CVRNUM, !CVRMIN and !CVRMAX qualifiers.
Currently, only a covariate table containing Legendre polynomials can be
created, by specifying LEG as the argument of !CVRMAKE. The name of the
new covariate table file is `hpTable\emph{tt}.txt', for example
hpTable01.txt. If !CVRSingleCov is specified, a separate file is created
for each covariate. In that case, the names of the new covariate table
files are `hpTable\emph{tt}\_\emph{nn}.txt', for example
hpTable01\_00.txt, where \emph{tt} is the number in the label of the
covariate table and \emph{nn} the number of the covariate of the
polynomial specified for the covariate table \emph{tt}, ranging from 0
to the order specified.

\paragraph{4.2.3. Associated output files}\label{Obse15}

{\def\LTcaptype{none} % do not increment counter
\begin{longtable}[]{@{}ll@{}}
\toprule\noalign{}
Output file & Description \\
\midrule\noalign{}
\endhead
\bottomrule\noalign{}
\endlastfoot
cvrtable.txt & covariate table, if created by MiXBLUP \\
\end{longtable}
}

\subsubsection{4.3. General covariate files}\label{Obse16}

\paragraph{4.3.1. General}\label{Obse17}

Some covariates are individual-specific: they never change for an
individual, but vary across individuals. They are more associated with
the individual than with its data records. Examples are breed
composition, genetic groups, heterosis and recombination. Such
covariates can be stored in a covariate file, in which all individuals
in the analysis have a record. MiXBLUP converts the covariate file with
all individuals to a data covariate file that exactly matches the data
file, including repeated records.

\paragraph{4.3.2. Input file}\label{Obse18}

General covariate files contain at least the ID of the animal and any
number of covariates, but all records should have the same number of
covariates. General covariate files must be provided in space-separated
format. Covariates are read as real numbers, regardless of whether a
decimal point is present in the corresponding field. General covariate
files contain at least all individuals with a phenotype for any of the
traits in the statistical model. Individuals without any phenotypes will
be ignored, except in the case of genetic group covariates.

\pandocbounded{\includegraphics[keepaspectratio]{https://raw.githubusercontent.com/JantN1966/MasterManual/main/Images/Observ03.jpg}}\\

\emph{Example}. Covariate file with breed fractions in a mixed breed
population

\paragraph{4.3.3. Syntax}\label{Obse19}

\subparagraph{4.3.3.1. Syntax of a general covariate file and associated
variance-covariance file}\label{Obse20}

\begin{quote}
REGFILE\\
\textless field animal\textgreater{} \textless field type I or
A\textgreater{}\\
REG01 \textless file name REG01\textgreater{} !REGTYPE F/R/H {[}!IDCOL
1{]} {[}!STARTCOV 2{]} {[}!LASTCOV 7{]}\\
REG02 \textless file name REG02\textgreater{} !REGTYPE F/R/H {[}!IDCOL
1{]} {[}!STARTCOV 2{]} {[}!LASTCOV 7{]}\\
\ldots{}\\
REG99 \textless file name REG99\textgreater{} !REGTYPE F/R/H {[}!IDCOL
1{]} {[}!STARTCOV 2{]} {[}!LASTCOV 7{]}\\
\strut \\
REGPARFILE\\
REG01 \textless file name REG01\textgreater{}\\
REG02 \textless file name REG02\textgreater{}\\
\ldots{}\\
REG99 \textless file name REG99\textgreater{}
\end{quote}

Sections:

\textbf{REGFILE}\\
The REGFILE section specifies the name of one or more general covariate
files and its attributes, such as column numbers and whether one
variance for all covariates is used or an individual variance for each
covariate.

\textbf{REGPARFILE}\\
The REGPARFILE section is used to specify a file with components of
variance and covariance among traits associated with general covariates.
A general covariate file labelled in REGFILE needs a corresponding entry
in a REGPARFILE section if the regression type is R for random or H for
heterogeneous variances. There are no file-independent qualifiers.

Qualifiers:

The file-dependent qualifiers of REGFILE can be specified for each
covariate file. These qualifiers are explained below.

\textbf{!REGTYPE}\\
The file-specification line must contain the !REGTYPE qualifier. It
specifies how the covariates in the file are fitted in the model. If `f'
is specified, the covariates in the file are fitted as a fixed
regression. Covariates fitted as a fixed effect do not have a variance
associated with it, so it is not necessary to specify a parameter file
in the REGPARFILE section. If it is present, it is ignored. If `r' is
specified, the covariates in the file are fitted as a random regression
with a single variance for all covariates in the file. The variance is
specified in the corresponding parameter file in the REGPARFILE section.
If `h' is specified, the covariates in the file are fitted as a random
regression, each with their own variance. The covariate-specific
variances are specified in the corresponding parameter file in the
REGPARFILE section.

\textbf{!IDCOL}\\
The !IDCOL qualifier is optional and specifies which field in the
covariate file contains the ID of the individual. If it is omitted, it
is assumed that the ID is in the first field of the record (so the
default is !IDCOL 1).

\textbf{!STARTCOV}\\
The !STARTCOV qualifier is optional and specifies which field contains
the first covariate. If it is omitted, it is assumed that the covariates
start in the second field of the record (so !STARTCOV 2).

\textbf{!LASTCOV}\\
The !LASTCOV qualifier is optional and specifies which field contains
the last covariate of the file to include in the model. If it is
omitted, it is assumed that all fields after the first covariate contain
covariates to include in the model.

\subparagraph{4.3.3.2. Syntax of fitting a general covariate file in the
model for the MiX99 solver}\label{Obse23}

\begin{quote}
MODEL trait \textasciitilde{} fixed !RANDOM REG(1,2..5)
\end{quote}

Qualifiers:

\textbf{REG(\ldots)}\\
The REG function is used in the MODEL section and can be used to specify
which general covariate files should be fitted in the model of a trait.
If a covariate file is specified, then all specified covariates in the
file will be fitted simultaneously. The numbers in the REG(\ldots)
function link to the number in the label of the general covariate file
in the REGFILE section (and the REGPARFILE section). The numbers may be
specified individually as (1, 2, 3, 4) or as a range, indicated by two
subsequent full stops, for example (1..4), or a combination of both. If
a covariate file is fitted for any trait through REG(\ldots), the
covariates will be fitted for all traits, even the ones for which
REG(\ldots) is not specified.

\subparagraph{4.3.3.3. Syntax of fitting a general covariate file in the
model for the hpblup solver}\label{Obse25}

\begin{quote}
MODEL trait \textasciitilde{} \textless fixed\textgreater{} !RANDOM
hpREG(1,\textless field index\textgreater)
\end{quote}

Qualifiers:

\textbf{hpREG(\textless number in label of covariate file\textgreater,
\textless field index\textgreater)}\\
The hpReg function is used to fit a general covariate file in the model
of a trait, for which it is specified. For a random effect, REGTYPE
needs to be set to R or H and hpREG needs to be specified after the
!Random qualifier. For a fixed effect, REGTYPE needs to be set to F and
hpREG needs to be specified before the !Random qualifier.

\paragraph{4.3.4. Associated output files}\label{Obse27}

{\def\LTcaptype{none} % do not increment counter
\begin{longtable}[]{@{}
  >{\raggedright\arraybackslash}p{(\linewidth - 2\tabcolsep) * \real{0.5000}}
  >{\raggedright\arraybackslash}p{(\linewidth - 2\tabcolsep) * \real{0.5000}}@{}}
\toprule\noalign{}
\begin{minipage}[b]{\linewidth}\raggedright
Output file
\end{minipage} & \begin{minipage}[b]{\linewidth}\raggedright
Description
\end{minipage} \\
\midrule\noalign{}
\endhead
\bottomrule\noalign{}
\endlastfoot
RegCov\%\%.txt & temporary file; data covariate file \\
RegCov\%\%NoDat.txt & temporary file; covariates of individuals without
any phenotypes \\
Solreg\_mat.txt & solutions of all covariates in any general or SNP
covariate file \\
\end{longtable}
}

\subsubsection{4.4. Random effects with correlated level
effects}\label{Obse28}

\paragraph{4.4.1. General}\label{Obse29}

For non-genetic random effects, it is often assumed that level effects
are uncorrelated. In practice, this may not be a valid assumption, for
example for subsequent year-seasons within a herd. For these cases, the
user may provide a correlation matrix to model that some level effects
are more similar than others.

\paragraph{4.4.2. Input file}\label{Obse30}

The inverse of the correlation matrix has to be provided as a sparse
matrix in I-J-Value format. The file contains a line for each non-zero
element in the matrix. The line contains effect label of row, effect
label of column, non-zero element. It has to be provided in
upper-triangular format, so I is equal to or lower than J in the
I-J-Value format. The user has to verify that the inverse correlation
matrix is positive definite and not close to singularity.

\subsubsection{4.5. Syntax}\label{Obse31}

\begin{quote}
CORRFILE\\
RCE01 \textless file name RCE01\textgreater{}\\
RCE05 \textless file name RCE05\textgreater{}\\
\strut \\
MODEL\\
\textless trait\textgreater{} \textasciitilde{}
\textless fixed\textgreater{} !RANDOM RCE(\textless random effect name
for RCE01\textgreater,1) RCE(\textless random effect name for
RCE05\textgreater,5) G(\ldots)
\end{quote}

Sections:

\textbf{CORRFILE}\\
The CORRFILE section specifies the name of one or more inverse
correlation matrix files for non-genetic random correlated effects. The
CORRFILE section does not have qualifiers for non-genetic random
effects. For use of CORRFILE for specifying additional genetic
relationship matrices.

\paragraph{4.5.1. Associated output files}\label{Obse32}

Output files are the same as for non-genetic random uncorrelated
effects.

\hyperref[Tabl01]{Back to Table of Contents}

\newpage

\subsection{5. Genetic similarity among individuals}\label{Gene01}

\textbf{Two individuals that have an ancestor in common are more similar
than two unrelated individuals. This genetic similarity can be specified
in various ways. This chapter describes the recommended methods in
MiXBLUP to specify genetic similarity.}

Chapter \hyperref[Gene02]{5.1.} describes the format of pedigree
information that is used to build A\textsuperscript{-1}. If only a
pedigree is available, MiXBLUP will calculate the expected genetic
relationships between individuals as they appear in the inverse pedigree
relationship matrix (A\textsuperscript{-1}), without the need to specify
this matrix explicitly (chapter \hyperref[Gene26]{5.3.}).\\
Chapter \hyperref[Gene21]{5.2.} describes the recommended format of
genomic data. If all or part of the individuals were genotyped for many
genetic markers, such as SNPs, MiXBLUP can be used to estimate true
genetic similarity from genomic data (chapter \hyperref[Gene29]{5.4.}).
One method is to calculate the estimated true genetic relationships in a
genomic relationship matrix. This inverse genomic relationship matrix
can be used on its own if no pedigree information is available (chapter
\hyperref[Gene32]{5.4.1.2.}). It can also be combined with pedigree
information to analyse genotyped and non-genotyped individuals
simultaneously (chapter \hyperref[Gene37]{5.4.2.2.}). Pedigree
information can also be used if all individuals are genotyped.\\
An equivalent method to use estimated true genetic relationships
implicitly, without the need to construct and invert a genomic
relationship matrix, is random regression of all SNPs simultaneously on
the data (chapter \hyperref[Gene33]{5.4.1.3.} and chapter
\hyperref[Gene34]{5.4.1.4.}) and \hyperref[Gene47]{5.4.2.4.}).\\
Genetic similarity in case of multiple breeds and crosses can be
addressed with breed-specific allele frequencies, breed-specific genetic
groups or a fixed effect of breed composition in the model (chapter
\hyperref[Gene58]{5.6.})\\
It can be necessary to provide an existing inverse relationship matrix
if, for example, the Henderson rules to calculate the inverse pedigree
relationship matrix directly do not apply. MiXBLUP will use this matrix
to model genetic similarity between individuals (chapter
\hyperref[Gene54]{5.5.}). The G(\ldots) function and the SNP(\ldots) or
hpSNP(\ldots) functions in the MODEL section are used to link genetic
similarity to data records (Chapter \hyperref[Stat01]{7.}).

\subsubsection{5.1. Preparing pedigree data}\label{Gene02}

\paragraph{5.1.1. General}\label{Gene03}

Expected genetic similarity between individuals can be based on observed
pedigree relationships. Any individual occurring in the data file,
regardless of whether it occurs with a record or as a maternal, paternal
or group mate genetic effect (in case of a social interaction model),
must be present in the pedigree file. Any individual that does not
appear in the data file, but exists as an ancestor in the pedigree file,
must also have its own record in the pedigree file.

\paragraph{5.1.2. Recommended file formats}\label{Gene04}

\subparagraph{5.1.2.1. Pedigree file}\label{Gene05}

The pedigree file consists of the individual identification code (ID)
and the IDs of its sire and dam in the first three columns. The columns
must be separated by at least one space. The IDs in the pedigree file
must be of same type as the IDs in the data file (either numeric or
text). The pedigree file may contain other information in any number of
additional columns, as long as the number of columns is the same for all
records.\\
Calculating reliabilities requires a block variable to be present in the
pedigree file (see Chapter \hyperref[Reli01]{9.}). In that case the
pedigree file, as well as the data file, will be sorted on the block
variable. If a block group variable is added to the pedigree, it must be
marked with the qualifier !BLOCK. It does not have to be in the fourth
column, as in older versions of MiXBLUP. The pedigree file does not need
to be sorted. MiXBLUP takes care of any required sorting.

\pandocbounded{\includegraphics[keepaspectratio]{https://raw.githubusercontent.com/JantN1966/MasterManual/main/Images/GenSim01.jpg}}\\

\emph{Example}. Pedigree file with a single code for unknown parents/

\subparagraph{5.1.2.2. Pedigree inbreeding coefficient
file}\label{Gene06}

A file with previously calculated pedigree inbreeding coefficients can
be any free-format text file with any number of columns, as long as it
contains the ID of each individual in the analysis and its inbreeding
coefficient. This may be the pedigree file with an additional column of
inbreeding coefficients.

\pandocbounded{\includegraphics[keepaspectratio]{https://raw.githubusercontent.com/JantN1966/MasterManual/main/Images/GenSim02.jpg}}\\

\emph{Example}. File with inbreeding coefficients

\paragraph{5.1.3. Pedigree inbreeding coefficients}\label{Gene07}

\subparagraph{5.1.3.1. General}\label{Gene08}

Inbreeding coefficients are often ignored in breeding value estimation
using pedigree relationships only. The internally calculated numerator
relationship matrix (A-1) is by default set up without taking into
account inbreeding. Inbreeding can be included by providing the kernel
with a file with the inbreeding coefficient of each individual in the
pedigree file. This file may be provided as an existing input file or
calculated within MiXBLUP as a preparation step. Note that inbreeding
coefficients do not affect the reliability calculation and will be
ignored.

\subparagraph{5.1.3.2. Syntax of calculating inbreeding
coefficients}\label{Gene09}

\begin{quote}
PEDFILE \textless pedigree file\textgreater{} {[}!CALCINBR
\textless method\textgreater{]}\\
\textless field animal\textgreater{} \textless field
type\textgreater{}\\
\textless field sire\textgreater{} \textless field type\textgreater{}\\
\textless field dam\textgreater{} \textless field type\textgreater{}
\end{quote}

Qualifier:

\textbf{!CALCINBR \textless method\textgreater{}}\\
The qualifier CALCINBR is optional and is used to indicate that
inbreeding coefficients should be calculated and included in the
calculation of the inverse pedigree relationship matrix (A-1). If
!CALCINBR has been specified, the section INBRFILE is ignored.\\
If neither !CALCINBR, nor INBRFILE is specified for a genetic evaluation
for which only pedigree information is available, then inbreeding
coefficients are not included in the inverse pedigree relationship
matrix. For genomic evaluations, however, the default setting is
different for the two solvers. For the MiX99 solver, the default is that
inbreeding coefficients are not taken into account if it is not
specified to include them. For the hpblup solver, the default is to
always calculate inbreeding coefficients if pedigree information is
available for the evaluation.\\
There are two methods available to calculate inbreeding coefficients.
The default method is published by Sargolzaei et al.~(2005) and can be
specified as !CalcInbr or !CalcInbr S{[}argolzaei{]}. The alternative
method is published by Meuwissen and Luo (1992) and can be specified as
!CalcInbr M{[}euwissen{]}. Which algorithm is fastest, depends on the
structure of the pedigree.

\subparagraph{5.1.3.3. Syntax of using file with inbreeding
coefficients}\label{Gene10}

\begin{quote}
INBRFILE \textless inbreeding coefficient file \textgreater{} {[}!IDCOL
\textless field number\textgreater{]} {[}!INBRCOL \textless field
number\textgreater{]}
\end{quote}

Qualifier:

\textbf{!IDCOL }\\
The optional qualifier !IDCOL can be used to specify the field number in
the inbreeding coefficient file that contains the animal ID. The default
field number is 1.

\textbf{!INBRCOL }\\
The optional qualifier !INBRCOL can be used to specify the field number
in the inbreeding coefficient file that contains the inbreeding
coefficient. The default field number is 4.

\paragraph{5.1.4. Unknown parents are from a single large base
population}\label{Gene12}

It is inevitable that for at least some individuals in the pedigree, the
parents are unknown. If it is reasonable to assume that these unknown
parents come from the same large population, they should be coded with a
zero (0).

\paragraph{5.1.5. Unknown parents are from multiple large base
populations}\label{Gene13}

For pedigrees with unknown parents from various known origins, or many
individuals without known parents across generations, it may be
desirable to specify that some individuals with unknown parents are more
similar than average. For example, in case of genetic selection, two
individuals born in the same year are more similar than two individuals
born in different years. In case of a large difference in selection
differential between males and females, it may be useful to distinguish
males and females born in the same year. In case of mixed-breed or
mixed-line evaluations, it may be useful to group individuals by breed,
line or type of cross. This can be done by assigning individuals with
one or two unknown parents to an appropriate genetic (or phantom parent)
group.\\
Genetic groups can be included in the analysis in two ways: (1) Westell
grouping and (2) genetic group covariates. Westell grouping augments the
pedigree relationship matrix with the number of genetic groups. For
genetic group covariates, a covariate matrix Q is set up that contains
the proportion of each genetic group for each animal. Genetic group
covariates can be provided as an existing covariate file or calculated
as part of the evaluation. If genetic group covariates are provided as
an existing covariate file, then the pedigree does not need to contain
genetic groups, as well. In that case, genetic groups will be replaced
with unknown parents. For both Westell grouping and genetic group
covariates, the genetic solutions include the genetic group effect.\\
In the pedigree file, the genetic group of the individual is entered on
the position of the unknown parent. Genetic groups must be coded as
negative integers, but do not have to be sequentially numbered.\\
Genetic groups can be modelled either as fixed, pseudo-random (Westell
grouping) or random effects. For Westell grouping, the specified value
will be added to the diagonal elements of the genetic group effects in
the inverse coefficient matrix. If a value of zero is added, genetic
group effects are modelled as fixed effects. For values larger than
zero, genetic groups are modelled as pseudo-random effects. The larger
the value, the more estimates are regressed towards the mean. For
genetic group covariates, a variance component can be specified for each
genetic group covariate separately or one for all genetic group
covariates. It is also possible to fit the covariates as fixed effects.

\pandocbounded{\includegraphics[keepaspectratio]{https://raw.githubusercontent.com/JantN1966/MasterManual/main/Images/GenSim03.jpg}}\\

\emph{Example}. Pedigree file with genetic groups for unknown parents

\subparagraph{5.1.5.1. Syntax of multiple large base populations using
Westell grouping}\label{Gene14}

\begin{quote}
PEDFILE \textless pedigree file\textgreater{} {[}!GROUPS
\textless value\textgreater{]}\\
\textless field animal\textgreater{} \textless field
type\textgreater{}\\
\textless field sire\textgreater{} \textless field type\textgreater{}\\
\textless field dam\textgreater{} \textless field type\textgreater{}
\end{quote}

Qualifier:

\textbf{!Groups \textless value\textgreater{}} The qualifier GROUPS
means that genetic groups are included in the pedigree. Genetic groups
need to be coded with negative integer values. With
\textless value\textgreater, it is possible to specify whether these
Genetic group effects should be modelled as fixed (value = 0.0) or as
random (value \textgreater{} 0.0). In practice, !GROUPS does not need to
be set at a much higher value than about 3.

\subparagraph{5.1.5.2. Associated output files for Westell
grouping}\label{Gene15}

{\def\LTcaptype{none} % do not increment counter
\begin{longtable}[]{@{}
  >{\raggedright\arraybackslash}p{(\linewidth - 2\tabcolsep) * \real{0.5000}}
  >{\raggedright\arraybackslash}p{(\linewidth - 2\tabcolsep) * \real{0.5000}}@{}}
\toprule\noalign{}
\begin{minipage}[b]{\linewidth}\raggedright
Output file
\end{minipage} & \begin{minipage}[b]{\linewidth}\raggedright
Description
\end{minipage} \\
\midrule\noalign{}
\endhead
\bottomrule\noalign{}
\endlastfoot
Solani.txt & Solutions of the direct genetic effect including the
genetic group effects \\
& when the field type of the ID is integer \\
Solani.out & Solutions of the direct genetic effect including the
genetic group effects \\
& when the field type of the ID is alphanumerical \\
Relani.txt & Approximate reliabilities when the field type of the ID is
integer \\
Relani.out & Approximate reliabilities when the field type of the ID is
alphanumerical \\
\end{longtable}
}

\subparagraph{5.1.5.3. Syntax of multiple large base populations using
precalculated genetic group covariates}\label{Gene16}

\begin{quote}
PEDFILE \textless pedigree file\textgreater{}\\
\textless field animal\textgreater{} \textless field
type\textgreater{}\\
\textless field sire\textgreater{} \textless field type\textgreater{}\\
\textless field dam\textgreater{} \textless field type\textgreater{}\\
REGFILE\\
\textless field animal\textgreater{} \textless field type I or
A\textgreater{}\\
REG01 \textless covariate file name REG01\textgreater{} !GGCOV !REGTYPE
F/R/H\\
{[}REGPARFILE{]}\\
{[}REG01 \textless parameter file name REG01\textgreater{]}\\
MODEL\\
\textless trait\textgreater{} \textasciitilde{} \textless fixed
effects\textgreater{} !RANDOM REG(1) \textless other random
effects\textgreater{}
\end{quote}

Qualifier:

\textbf{!GGcov}\\
The qualifier !GGcov specifies which external covariate file contains
genetic group covariates. If !MakeGGcov is specified, there is no need
to specify a file name for the covariate file with !GGcov

\textbf{REG(\ldots) or hpReg(\ldots)}\\
When using the MiX99 solver, the REG function can be used to fit a
genetic group covariate file in the model of a trait. If the genetic
group covariate file is fitted for any trait through REG(\ldots), the
covariates will be fitted for all traits, even the ones for which
REG(\ldots) is not specified.\\
The numbers in the REG(\ldots) function link to the number in the label
of the general covariate file in the REGFILE section (and the REGPARFILE
section). The numbers may be specified individually as (1, 2, 3, 4) or
as a range, indicated by two subsequent full stops, for example (1..4),
or a combination of both. The index is the individual's ID in the data
file.\\
When using the hpblup solver, the hpReg function can be used to fit a
genetic group covariate file in the model of a trait. Note that a
genetic group covariate file fitted through hpReg(\ldots) is only fitted
for the traits for which it is in the model.\\
The hpReg function has two parameters. The first one is the label number
of the covariate file in the REGFILE section. The second parameter is
the field name in the data file of the index of the covariate file.

\subparagraph{5.1.5.4. Syntax of multiple large base populations using
calculated genetic group covariates}\label{Gene99}

\begin{quote}
PEDFILE \textless pedigree file\textgreater{} !MAKEGGCOV\\
\textless field animal\textgreater{} \textless field
type\textgreater{}\\
\textless field sire\textgreater{} \textless field type\textgreater{}\\
\textless field dam\textgreater{} \textless field type\textgreater{}\\
REGFILE\\
\textless field animal\textgreater{} \textless field type I or
A\textgreater{}\\
REG01 !GGCOV !REGTYPE F/R/H\\
{[}REGPARFILE{]}\\
{[}REG01 \textless parameter file name REG01\textgreater{]}\\
MODEL\\
\textless trait\textgreater{} \textasciitilde{} \textless fixed
effects\textgreater{} !RANDOM REG(1) \textless other random
effects\textgreater{}
\end{quote}

Qualifier: \textbf{!MakeGGcov}\\
The qualifier !MakeGGcov is optional and triggers MiXBLUP to set up a
covariate matrix Q of the number of genetic groups by the number of
individuals in the analysis. The covariates are stored in a standard
covariate file.

\subparagraph{5.1.5.5. Associated output files for genetic group
covariates}\label{Gene17}

{\def\LTcaptype{none} % do not increment counter
\begin{longtable}[]{@{}
  >{\raggedright\arraybackslash}p{(\linewidth - 2\tabcolsep) * \real{0.5000}}
  >{\raggedright\arraybackslash}p{(\linewidth - 2\tabcolsep) * \real{0.5000}}@{}}
\toprule\noalign{}
\begin{minipage}[b]{\linewidth}\raggedright
Output file
\end{minipage} & \begin{minipage}[b]{\linewidth}\raggedright
Description
\end{minipage} \\
\midrule\noalign{}
\endhead
\bottomrule\noalign{}
\endlastfoot
Solani.txt & Solutions of the direct genetic effect when the field type
of the ID is integer \\
SolaniGG.txt & Solutions of the accumulated genetic group effects for
each individual \\
Solanitot.txt & Solutions of the direct genetic effect including the
genetic group effects \\
& when the field type of the ID is integer \\
GeneticGroupsInQ.txt & Original genetic group label by column in the
covariate file \\
Solani.out & Solutions of the direct genetic effect when the field type
of the ID is alphanumerical \\
Solreg\_mat.txt & Solutions of genetic group covariates, along with
solutions \\
& of any other external covariates \\
SolaniGG.out & Solutions of the accumulated genetic group effects for
each individual \\
Solanitot.out & Solutions of the direct genetic effect including the
genetic group effects \\
& when the field type of the ID is alphanumerical \\
Relani.txt & Approximate reliabilities when the field type of the ID is
integer \\
Relani.out & Approximate reliabilities when the field type of the ID is
alphanumerical \\
\end{longtable}
}

\paragraph{5.1.6. Unknown parents are from multiple related base
populations (metafounders)}\label{Gene18}

If individuals without known parents originate from multiple base
populations, it may be reasonable to assume that these base populations
are not unrelated. Information on average relationships within and
between base populations may come from genotyped descendants of the
unknown parents or from discarded pedigree information. Related base
populations are generally referred to as metafounders. Average genetic
relationships within and between metafounders are presented to the
solver in an inverse gamma matrix. In case of relationships because of
discarded pedigree and/or genotype information, the user needs to
calculate the gamma matrix prior to the evaluation and specify the name
of the matrix. If the relationships are entirely based on genotypes
included in the evaluation, then MiXBLUP can also calculate the gamma
matrix. Metafounders are presented in the pedigree file in the same way
as genetic groups, i.e.~as negative integers. The file with the gamma
matrix should be a text file in I-J-Value format, i.e.~metafounder ID of
row, metafounder ID of column, average genetic relationship. Only
non-zero genetic relationships of the lower triangular part of the
matrix need to be specified.

\subparagraph{5.1.6.1. Syntax of multiple related base populations using
metafounders}\label{Gene19}

\begin{quote}
PEDFILE \textless pedigree file\textgreater{} !Metafounders\\
\textless field animal\textgreater{} \textless field
type\textgreater{}\\
\textless field sire\textgreater{} \textless field type\textgreater{}\\
\textless field dam\textgreater{} \textless field type\textgreater{}
\end{quote}

Qualifier:

\textbf{!Metafounders}\\
or\\
\textbf{!Metafounders \textless file with gamma matrix\textgreater{}}\\
This qualifier indicates that base populations in the pedigree are
related. If a file with the gamma matrix is specified, the gamma matrix
is coded for either the default or the hpblup solver. If a file with the
gamma matrix is not specified, the gamma matrix with genomic
relationships within and between metafounders is estimated from
available genomic information. Metafounders are fitted using this gamma
matrix with coded IDs, and QP transformation.

\subparagraph{5.1.6.2. Associated output files for
metafounders}\label{Gene20}

{\def\LTcaptype{none} % do not increment counter
\begin{longtable}[]{@{}
  >{\raggedright\arraybackslash}p{(\linewidth - 2\tabcolsep) * \real{0.5000}}
  >{\raggedright\arraybackslash}p{(\linewidth - 2\tabcolsep) * \real{0.5000}}@{}}
\toprule\noalign{}
\begin{minipage}[b]{\linewidth}\raggedright
Output file
\end{minipage} & \begin{minipage}[b]{\linewidth}\raggedright
Description
\end{minipage} \\
\midrule\noalign{}
\endhead
\bottomrule\noalign{}
\endlastfoot
Solani.txt & Solutions of the direct genetic effect when the field type
of the ID is integer \\
Solani.out & Solutions of the direct genetic effect when the field type
of the ID is alphanumerical \\
Relani.txt & (Weighted) inverse genomic relationship matrix \\
Relani.out & Approximate reliabilities when the field type of the ID is
integer \\
ExtRelMat\_tri.txt & Weighted inverse genomic relationship matrix \\
gamma.dat & gamma matrix with relationships within and across
metafounders \\
\end{longtable}
}

\subsubsection{5.2. Preparing genomic data}\label{Gene21}

\paragraph{5.2.1. General}\label{Gene22}

Genomic data can be used to estimate true genetic similarity between
individuals. Genomic data consists of a large number of bi-allelic SNP
markers.

\paragraph{5.2.2. Recommended file formats}\label{Gene23}

\subparagraph{5.2.2.1. Genomic data}\label{Gene24}

It is recommended to provide genomic data as genotypes, which is the
count of one of the two alleles. This may be provided as a
space-separated or dense text file or in binary format.

\pandocbounded{\includegraphics[keepaspectratio]{https://raw.githubusercontent.com/JantN1966/MasterManual/main/Images/GenSim04.jpg}}\\

\emph{Example}. Genotype file with marker genotype data per animal in
dense format. It contains the number of copies per locus of the allele
with the highest number (11=0, 12=1 and 22=2).

The recommended way to provide genomic data is binary plink format. It
consists of three files. The \textbf{.bim} file contains the SNP marker
details in text format, the \textbf{.fam} file contains details of
genotyped individuals in text format and the \textbf{.bed} file contains
the genotype of each individual for each SNP marker in compressed binary
format. For GBLUP and single-step GBLUP, genomic data may be also
provided as alleles, either using pairs of alleles on a single record
per individual or splitting pairs of alleles onto two records per
individual. This may be useful for using multi-allele loci. See appendix
\textless Genomic data files\textgreater{} for details.

\subparagraph{5.2.2.2. Allele frequencies}\label{Gene25}

If the user does not want to use allele frequencies calculated from the
data, then pre-calculated allele frequencies can be supplied as an
additional input file, The file specified should contain for each locus
the allele frequency of the allele with the highest integer code, if the
genetic marker file contains alleles. The file specified should contain
for each locus the frequency of the allele of which the homozygote
genotypes are coded as 2. The structure of the file is \textless locus
number in order of the genetic marker file\textgreater{}
\textless allele frequency\textgreater. Pre-calculated allele
frequencies are supported for GBLUP, single-step GBLUP, SNPBLUP and
single-step SNPBLUP.

\pandocbounded{\includegraphics[keepaspectratio]{https://raw.githubusercontent.com/JantN1966/MasterManual/main/Images/GenSim08.jpg}}\\

\emph{Example}. Pre-calculated allele frequency per locus of allele
coded with the highest integer code for 6 loci

\subsubsection{5.3. Genetic similarity from pedigree only}\label{Gene26}

\paragraph{5.3.1. General}\label{Gene27}

MiXBLUP supports analyses using a pedigree that consists of individuals
and their parents (animal model). A sire model with sires and maternal
grandsires in the pedigree file is currently not supported in MiXBLUP.

\paragraph{5.3.2. Syntax of using pedigree BLUP}\label{Gene28}

\begin{quote}
PEDFILE \textless pedigree file\textgreater{} {[}!SKIP \textless n
lines\textgreater{]}\\
\textless field animal\textgreater{} \textless field
type\textgreater{}\\
\textless field sire\textgreater{} \textless field type\textgreater{}\\
\textless field dam\textgreater{} \textless field type\textgreater{}\\
{[}\textless field block variable\textgreater{} \textless field
type\textgreater{} !BLOCK{]}
\end{quote}

Qualifiers:

\textbf{!SKIP }\\
The SKIP qualifier may be used to skip the first n lines of the pedigree
file. This is useful for ignoring a header.

IF(HPB)

\subsubsection{5.4. Genetic similarity from genomic data}\label{Gene29}

\paragraph{5.4.1. Genetic similarity from genomic data with unknown
pedigree}\label{Gene30}

\subparagraph{5.4.1.1. General}\label{Gene31}

In an evaluation with genomic data only, all individuals are genotyped.
Genomic data can be used to calculate a genomic relationship matrix. The
inverse of this matrix is used to include genetic similarity between
individuals. A range of inverse genomic relationship matrix files can be
created by MiXBLUP to be incorporated in the evaluation (chapter
\hyperref[Gene32]{5.4.1.2.}). The marker effect solutions can be
estimated afterwards (chapter \hyperref[Gene53]{5.4.4.}).

An alternative method to estimate genomic breeding values is to model
the direct genetic effect with a random regression on number of copies
of a SNP allele for a large number of loci (chapter
\hyperref[Gene33]{5.4.1.3.} and chapter \hyperref[Gene34]{5.4.1.4.}).
Direct genomic values for genotyped individuals without data can be
estimated afterwards from the marker effect solutions.

\subparagraph{5.4.1.2. Syntax of using GBLUP}\label{Gene32}

\begin{quote}
ERMFILE \textless Name file with genetic markers\textgreater{}
!CONSTRUCT Ginv\\
\textless animal ID\textgreater{} \textless field type\textgreater{}\\
!METHOD \textless Yang, VanRaden or VanRaden2\textgreater{} (optional;
default VanRaden2)\\
!ALFREQ \textless file name\textgreater{} (optional; default calculated
from data)\\
!CROSSBRED \textless number of breeds\textgreater{} \textless file
name\textgreater{} (optional; default single breed)\\
!INFORMATIVE (optional; default is to use all SNPs)\\
!DENSE \textless field number\textgreater{} (optional; default string of
markers in free format, field number is column with dense genotypes)\\
!MAF \textless minimum allele frequency\textgreater{} (optional; default
is 0.005)\\
!NUMPROC \textless number of processors to be used by
calc\_grm\textgreater{} (optional; default is 1)\\
!SKIP \textless n lines\textgreater{} (optional; default is reading all
lines)\\
!GFROMDISK (optional; default is to store relationship matrix in memory
during solving)
\end{quote}

Qualifiers:

\textbf{!CONSTRUCT Ginv}\\
The !CONSTRUCT qualifier is optional and indicates that the external
relationship matrix has not been calculated yet and needs to be
calculated in the MiXBLUP parser. For a GBLUP analysis, the argument of
!CONSTRUCT is Ginv, for an inverse genomic relationship matrix.

\textbf{!METHOD \textless Yang, VanRaden or VanRaden2\textgreater{}}\\
The !METHOD qualifier is optional and specifies whether the method of
Yang (Nat Genet 42:565-569) or the method of VanRaden (J Dairy Sci 91:
4414-4423) is used. The VanRaden2 method is the default.

\textbf{!ALFREQ \textless file name\textgreater{}}\\
The !ALFREQ qualifier is optional and allows the use of pre-defined
allele frequencies per locus from the file specified. By default, the
base population allele frequencies are calculated from the data.

\textbf{!CROSSBRED \textless number of breeds\textgreater{}
\textless name file with breed composition per animal\textgreater{}}\\
The !CROSSBRED qualifier is optional and can be used for multi-breed
analyses that may or may not include crossbred animals. There are two
arguments. The one argument is the number of pure breeds in the
analysis. The other argument is the name of the file with the breed
composition of the animals in the genetic marker file. The !CROSSBRED
option will consider relationships between all animals, regardless of
their breed composition, using for each animal allele frequencies that
are specific for their breed composition.

\textbf{!INFORMATIVE}\\
The !INFORMATIVE qualifier is used to include only genetic markers with
all three genotypes present in the population of genotyped individuals.
The default is to include all genetic markers with more than just one
allele in the data.

\textbf{!DENSE {[}\textless field number of dense
column\textgreater{]}}\\
The !DENSE qualifier must be specified if the genetic marker data is
presented as a sequence of genetic markers without spaces. If the dense
column is not the second field in the record, the field number of the
dense column needs to be specified after the qualifier, for example
!DENSE 4. If !DENSE is not specified for a file with dense genetic
marker data, MiXBLUP will give a column-width error, as it attempts to
read the dense genetic markers as a single column. By default, MiXBLUP
expects space-separated genetic markers.

\textbf{!MAF }\\
The qualifier !MAF is optional and is used to set the minimum allele
frequency of genetic markers to be included in the analysis. The default
value is 0.005. Use !MAF 0.0 to include all SNP with more than one
allele in the data.

\textbf{!NUMPROC}\\
The !NUMPROC qualifier can be used to specify the number of threads to
be used by calc\_grm.

\textbf{!GFROMDISK}\\
The !GFROMDISK qualifier instructs the solver to read the inverse
genomic relationship matrix from disk during solving. This was the only
option in previous versions of MiXBLUP. The default is to keep this
matrix in memory, which is more demanding for memory requirement, but it
saves the time to read this matrix every iteration. It is specified in
the SOLVING section of the MiXBLUP instruction file.

\subparagraph{5.4.1.3. Syntax of using SNPBLUP for the MiX99
solver}\label{Gene33}

\begin{quote}
SNPFILE {[}!CENTER{]} {[}!NOIMPUTE{]} {[}!MISSCOMB 0.01{]}
{[}!MISSPERLOC 0.01{]} {[}!NOPRUNE{]} {[}!CALCSNPVAR{]}
{[}!MINGENFREQ{]} {[}!GBSORTSNP {]} {[}!SAMEORDER{]}\\
\textless field animal\textgreater{} \textless field type I or
A\textgreater{}\\
SNP01 \textless file name SNP01\textgreater{} !REGTYPE R {[}!IDCOL
\textless n\textgreater{]} {[}!STARTCOV \textless n\textgreater{]}
{[}!LASTCOV \textless n\textgreater{]}\\
SNP02 \textless file name SNP02\textgreater{} !REGTYPE R {[}!IDCOL
\textless n\textgreater{]} {[}!STARTCOV \textless n\textgreater{]}
{[}!LASTCOV \textless n\textgreater{]}\\
\textless\ldots\textgreater{}\\
SNP99 \textless file name SNP99\textgreater{} !REGTYPE R {[}!IDCOL
\textless n\textgreater{]} {[}!STARTCOV \textless n\textgreater{]}
{[}!LASTCOV \textless n\textgreater{]}\\
{[}SNPPARFILE (required only for !REGTYPE H or for !REGTYPE R if
!CALCSNPVAR is not specified)\\
SNP01 \textless file name SNP01\textgreater{}\\
SNP02 \textless file name SNP02\textgreater{}\\
\textless\ldots\textgreater{}\\
SNP99 \textless file name SNP99\textgreater{]}\\
MODEL\\
\textless trait\textgreater{} \textasciitilde{}
\textless fixed\textgreater{} !RANDOM SNP(1,2,8..15,23) \# so no need
for G(\ldots) in the model
\end{quote}

Sections:

\textbf{SNPFILE}\\
The SNPFILE section specifies the name of one or more SNP covariate
files and its attributes, such as column numbers, dense or
space-separated SNPs and whether one variance for all SNPs is used or an
individual variance for each SNP. The section also has a number of
qualifiers that apply to all SNP covariate files.

\textbf{SNPPARFILE}\\
The SNPPARFILE section specifies the name of a parameter file for each
SNP covariate file for which the SNP covariates are fitted as a random
regression (so !REGTYPE is either `r' or `h'). The SNPPARFILE section
does not have any associated qualifiers. The lines of the SNPPARFILE
section each contain two columns. The first column is the label that
links the parameter file to the SNP covariate file. The second column is
the name of the file.

Qualifiers:

The file-independent qualifiers of SNPFILE are typically specified on
the first line of the SNPFILE section. These are:\\
\textbf{!CENTER}\\
The !CENTER qualifier is optional and scales all SNP's to a mean of 0
and standard deviation of 1. For details, see Stranden and Christensen
(2011). Centring the SNPs affects the fixed effect solutions, but not
the SNP effect solutions. Centering may enhance convergence of the PCG
iteration.

\textbf{!MISSCOMB \textless maximum fraction of SNPs
missing\textgreater{}}\\
The MISSCOMB qualifier is optional and can be used to specify the
tolerance level of missing combinations of animal and SNP. Above the
tolerance level, a warning is printed that the analysis may not yield
meaningful results, but the analysis continues. If !MISSCOMB is not
specified, the tolerance level is 0.001 of all combinations of animal
and SNP marker.

\textbf{!MISSPERLOC \textless maximum fraction of SNPs missing per
locus\textgreater{}}\\
The MISSPERLOC qualifier is optional. It specifies the tolerance level
of missing SNPs per locus. Loci with too many missing SNPs are written
to CheckDataSNP.log and a warning is printed. If !MISSPERLOC is not
specified, the tolerance level is 0.05 of all genotypes for the locus
(call rate of 95\%).

\textbf{!NOIMPUTE}\\
The !NOIMPUTE qualifier can be used to avoid automatic imputation of
missing SNPs with the average SNP value of the locus. If !NOIMPUTE is
specified, then animals with one or more missing SNPs get a genomic
breeding value of -99999 in the solanigen.txt file. If !NOIMPUTE is not
specified, then the average SNP value of the locus is used in the
calculation of the genomic breeding value.

\textbf{!NOCHECK}\\
The !NOCHECK qualifier can be used omit any imputing, pruning, centring
or checks of SNP covariates, which must be on the scale 0 to 2.

\textbf{!NOPRUNE}\\
The !NOPRUNE qualifier can be used omit the verification that all SNPs
are informative and the exclusion of non- or less-informative SNPs.

\textbf{!MINGENFREQ}\\
The !MINGENFREQ qualifier is optional and can be used to vary the
definition of a less-informative SNP. If the frequency of the minor SNP
genotype is below the threshold, it is considered to be less-informative
and it will be excluded from the analysis, unless !NOPRUNE has been
specified. The default threshold is 0, which does not remove any SNP
markers.

\textbf{!CALCSNPVAR}\\
The optional !CALCSNPVAR qualifier can be used to calculate the SNP
variance from the direct genetic variance in the parameter file
specified in the PARFILE section. The CALCSNPVAR qualifier must not be
specified if one or more SNP files have SNP-specific variances (!REGTYPE
H). If one or more SNP files are fixed (!REGTYPE F), the SNP variance is
calculated using the remaining SNP files with !REGTYPE R. When
!CALCSNPVAR is specified, the SNPPARFILE section is ignored.

\textbf{!GBSORTSNP \textless amount of memory in Gb\textgreater{}}\\
The qualifier !GbSortSNP only applies to the MiX99 solver. It is
optional and can be used to control the use of memory for sorting SNP
covariate records in the order of the data file. SNP covariate records
are sorted in blocks of records. !GBSORTSNP determines the size of such
a block of records to be sorted simultaneously. The default allocation
is 16 Gb. The number of blocks of records is the number of times a SNP
covariate file has to be read, so a small memory allocation increase the
time needed to sort the SNP covariates.

\textbf{!SAMEORDER}\\
The qualifier !SameOrder only applies to the MiX99 solver. If there are
multiple SNP covariate files and records in each file are in the same
order of individual ID, then !SAMEORDER can be used. MiXBLUP will sort
all SNP covariate files simultaneously. Note that memory allocated with
!GBSORTSNP is now used for multiple SNP covariate files instead of a
single file.

The second line of the SNPFILE section specifies the animal ID code that
can be used to link the data and the SNP covariate files.\\
The following lines each specify a SNP covariate file. Each line starts
with a label. The label links the SNP covariate file to the
corresponding SNP parameter file. The label must have the form `SNPxx'
where xx is a number between 01 and 99. The second field contains the
name of the SNP covariate file. The additional fields contain one or
more file-specific qualifiers. These are:

\textbf{!REGTYPE}\\
The file-specification line must contain the !REGTYPE qualifier. It
specifies how the covariates in the file are fitted in the model. If `f'
is specified, the covariates in the file are fitted as a fixed
regression. Covariates fitted as a fixed effect do not have a variance
associated with it, so it is not necessary to specify a parameter file
in the SNPPARFILE section. If it is present, it is ignored. If `r' is
specified, the covariates in the file are fitted as a random regression
with a single variance for all covariates in the file. The variance is
specified in the corresponding parameter file in the SNPPARFILE section.
If `h' is specified, the covariates in the file are fitted as a random
regression, each with their own variance. The covariate-specific
variances are specified in the corresponding parameter file in the
SNPPARFILE section.

\textbf{!IDCOL}\\
The !IDCOL qualifier is optional and specifies which field in the SNP
covariate file contains the ID of the genotyped animal. If it is
omitted, it is assumed that the ID is in the first field of the record
(so the default is !IDCOL 1).

\textbf{!STARTCOV}\\
The !STARTCOV qualifier is optional and specifies which field contains
the first SNP covariate. If it is omitted, it is assumed that the SNP
covariates start in the second field of the record (so !STARTCOV 2).

\textbf{!LASTCOV}\\
The !LASTCOV qualifier only applies to the MiX99 solver. It is optional
and specifies which field contains the last SNP covariate of the file to
include in the model. If it is omitted, it is assumed that all fields
after the first SNP covariate contain SNP covariates to include in the
model.

\textbf{SNP(\ldots)}\\
The SNP function only applies to the MiX99 solver. It can be used in the
MODEL section to specify which SNP covariate files should be fitted in
the model of a trait. If a SNP covariate file is specified, then all
specified SNP covariates in the file will be fitted. The number in the
SNP(\ldots) function links to the number in the label of the SNP
covariate file. The numbers may be specified individually as (1, 2, 3,
4) or as a range, indicated by two subsequent full stops, for example
(1..4), or a combination of both. When SNP(\ldots) is specified, it is
not necessary to specify the G(\ldots) function to specify a genetic
effect, but it is possible, for example to specify a maternal genetic
effect. Despite what the use of SNP(\ldots) in the MODEL section may
suggest, all SNP covariate files used in any trait are fitted for all
traits.

\subparagraph{5.4.1.4. Syntax of using SNPBLUP for hpblup
solver}\label{Gene34}

\begin{quote}
SNPFILE {[}!CENTER{]} {[}!NOIMPUTE{]} {[}!MISSCOMB 0.01{]}
{[}!MISSPERLOC 0.01 {[}!NOPRUNE{]} {[}!CALCSNPVAR{]} {[}!MINGENFREQ{]}\\
\textless field animal\textgreater{} \textless field type I or
A\textgreater{}\\
SNP02 \textless file name SNP02\textgreater{} !REGTYPE R {[}!IDCOL 1{]}
{[}!STARTCOV 2{]} {[}!PLINK{]}\\
{[}SNPPARFILE (required only for !REGTYPE H or for !REGTYPE R if
!CALCSNPVAR is not specified)\\
SNP02 {]}\\
MODEL\\
\textless trait\textgreater{} \textasciitilde{}
\textless fixed\textgreater{} !RANDOM hpSNP(2,\textless field
animal\textgreater) {[}hpSNP(2,\textless field dam\textgreater){]}
\end{quote}

Qualifiers:

Please note: the qualifiers !GbSortSNP, !SameOrder and !LastCov have no
effect when using the hpblup solver.

\textbf{hpSNP(\textless label number\textgreater,\textless index
field\textgreater)}\\
The hpSNP function is used to specify which SNP covariate files should
be fitted in the model of a trait. Unlike for the MiX99 solver, the SNP
covariate file is not fitted automatically for all traits. The first
parameter of the hpSNP function is the number in the label of the SNP
covariate file. The second parameter is the index field in the data
file. Every combination of label and index field requires a separate
hpSNP function in the model of a trait.

\paragraph{5.4.2. Genetic similarity from pedigree and genomic
data}\label{Gene35}

\subparagraph{5.4.2.1. General}\label{Gene36}

If pedigree data is available or if some individuals are not genotyped,
a single-step genomic BLUP model can be used. There are three
approaches. In the first approach, an inverse genomic relationship
matrix is used (ssGBLUP), either explicitly or implicitly in decomposed
form. The inverse genomic and pedigree relation matrices are blended.
There are various ways to do this (Appendix 3). If the number of
genotyped individuals is below 40,000, it is recommended to use the full
inverse of weighted G (chapter \hyperref[Gene37]{5.4.2.2.}. If the
number of genotyped individuals is higher, it is recommended to use the
component-wise Ta decomposition of G (ssGTacBLUP; chapter
\hyperref[Gene46]{5.4.2.3.}. The second approach is mathematically
equivalent to ssGBLUP and fits every SNP marker as a covariate. The
method implemented was developed by Liu et al.~(2014) and contains SNP
covariates and a residual polygenic effect. A genomic estimated breeding
value (GEBV) is calculated for all individuals, which is the sum of the
direct genomic value, calculated from the SNP effects, and the residual
polygenic breeding value. This method is called single-step SNP BLUP
(ssSNPBLUP; chapter \hyperref[Gene47]{5.4.2.4.}). If the number of
animals with both a phenotype and a genotype is sufficiently large for
all traits, then SNP effect estimates are quite stable for a number of
subsequent evaluations. The third approach uses previously estimated SNP
effects to calculate direct genomic values (DGV) for genotyped animals.
These DGV are fitted as prior information in a pedigree BLUP evaluation
(DGV-PBLUP; chapter \hyperref[Gene48]{5.4.2.5.} and
\hyperref[Gene49]{5.4.2.6.}). DGV-PBLUP can be used with ssGTacBLUP and
ssSNPBLUP. A full genomic evaluation is needed periodically to
re-estimate SNP effects. DGV-PBLUP provides a substantial reduction in
runtime and computing resources. The number of subsequent evaluations
for which DGV-PBLUP can be used, depends on multiple factors, such as
the size of the phenotype and genotype datasets. DGV-PBLUP is only
available for the hpblup solver.

\subparagraph{5.4.2.2. Syntax of SSGBLUP: single-step genomic BLUP with
full inverse of a weighted G}\label{Gene37}

\begin{quote}
ERMFILE \textless Name file with genetic markers\textgreater{}
!CONSTRUCT SSmat\\
\textless animal ID\textgreater{} \textless field type\textgreater{}
!LAMBDA \textless weighting factor G matrix, 0.0-1.0\textgreater{}
(optional; default 1.0)\\
!ALPHA \textless weighting factor, 0.0-1.0\textgreater{} (optional;
default is 0.95)\\
!BETA \textless weighting factor, 0.0-1.0\textgreater{} (optional;
default is 1.0-ALPHA)\\
!OMEGA \textless weighting factor, 0.0-1.0\textgreater{} (optional;
default is LAMBDA)\\
!SINGLESTEP\\
{[}INBRFILE \textless inbreeding coefficient file\textgreater{} !IDCOL
\textless number\textgreater{}
!INBRCOL~\textless number\textgreater{]}\\
PEDFILE \textless pedigree file\textgreater{} {[}!CALCINBR{]}\\
\textless individual ID\textgreater{} \textless field
type\textgreater{}\\
\textless sire ID\textgreater{} \textless field type\textgreater{}\\
\textless dam ID\textgreater{} \textless field type\textgreater{}
\end{quote}

Any qualifier of ERMFILE described in chapter
\hyperref[Gene31]{5.4.1.1.} can also be used for a weighted inverse
genomic relationship matrix. Specific qualifiers:

\textbf{!CONSTRUCT SSmat}\\
The !CONSTRUCT qualifier indicates that the external relationship matrix
has not been calculated yet and needs to be calculated in the MiXBLUP
parser. For a ssGBLUP analysis with a weighted inverse genomic
relationship matrix, the argument of !CONSTRUCT is SSmat.

\textbf{!LAMBDA }\\
\textbf{!ALPHA }\\
\textbf{!BETA }\\
\textbf{!OMEGA }\\
The !LAMBDA, !ALPHA, !BETA and !OMEGA qualifiers are the fudge
parameters suggested by Misztal and coworkers. They are optional and can
be used to give more weight to numerator relationship matrix (A-1) in
the construction of the blended matrix (H-1):

\pandocbounded{\includegraphics[keepaspectratio]{https://raw.githubusercontent.com/JantN1966/MasterManual/main/Images/GenSim04.jpg}}\\

By default, lambda is set to 1, omega to lambda, alpha to 0.95 and beta
to 0.05.

\textbf{!SINGLESTEP}\\
The qualifier !SINGLESTEP is used to indicate that the MiXBLUP kernel
should calculate elements of the H-inverse from a G-inverse, the
pedigree file and a file with inbreeding coefficients. This option
requires a matrix that is set up using !CONSTRUCT with argument SSmat.

\subparagraph{5.4.2.3. Syntax of ssGTacBLUP: single-step GBLUP with
component-wise Ta decomposition of a weighted G}\label{Gene46}

\begin{quote}
ERMFILE \textless Name file with genetic markers\textgreater{}
!CONSTRUCT ssMat !Tac !SingleStep\\
\textless animal ID\textgreater{} \textless field type\textgreater{}
\end{quote}

Any qualifier of ERMFILE described in chapter
\hyperref[Gene32]{5.4.1.2.} and chapter \hyperref[Gene37]{5.4.2.2.} can
also be used for a decomposition of a weighted inverse genomic
relationship matrix. Specific qualifiers:

\textbf{!CONSTRUCT SSmat}\\
The !CONSTRUCT qualifier indicates that the external relationship matrix
has not been calculated yet and needs to be calculated in the parser.

\textbf{!Tac}\\
For a ssGBLUP analysis with a decomposition of a weighted inverse
genomic relationship matrix, add the qualifier !Tac. Component-wise Ta
decomposition of a weighted G is only more efficient if the number of
genotyped individuals is substantially larger than the number of SNP
markers. If the number of genotyped individuals is between 40,000 and
say 1.5 times the number of SNP markers, then ordinary Ta decomposition
using the qualifier !Ta instead of !Tac may be more efficient.

\subparagraph{5.4.2.4. Syntax of ssSNPBLUP: single-step BLUP using SNP
genotypes as covariates}\label{Gene47}

\begin{quote}
PEDFILE \textless file name pedigree file\textgreater{} !CalcInbr !Beta
\textless value\textgreater{}\\
\textless field individual ID\textgreater{} \textless field type I or
A\textgreater{}\\
\textless field sire ID\textgreater{} \textless field type I or
A\textgreater{}\\
\textless field dam ID\textgreater{} \textless field type I or
A\textgreater{}\\
\textless\ldots\textgreater{}\\
SNPFILE {[}!CalcSNPvar{]} !NoCheck !NoPrune\\
\textless field individual ID\textgreater{} \textless field type I or
A\textgreater{}\\
SNP01 \textless file name SNP01\textgreater{} !REGTYPE R {[}!IDCOL
\textless value\textgreater{]} {[}!STARTCOV
\textless value\textgreater{]}\\
\textless\ldots\textgreater{}\\
MODEL\\
\textless trait\textgreater{} \textasciitilde{}
\textless\ldots\textgreater{} !RANDOM hpSNP(1,\textless individual ID
field\textgreater) {[}hpSNP(1,\textless dam ID field\textgreater){]}
G(\textless field individual ID\textgreater{[},\textless dam ID
field\textgreater{]})
\end{quote}

Qualifiers:

\textbf{!Beta }\\
The qualifier !Beta is optional and can be used to specify the fraction
of residual polygenic variation. It must not be zero and should be at
least 0.05, which is the default value if !Beta is not specified. The
qualifier !Beta can be specified in the PEDFILE or the SOLVING section
for ssSNPBLUP. The qualifier !Alpha cannot be specified for ssSNPBLUP
but is calculated as 1 -- Beta.

\textbf{!CalcSNPvar}\\
The optional !CALCSNPVAR qualifier can be used to calculate the SNP
variance from the direct genetic variance in the parameter file
specified in the PARFILE section. The CALCSNPVAR qualifier must not be
specified if one or more SNP files have SNP-specific variances (!REGTYPE
H). If one or more SNP files are fixed (!REGTYPE F), the SNP variance is
calculated using the remaining SNP files with !REGTYPE R. When
!CALCSNPVAR is specified, the SNPPARFILE section is ignored.

\textbf{!NoCheck, !NoPrune}\\
It is advised to use !NoCheck and !NoPrune if GEBV of young individuals
with genotype, but not phenotype or progeny, are calculated in a
separate analysis. Not using these options may lead to non-compatible
SNP covariate files of the main evaluation and the evaluation of young
genotyped individuals, which will result in an error.

\subparagraph{5.4.2.5. Syntax of DGV-Pedigree BLUP from Tac: using
previously estimated SNP effects as prior information in pedigree
BLUP}\label{Gene48}

\begin{quote}
PEDFILE \textless file name pedigree file\textgreater{} !CalcInbr\\
\textless field individual ID\textgreater{} \textless field type I or
A\textgreater{}\\
\textless field sire ID\textgreater{} \textless field type I or
A\textgreater{}\\
\textless field dam ID\textgreater{} \textless field type I or
A\textgreater{}\\
\textless\ldots\textgreater{}\\
ERMFILE \textless genotype file in plink format\textgreater{} !Plink
!CONSTRUCT ssMat !SingleStep {[}!Tac or !Ta{]} \textless field
individual ID\textgreater{} \textless field type I or A\textgreater{}\\
!METHOD VanRaden\\
!MAF 0.000\\
!ALPHA 0.95\\
!DGVPBLUP \textless name file with SNP effect solutions from Tac/Ta
evaluation\textgreater{}\\
!AlFreq \textless name file with estimated allele frequencies from
Tac/Ta evaluation\textgreater{}\\
MODEL\\
\textless trait\textgreater{} \textasciitilde{}
\textless\ldots\textgreater{} !RANDOM G(\textless field individual
ID\textgreater{[},\textless dam ID field\textgreater{]})\\
SOLVING\\
!hpblup
\end{quote}

Qualifiers:

\textbf{!DGVPBLUP }\\
The qualifier !DGVPBLUP can be used to specify that old solutions for
SNP effects should be used as prior information. For a component-wise Ta
decomposition, SNP effect estimates are written to a file called
snpeffects\_ta.dat. The specified file should be a copy of this file.

\textbf{!AlFreq }\\
The qualifier !AlFreq is mandatory if !DGVPBLUP is used. The file
specified should come from the same evaluation as the SNP effect
estimates. Allele frequencies are written to calculated\_all\_freq.dat.
The specified file should be a copy of this file.

\textbf{!MAF 0.000}\\
It is advised to use a minor allele frequency of 0.0\%. Not using this
option may lead to non-compatible SNP and allele frequency files, which
will result in an error.

\subparagraph{5.4.2.6. Syntax of DGV-Pedigree BLUP from ssSNPBLUP: using
previously estimated SNP effects as prior information in pedigree
BLUP}\label{Gene49}

\begin{quote}
PEDFILE \textless file name pedigree file\textgreater{} !CalcInbr !Beta
\textless value\textgreater{}\\
\textless field individual ID\textgreater{} \textless field type I or
A\textgreater{}\\
\textless field sire ID\textgreater{} \textless field type I or
A\textgreater{}\\
\textless field dam ID\textgreater{} \textless field type I or
A\textgreater{}\\
\textless\ldots\textgreater{}\\
SNPFILE {[}!CalcSNPvar{]} !NoCheck !NoPrune !DGVPBLUP \textless name
file with SNP effects solutions\textgreater{} !AlFreq \textless name
file with estimated allele frequencies\textgreater{}\\
\textless field individual ID\textgreater{} \textless field type I or
A\textgreater{}\\
SNP01 \textless file name SNP01\textgreater{} !REGTYPE R {[}!IDCOL
\textless value\textgreater{]} {[}!STARTCOV
\textless value\textgreater{]}\\
\textless\ldots\textgreater{}\\
MODEL\\
\textless trait\textgreater{} \textasciitilde{}
\textless\ldots\textgreater{} !RANDOM hpSNP(1,\textless individual ID
field\textgreater) {[}hpSNP(1,\textless dam ID field\textgreater){]}
G(\textless field individual ID\textgreater{[},\textless dam ID
field\textgreater{]})
\end{quote}

Qualifiers:

\textbf{!DGVPBLUP \textless name of file with estimated SNP
effects\textgreater{}}\\
The qualifier !DGVPBLUP can be used to specify that old solutions for
SNP effects should be used as prior information. For a single-step
SNPBLUP evaluation, SNP effect estimates are written to
solutions\_mixblup.dat. The file Solreg\_mat.txt can also be used and is
considerably smaller. The specified file should be a copy of one of
these file.

\textbf{!AlFreq }\\
The qualifier !AlFreq is mandatory if !DGVPBLUP is used. The file
specified should come from the same evaluation as the SNP effect
estimates. Allele frequencies are written to calculated\_all\_freq.dat.
The specified file should be a copy of this file.

\textbf{!MAF 0.000 !NoCheck !NoPrune}\\
It is advised to use a minor allele frequency of 0.0\%, and the !NoCheck
and !NoPrune options. Not using these options may lead to non-compatible
SNP and allele frequency files, which will result in an error.

\paragraph{5.4.3. Modelling a genetic difference between genotyped and
non-genotyped individuals with the hpblup solver}\label{Gene50}

\subparagraph{5.4.3.1. General}\label{Gene51}

The group of genotyped individuals without genotyped ancestors may not
be representative for the base population in the evaluation, which is
the group of individuals without known parents. As a result, genomic
estimated breeding values are biased in the presence of selection and/or
selective genotyping (Vitezica et al., 2011). When genomic relationships
were shifted by a constant in their study, the single-step method was
unbiased and the most accurate of the methods compared. For ssGBLUP
using either a full or APY inverse of G, this constant is automatically
applied by calc\_grm, unless the !NoScale and !NoReg options are used.
For ssGTacBLUP and ssSNPBLUP, this constant can be modelled with a
so-called J factor that quantifies for any individual the pedigree
relationship with genotyped individuals. For genotyped individuals, the
J factor is set to -1. For base animals, the J factor is initialised at
0. For ancestors of genotyped individuals, the J factor (Jn) is
estimated using Ann Jn = -Ang Jg, where Ang and Ann are partitions of
A-inverse of non-genotyped ancestors by genotyped individuals and
non-genotyped ancestors by non-genotyped ancestors, respectively, and
Jg, as J factor for genotyped individuals, is a vector containing -1 for
all. For all remaining individuals, the J factor is the average of the J
factor of the parents. It is possible to do this calculation in the
parser. The regression coefficient of each trait on the J factor is
estimated on all available records. The impact of a priori assuming that
estimates originate from a normal distribution with a given variance
(i.e.~fitting it as a random effect) is very limited. It is therefore
recommended to fit J-factor covariate as a fixed effect. The individual
correction for bias, calculated for each trait as regression coefficient
times J-factor covariate, is added to the GEBV of the trait for each
animal. Note that the J factor will change for non-genotyped individuals
if new genotyped individuals are added to an evaluation. The calculation
of J-factor covariates requires the genotype file and the pedigree file.
An existing J-factor covariate file needs to contain all animals in the
pedigree.

\subparagraph{5.4.3.2. Syntax of fitting J factor
covariate}\label{Gene52}

\begin{quote}
PEDFILE \textless file name pedigree file\textgreater{} !CalcInbr !Beta
\textless value\textgreater{} !MakeJcov\\
\textless field individual ID\textgreater{} \textless field type I or
A\textgreater{}\\
\textless field sire ID\textgreater{} \textless field type I or
A\textgreater{}\\
\textless field dam ID\textgreater{} \textless field type I or
A\textgreater{}\\
\textless\ldots\textgreater{} REGFILE\\
\textless field individual ID\textgreater{} \textless field type I or
A\textgreater{}\\
REG01 !REGTYPE F !Jcov\\
\textless\ldots\textgreater{}\\
MODEL\\
\textless trait\textgreater{} \textasciitilde{}
\textless\ldots\textgreater{} hpReg(1, \textless field individual
ID\textgreater) !RANDOM \textless\ldots\textgreater{}
\end{quote}

Qualifiers:

\textbf{!MakeJcov}\\
The qualifier !MakeJcov is optional and is used to specify the
calculation of J-factor covariates.

\textbf{!Jcov}\\
The qualifier !Jcov marks the covariate file that contains the J-factor
covariates. If !MakeJcov is specified, it is not necessary to specify a
file name.

\textbf{hpReg(\textless number in label of covariate file\textgreater,
\textless field individual ID\textgreater)}\\
The hpReg function is used to fit the J-factor covariate file in the
model of each trait. It is recommended that it be fitted as a fixed
effect, by specifying the RegType F qualifier in the REGFILE section and
by specifying the hpReg function before the !Random qualifier on each
line in the MODEL section.

\paragraph{5.4.4. Obtaining solutions for genetic marker effects with
the hpblup solver}\label{Gene53}

Genetic marker effect solutions are provided with the recommended
methods of a genomic evaluation, ssGTaBLUP, ssGTacBLUP and ssSNPBLUP.
Solutions are present in the file Solreg\_mat.txt, which contains just
the centred genetic marker effect solutions. The order of fields in
Solreg\_mat.txt is Trait - Matrix - Covariate - EffectID - Solution -
Matrix name. Other supported methods of a genomic evaluation (Appendix
XX) require back-solving to obtain genetic marker effect solutions.
Back-solving needs to be specified at the start of a genomic evaluation.
It cannot be done retrospectively. The option !BackSolve in the SOLVING
section can be used for this purpose. Back-solving only yields
approximate genetic marker effect solutions, so ssGTacBLUP and ssSNPBLUP
are the two recommended methods to obtain genetic marker effect
solutions.

\subsubsection{5.5. External genetic relationship matrix}\label{Gene54}

\paragraph{5.5.1. General}\label{Gene55}

A range of inverse genetic relationship matrix files can be created by
MiXBLUP explicitly or are implicitly incorporated in the analysis. It is
also possible to use a previously created inverse genetic relationship
matrix or one that as yet cannot be created by MiXBLUP itself. It is not
possible to calculate reliabilities with MiXBLUP when using a full
external relationship matrix.

\paragraph{5.5.2. Recommended file format}\label{Gene56}

An external relationship matrix must contain the original ID of each
individual. This can be a decoded matrix from a previous evaluation or
one from an external source. The recommended file format for a dense
inverse relationship matrix (most matrix elements are non-zero) is the
dense or lower-triangular matrix format. The first line of the dense
format contains the number of genotyped individuals and the number of
individuals in the core in case of an APY-inverse or 0 otherwise. The
second line contains the original ID of each individual separated by one
or more spaces. The next lines contain all elements of the
lower-triangular part of the matrix up to and including the diagonal
element. An external dense relationship matrix may also be provided in
stream binary format, recognized from the .stream filename extension.
The recommended file format for a sparse inverse relationship matrix is
the sparse or I-J-value matrix format. Each line consists at least of
three fields: original individual ID of row, original individual ID of
column, matrix element. Any other fields on the line are ignored.

\paragraph{5.5.3. Syntax of using an external relationship
matrix}\label{Gene57}

\begin{quote}
ERMFILE \textless external relationship matrix file\textgreater{}
{[}!SKIP \textless n lines\textgreater{]} {[}!ASIS{]} {[}!NOORIG{]}
\textless field individual ID\textgreater{} \textless field
type\textgreater{}
\end{quote}

Qualifier:

\textbf{!SKIP \textless n lines\textgreater{}}\\
The optional !SKIP qualifier may be used to skip the first n lines of
the external relationship matrix file. This is useful for ignoring a
header line.

\textbf{!ASIS}\\
The !AsIs qualifier is optional. It is used to write the external
inverse relationship matrix to the kernel without any checks or sorting.
This can be specified if the external relationship matrix file is known
to be correct, for example because it was created or checked by MiXBLUP
in a previous run. The !AsIs qualifier can only be used if the field
type of the individual ID is integer. It is ignored when individual ID
has alphanumerical field type.

\textbf{!NOORIG}\\
If the MiX99 solver is used, MiXBLUP converts the coded external inverse
relationship matrix file back to the original individual IDs. If this
file is not needed for additional analyses, the !NoOrig qualifier can be
specified. Especially for very large analyses, the size of this file can
be substantial.

\emph{Table}. Decoded weighted genomic relationship matrices from
MiXBLUP evaluations

{\def\LTcaptype{none} % do not increment counter
\begin{longtable}[]{@{}
  >{\raggedright\arraybackslash}p{(\linewidth - 4\tabcolsep) * \real{0.3333}}
  >{\raggedright\arraybackslash}p{(\linewidth - 4\tabcolsep) * \real{0.3333}}
  >{\raggedright\arraybackslash}p{(\linewidth - 4\tabcolsep) * \real{0.3333}}@{}}
\toprule\noalign{}
\begin{minipage}[b]{\linewidth}\raggedright
Solver
\end{minipage} & \begin{minipage}[b]{\linewidth}\raggedright
Dense matrix format
\end{minipage} & \begin{minipage}[b]{\linewidth}\raggedright
Sparse matrix format
\end{minipage} \\
\midrule\noalign{}
\endhead
\bottomrule\noalign{}
\endlastfoot
MiX99 solver - integer ID & ExtRelMat\_tri.txt & ExtRelMat.txt \\
MiX99 solver - alphanumerical ID & ExtRelMatAlphaTri.stream &
ExtRelMatAlpha.stream \\
hpblup solver & Not decoded & Not decoded \\
\end{longtable}
}

\subsubsection{5.6. Genetic similarity in case of multiple breeds or
lines and crosses}\label{Gene58}

\paragraph{5.6.1. General}\label{Gene59}

There are several ways to correctly model that a genetic evaluation
consists of multiple breeds, lines or crosses (referred to as genetic
lines in this chapter). The effect of multiple genetic lines can be
fitted as a fixed effect in the model, it can be fitted as a separate
base population for each genetic line in the pedigree or allele
frequencies specific for the genetic lines may be provided.

\paragraph{5.6.2. Fixed effect in the model}\label{Gene60}

In case of the evaluation only containing purebred individuals, the
genetic line of the individual may be fitted as a fixed class effect. In
case of unstructured crossbreeding, it is recommended to fit a covariate
for each genetic line with the percentage of genes for that line, except
for one genetic line: a covariate of the most prevalent genetic line
should not be fitted in order to avoid singularity of the coefficient
matrix. Its effect is then included in the overall mean. In case of
purebred breeding and structured crossbreeding, either fixed class
effects can be used, or fixed genetic line covariates. See Chapter
\hyperref[Stat01]{7.} for the syntax of fitting fixed class effects or
fixed covariates. Note that fixed effect estimates of genetic line are
not automatically included in genetic effect solutions.

\paragraph{5.6.3. Base population for each genetic line}\label{Gene61}

An alternative way of fitting effects of genetic line is to assign
individuals in the base generation to a base population (i.e.~genetic
group) that is specific for a genetic line. See Chapter
\hyperref[Gene13]{5.1.5.} for the syntax of fitting multiple base
populations. Solutions of base populations are automatically added to
genetic effect solutions.

\paragraph{5.6.4. Allele frequencies specific for a genetic
line}\label{Gene62}

A breed composition file can be used to indicate that allele frequencies
should be estimated for each line in the breed composition file
separately. The breed composition file contains the original animal ID
in the first column and contains a number of additional columns that is
equal to the number of genetic lines specified. The breed composition
may be presented as a number, for example 4 (out of 8 or any other
number), as a percentage, for example 50, or as a fraction, for example
0.50. MiXBLUP converts the breed composition of an animal to the value
of one breed over the sum of values across breeds. For example in an
analysis with four breeds, animal X having 4 0 2 2 as the breed
composition will be converted to X 0.500 0.000 0.250 0.250. It is
therefore essential that the breed information is complete, so add a
column for `unknown or other', if necessary. All columns must be
separated by at least one space.

\pandocbounded{\includegraphics[keepaspectratio]{https://raw.githubusercontent.com/JantN1966/MasterManual/main/Images/GenSim06.jpg}}\\

\emph{Example}. Breed composition file with the percentage of four
breeds per animal

\pandocbounded{\includegraphics[keepaspectratio]{https://raw.githubusercontent.com/JantN1966/MasterManual/main/Images/GenSim07.jpg}}\\

\emph{Example}. Breed composition file in parts of one eighth of four
breeds per animal

\subsubsection{5.7. Non-additive genetic similarity}\label{Gene63}

Multiple genetic or genomic relationship matrices will be supported
shortly. This includes a genomic dominance or epistasis relationship
matrix. This section is pending the release of this feature.

\paragraph{5.7.1. General}\label{Gene64}

Systematic change of populations through genetic selection focuses on
additive genetic effects. In reality, there may be interactions between
alleles within a locus (dominance), interactions between loci
(epistasis), and expression that depends on environmental conditions.
Loss of heterozygosity because of inbreeding causes loss of dominance
effects and leads to inbreeding depression. Recovering from inbreeding
depression after crossbreeding is referred to as heterosis. Breaking up
of favorable epistatic effects after repeated crossbreeding is referred
to as recombination.

\paragraph{5.7.2. Expected heterosis and recombination in
crossbreds}\label{Gene65}

In case of unstructured crossbreeding, it is meaningful to correct for
expected heterosis and recombination if genotype information is not
available for a large part of the population. Both are a function of the
breed fractions of the parents.

Expected heterosis (Het) for genetic line A and B of a crossbred
individual is calculated as: \textgreater Het\textsubscript{A,B} =
p\textsubscript{sire,A}*p\textsubscript{dam,B} +
p\textsubscript{dam,A}*p\textsubscript{sire,B}

Expected recombination (Rec) for genetic line A and B of a crossbred
individual is calculated as: \textgreater Rec\textsubscript{A,B} =
p\textsubscript{sire,A}*p\textsubscript{sire,B} +
p\textsubscript{dam,A}*p\textsubscript{dam,B}

Here, p\textsubscript{sire,A} and p\textsubscript{dam,A} are the breed
fractions of genetic line A for the sire and the dam of the individual.
An expected heterosis and an expected recombination term should be
fitted for each combination of genetic lines present in the evaluation.
Expected heterosis and recombination can be fitted as fixed covariates
for each individual in the genetic evaluation. See Chapter
\hyperref[Stat01]{7.} for the syntax of fitting a fixed covariate.

\paragraph{5.7.3. Genomic dominance effects (to complete)}\label{Gene66}

\subparagraph{5.7.3.1. General}\label{Gene67}

Dominance effects can be estimated from heterozygous loci if all
individuals in the evaluation are genotyped. Heterozygosity can be used
to calculate an inverse dominance relationship matrix or create a
covariate. Genomic dominance can be fitted as an additional random
genetic correlated effect.

\subparagraph{5.7.3.2. Syntax}\label{Gene68}

\begin{quote}
ERMFILE \textless genotype file\textgreater{}\\
\textless ID field\textgreater{} I/A\\
!ConstructDom\\
CORRFILE\\
RCE01 !GenRCE !Dominance \# no file needed\\
\textless\ldots\textgreater{}\\
SOLVING\\
!BackSolveDom
\end{quote}

Qualifiers:

\textbf{!Dominance} The !Dominance qualifier is used to specify that a
random effect with correlated levels

\textbf{!ConstructDom} \textless\ldots\textgreater{}

\textbf{!BackSolveDom} \textless\ldots\textgreater{}

\subparagraph{5.7.3.3. Output files}\label{Gene69}

{\def\LTcaptype{none} % do not increment counter
\begin{longtable}[]{@{}ll@{}}
\toprule\noalign{}
Output file & Description \\
\midrule\noalign{}
\endhead
\bottomrule\noalign{}
\endlastfoot
& \\
\end{longtable}
}

\paragraph{5.7.4. Genomic epistasis effects (to complete)}\label{Gene70}

\subparagraph{5.7.4.1. General}\label{Gene71}

Modelling epistasis aims to estimate how pairs or groups of markers
interact to affect a phenotype, such as in additive by additive or
dominance by dominance interactions. The epistasis relationship matrix
is usually constructed using the element-wise product of the
(non-inverted) additive or dominance relationship matrices. The inverse
of the epistasis relationship matrix is used in the genetic evaluation.
Genomic epistasis is fitted as an additional random correlated genetic
effect.

\subparagraph{5.7.4.2. Syntax}\label{Gene72}

\begin{quote}
ERMFILE \textless genotype file\textgreater{}\\
\textless ID field\textgreater{} I/A\\
!ConstructEpist\\
CORRFILE\\
RCE01 !GenRCE !Epistasis \# no file needed \textless\ldots\textgreater{}
SOLVING
\end{quote}

Qualifiers:

\textbf{!Epistasis}\\
\textless\ldots\textgreater{}

\textbf{!ConstructEpist}\\
\textless\ldots\textgreater{}

\subparagraph{5.7.4.3. Output files}\label{Gene73}

{\def\LTcaptype{none} % do not increment counter
\begin{longtable}[]{@{}ll@{}}
\toprule\noalign{}
Output file & Description \\
\midrule\noalign{}
\endhead
\bottomrule\noalign{}
\endlastfoot
& \\
\end{longtable}
}

\subsubsection{5.8. Genetic similarity in case of polyploidy or mixed
ploidy}\label{Gene74}

\paragraph{5.8.1. General}\label{Gene75}

The MiXBLUP software was designed originally to support breeding value
estimation for diploid species. Although rare in mammals, a mix of
diploid and polyploid individuals can be found in species like
salamander, frog, trout, salmon, and various insect species. Support for
mixed ploidy is currently only available for expected genetic similarity
from pedigree. The underlying assumption is that founder animals in the
pedigree are from a single large population. Unknown parent groups are
currently not supported for polyploidy or mixed ploidy.

\paragraph{5.8.2. Required format of pedigree file}\label{Gene76}

A pedigree file in which some or all individuals are non-diploid,
contains at least seven fields. The first seven fields are defined as
described in the table below.

\emph{Table}. Description of fields in a pedigree record for a mixed
ploidy pedigree

{\def\LTcaptype{none} % do not increment counter
\begin{longtable}[]{@{}
  >{\raggedright\arraybackslash}p{(\linewidth - 2\tabcolsep) * \real{0.5000}}
  >{\raggedright\arraybackslash}p{(\linewidth - 2\tabcolsep) * \real{0.5000}}@{}}
\toprule\noalign{}
\begin{minipage}[b]{\linewidth}\raggedright
Field in pedigree record
\end{minipage} & \begin{minipage}[b]{\linewidth}\raggedright
Description
\end{minipage} \\
\midrule\noalign{}
\endhead
\bottomrule\noalign{}
\endlastfoot
1 & Individual ID \\
2 & Sire ID \\
3 & Dam ID \\
4 & Ploidy level of the gamete transmitted by the sire \\
5 & Ploidy level of the gamete transmitted by the dam \\
6 & Probability that two gametic genes, drawn at random from a randomly
selected locus, are inherited from a single parental chromosome
(sire) \\
7 & Probability that two gametic genes, drawn at random from a randomly
selected locus, are inherited from a single parental chromosome (dam) \\
\end{longtable}
}

Ploidy level of the gamete of a parent of a diploid individual is 1. In
that case, the probability that two gametic genes, drawn at random from
a randomly selected locus, are inherited from a single parental
chromosome is 0. For polyploid individuals, the latter probability is 0
if there is no identity by descent within a gamete, it is 1 if the
polyploid gamete was transmitted by a double-haploid individual, and it
is between 0 and 1 if there is recombination.

\paragraph{5.8.3. Syntax}\label{Gene77}

\begin{quote}
PEDFILE pedigree\_ploidy.txt !polyploid !CalcInbr\\
\textless field individual ID\textgreater{} \textless field type I or
A\textgreater{}\\
\textless field sire ID\textgreater{} \textless field type I or
A\textgreater{}\\
\textless field dam ID\textgreater{} \textless field type I or
A\textgreater{}\\
ploidy\_sire I\\
ploidy\_dam I\\
prob\_sire I\\
prob\_dam I
\end{quote}

Qualifier:

\textbf{!Polyploid}\\
The !Polyploid qualifier is used to specify that the pedigree contains
non-diploid individuals. In that case, MiXBLUP expects each pedigree
record to consist of at least seven fields instead of three fields.
Without !Polyploid, the four additional fields in the pedigree are
ignored. The use of !Polyploid results in an (implicit) inverse
\emph{pedigree} relationship matrix that takes ploidy level of
individuals into account. An inverse \emph{genomic} relationship matrix
that takes differences in ploidy into account is currently not
supported.

\hyperref[Tabl01]{Back to Table of Contents}

\newpage

\subsection{6. Components of variance and covariance among
traits}\label{Comp01}

Components of variance and covariance among traits are normally
specified in the general parameter file. Additional covariance
components for covariates in a covariate file need to be specified in
separately labelled parameter files. Heterogeneous residual variances
also need to be specified in a separate file. This chapter describes how
to specify components of variance and covariance among traits.

\subsubsection{6.1. General parameter file}\label{Comp02}

\paragraph{6.1.1. General}\label{Comp03}

The trait (co)variance components file contains the between-trait
variance-covariance matrices of any random effects in the statistical
model. There are two options for the format of the general parameter
file: (1) in lower-triangular-matrix form and (2) in sparse-matrix form.
It is strongly recommended to use the lower-triangular-matrix format.
The instruction file specifies the name of the trait (co)variance
components file. The trait (co) variance components file is located by
default in the work directory, but can be in another folder if specified
in the name of the file.

\paragraph{6.1.2. Input file in lower-triangular-matrix
format}\label{Comp04}

The lower-triangular-matrix form is the default option and strongly
recommended. In this form, the trait covariance components file can be
specified as a lower-triangular matrix using trait names to identify the
components. This is the most user-friendly way. The name of the random
effect is given at the top of the matrix and the names of the traits are
given at the start of each line of the matrix.

\begin{itemize}
\tightlist
\item
  The lower triangular matrices and the traits within a matrix can be
  specified in any order. It means that the order given in the MODEL
  section of the instruction file is not leading.
\item
  The number of traits in the matrices can be larger than the number of
  traits specified in the model section. Only the lines for which the
  name has been specified in the model section will be used.
\item
  The order of the column names must be the same as the order of row
  names, so variance components are on the diagonal.
\item
  Restriction: in case of a marker-assisted BLUP model with the use of
  haplotype variance-covariance matrices, each matrix needs to be named
  and numbered, e.g.~GIV1, GIV2, etc. The name GIV refers to the use of
  the General Inverse Variance (GIV) function in the model. The order of
  matrices must be the same as the order of haplotypes given in the
  model lines of the instruction file. See Example 5.4 in the Appendix.
\item
  For all direct and indirect genetic effects (e.g.~animal, dam, mate),
  it should be specified immediately after the trait name and within
  brackets whether it is the genetic variance of animal, dam or mate.
\item
  In case of non-genetic random regression, the name of the class effect
  is specified at the top of the matrix and a line for each combination
  of trait and the full random regression term in the model of the trait
  should be specified. The syntax in previous versions of MiXBLUP with a
  separate matrix for each random regression term is still supported,
  but not recommended, as it ignores covariance components between
  different random regression terms of the same trait.
\item
  If the model contains genetic random regression, then all fitted
  regression terms should be specified in the variance covariance table
  (e.g.~animal*covar1 and animal*covar2).
\item
  For the MiX99 solver
\item
  If a covariate table file is used for random regression, then the
  columns should be referred to as cvr00 for the first covariate column,
  cvr01 for the second column and so on. The name should be lowercase:
  the use of CVR00 will give an error.
\item
  In case of a social interaction model, with multiple mate effects in
  the model, the first group mate effect in the model should be
  specified (e.g.~mate1*mate1\_x, where mate1\_x is a covariate that
  indicates whether mate2 is a real (1) or a dummy (0) group mate).
\item
  For the hpblup solver
\item
  If a covariate table file is used for random regression, then the
  columns should be referred to as TABLE01\_00 for the first covariate
  column of the file labelled TABLE01, TABLE01\_01 for the second column
  and so on.
\item
  In case of a social interaction model, with multiple mate effects in
  the model, the group mate in the G(\textless\ldots\textgreater,
  LINK(\textless\ldots\textgreater)) function in the model should be
  specified (e.g.~mate1 for G(animal, LINK(mate1))).
\item
  In case of group phenotypes, the effect in the LINK function in the
  model should be specified in the corresponding trait
  variance-covariance matrix. This applies to all random effects in the
  model.
\end{itemize}

\pandocbounded{\includegraphics[keepaspectratio]{https://raw.githubusercontent.com/JantN1966/MasterManual/main/Images/CompVar01.jpg}}\\

\emph{Example}. The lower triangular trait (co)variance components file
with two traits (body weight 1 and body weight 2) for non-genetic random
regression, animal genetic and residual effects.

\paragraph{6.1.3. Input file in sparse-matrix format for the MiX99
solver}\label{Comp05}

In the sparse matrix form, the order of the matrices must be the same as
the order of random effects in the model, with the restriction that the
genetic effect should be the last random effect in the model and the
elements of its (co)variance matrix should appear in the sparse matrix
file just before the elements of the residual (co)variance matrix. The
residual (co) variance matrix should be specified at the end of the
sparse matrix file. In summary, the order of matrices is: * Non-genetic
random effects in the same order as specified in the model * Genetic
effects as specified in the model * Residual effect The matrix elements
must be specified as the random effect number, row number, column number
and the value of the (co)variance. To avoid mistakes, it is recommended
to provide the elements of the lower triangle of the matrix, in other
words, any column number is smaller than or equal to the row number.
Off-diagonals only need to be specified if they are non-zero. When
haplotypes are used in the model for marker-assisted BLUP with the use
of an inverse IBD matrix, both haplotypes are counted as effects, but
the same variance components are used for the first and the second
haplotype, when haplotypes are combined with the AND function, so the
variance components should not be repeated for the second haplotype.
Effectively, the effect number corresponding to the second haplotype is
skipped from the list of inverse matrix elements. See Example 5.4 in the
Appendix.

\pandocbounded{\includegraphics[keepaspectratio]{https://raw.githubusercontent.com/JantN1966/MasterManual/main/Images/CompVar02.jpg}}\\

\emph{Example}. The trait (co) variance components file in sparse-matrix
format with two traits for the animal genetic and residual effects
Columns: random effect number, trait row number, trait column number and
variance or covariance component.

\paragraph{6.1.4. Syntax}\label{Comp06}

\begin{quote}
PARFILE \textless filename\textgreater{} {[}!SPARSE{]}
\end{quote}

Qualifier:

\textbf{!Sparse}\\
If !SPARSE is specified, the variance and covariance components are read
in sparse matrix form. If omitted, the matrix is read in lower
triangular form.

\subsubsection{6.2. Parameter files for general
covariates}\label{Comp07}

\paragraph{6.2.1. General}\label{Comp08}

The regression parameter file is specified for each general covariate
file that is fitted as random regression. The file may contain a single
set of variances and covariances between traits that apply to all
covariates or a set for each covariate separately. The MiXBLUP shell
checks whether scaling is necessary to avoid an error that the matrix is
not positive-definite and applies any required scaling automatically.
For the hpblup solver, the repression parameter file is provided as an
inverse.

\paragraph{6.2.2. Input file}\label{Comp09}

The format of the files with parameters of general covariates is the
lower-triangular-matrix format of the general parameter file. For the
MiX99 solver, every line of the variance covariance matrix starts with
the trait name, as it is used in MiXBLUP instruction file. Note that
trait names are case-sensitive. If !RegType R is specified for the
covariate file, a single trait-effect variance-covariance matrix can be
used for all covariates in the file. If !RegType H is used, a
trait-effect variance-covariance matrix has to be specified for each
covariate.

\pandocbounded{\includegraphics[keepaspectratio]{https://raw.githubusercontent.com/JantN1966/MasterManual/main/Images/CompVar03.jpg}}\\

\emph{Example}. Regression parameter file with a single set of variances
and covariances between traits for all covariates. A regression
parameter with covariate-specific variances and covariances contain such
a set for each covariate. The number in the label of the matrix is
linked with the position of the covariate in the record.

For the hpblup solver, a general covariate file may be fitted for
multiple indices, so it is necessary to specify the trait name followed
by the index name between brackets at the start of each line in the
variance covariance matrix. If !RegType H is used, a trait-effect
variance-covariance matrix has to be specified for each covariate.

\pandocbounded{\includegraphics[keepaspectratio]{https://raw.githubusercontent.com/JantN1966/MasterManual/main/Images/CompVar04.jpg}}\\

\emph{Example}. Regression parameter file with a single set of variances
and covariances between traits for all covariates, for the hpblup
solver.

\paragraph{6.2.3. Syntax}\label{Comp10}

\begin{quote}
REGPARFILE REG01 \textless file name REG01\textgreater{}\\
REG02 \textless file name REG02\textgreater{}\\
\textless\ldots\textgreater{}\\
REG99 \textless file name REG99\textgreater{}
\end{quote}

\textbf{REGPARFILE}\\
The REGPARFILE section must contain the name of a parameter file for
each covariate file for which the covariates are fitted as a random
regression (so !REGTYPE is either `r' or `h'). If the regression type is
fixed, the corresponding file in REGPARFILE is ignored. The REGPARFILE
section does not have any associated qualifiers. The lines of the
REGPARFILE section each contain two columns. The first column is the
label that links the parameter file to the covariate file. The second
column is the name of the file.

\subsubsection{6.3. Parameters for SNP covariate files}\label{Comp11}

\paragraph{6.3.1. General}\label{Comp12}

The SNP parameter file is specified for SNP covariate files that are to
be fitted for random regression. The file may contain a single set of
variances and covariances between traits for all SNP covariates or a set
for each SNP covariate separately. For a SNPBLUP model without a direct
genetic effect and SNP genotypes presented as 0, 1 and 2, the SNP
variance can be calculated from the direct genetic variance with\\

\pandocbounded{\includegraphics[keepaspectratio]{https://raw.githubusercontent.com/JantN1966/MasterManual/main/Images/CompVar05.jpg}}\\

where N is the number of informative SNPs and pi is the allele frequency
of the SNP allele counted on locus i. Non-informative SNPs must not be
included in this calculation. If variances smaller than 1.0E-06 are
specified, then the MiXBLUP kernel may give an error that the
variance-covariance matrix is not positive-definite. This can be
resolved by scaling the phenotypes with 10 or 100 and the variances with
100 or 10,000 accordingly. The MiXBLUP shell checks whether scaling is
necessary and applies any required scaling automatically.

\paragraph{6.3.2. Input file}\label{Comp13}

The format of the files with parameters of SNP covariates is the
lower-triangular-matrix format of the general parameter file. If a
single set of variances and covariances between traits is to be used for
all SNP covariates (so !REGTYPE is `r'), then only one matrix needs to
be specified. The matrix label needs to start with `SNP', but the number
is ignored. If SNP-specific variances and covariances are to be used (so
!REGTYPE is `h'), then a matrix has to be specified for every SNP
covariate separately. Depending on the number of SNP covariates in a
file, this could be many thousands. The label has to start with `SNP'.
The number in the label of the matrix is linked with the position of the
SNP covariate in the record of the corresponding file. The number must
be sequential and may be an integer between 1 and 2.1 billion. The label
of a matrix in a SNP parameter file refers to a SNP covariate in the
corresponding covariate file and should not be confused with the label
linking the SNP covariate and parameter files. For the MiX99 solver,
every line of the variance covariance matrix starts with the trait name,
as it is used in MiXBLUP instruction file.

\pandocbounded{\includegraphics[keepaspectratio]{https://raw.githubusercontent.com/JantN1966/MasterManual/main/Images/CompVar06.jpg}}\\

\emph{Example}. SNP parameter file with a single set of variances and
covariances between traits for all SNP covariates, for the MiX99 solver.

For the hpblup solver, a SNP covariate file may be fitted for multiple
indices, so it is necessary to specify the trait name followed by the
index name between brackets at the start of each line in the variance
covariance matrix.

\pandocbounded{\includegraphics[keepaspectratio]{https://raw.githubusercontent.com/JantN1966/MasterManual/main/Images/CompVar07.jpg}}\\

\emph{Example}. SNP parameter file with a single set of variances and
covariances between traits for all SNP covariates, for the hpblup
solver.

\paragraph{6.3.3. Syntax}\label{Comp14}

\begin{quote}
SNPPARFILE\\
SNP01 \textless file name SNP01\textgreater{}\\
SNP02 \textless file name SNP02\textgreater{}\\
\textless\ldots\textgreater{}\\
SNP99 \textless file name SNP99\textgreater{}
\end{quote}

\textbf{SNPPARFILE}\\
The SNPPARFILE section specifies the name of a parameter file for each
SNP covariate file for which the SNP covariates are fitted as a random
regression (so !REGTYPE is either `r' or `h'). If the regression type is
random and !CalcSNPvar has been specified, the corresponding file in
SNPPARFILE is ignored. The SNPPARFILE section does not have any
associated qualifiers. The lines of the SNPPARFILE section each contain
two columns. The first column is the label that links the parameter file
to the SNP covariate file. The second column is the name of the file.

\subsubsection{6.4. Parameters in case of heterogeneous residual
variances}\label{Comp15}

\paragraph{6.4.1. General}\label{Comp16}

The residual variance may not be the same for all observations. If this
is the case, observations can be grouped by their residual variance
prior to the analysis. A column in the data file links the observation
to the correct residual variance matrix. Modelling data with a random
regression approach often requires the use of multiple residual variance
classes.

\paragraph{6.4.2. Input file}\label{Comp17}

The file contains a matrix for every class number in the linking column
in the data file. The name of the matrix is Res followed by the class
number between brackets. The class number has to be an integer. The
example below gives the series of residual matrices for a situation with
observations being linked to one of three residual variances classes.

\pandocbounded{\includegraphics[keepaspectratio]{https://raw.githubusercontent.com/JantN1966/MasterManual/main/Images/CompVar08.jpg}}\\

\emph{Example}. The residual covariance file with three residual
variance-covariance matrices.

\paragraph{6.4.3. Syntax}\label{Comp18}

\begin{quote}
DATAFILE \textless filename\textgreater{}\\
\textless\ldots\textgreater{}\\
{[}\textless field j\textgreater{} I !RESVARCLASS{]}\\
\textless..\textgreater{}\\
RESFILE \textless filename\textgreater{}
\end{quote}

Section:

\textbf{RESFILE}\\
The RESFILE section specifies the name of the file with residual
between-trait variance-covariance matrices. The RESFILE section does not
have specific qualifiers.

Qualifier:

\textbf{!RESVARCLASS}\\
The qualifier !ResVarClass in the DATAFILE section links a data record
to the appropriate residual variance class, in case the residual
variance differs for groups of records. The field is integer. The
qualifier !RESVARCLASS must be used if the section RESFILE is specified.

\hyperref[Tabl01]{Back to Table of Contents}

\newpage

\subsection{7. Statistical models}\label{Stat01}

\textbf{Observed traits on two individuals may be similar due to genetic
effects and systematic non-genetic effects. The statistical model
contains all such effects known to explain variation in observed traits.
This chapter describes various statistical models to estimate genetic
effects on traits with as little bias as possible.}

\subsubsection{7.1. Basic models}\label{Stat02}

\paragraph{7.1.1. General}\label{Stat03}

The MODEL section specifies the start of the statistical models for the
traits in the analysis. Traits and statistical models start immediately
below the line with the MODEL keyword. For each trait, the statistical
model is specified on a separate line. MiXBLUP supports up to 63 traits
to be analysed simultaneously, if computer resources permit this. The
basic statistical model for a breeding value evaluation contains fixed
effects, uncorrelated, non-genetic random effects and a direct genetic
random effect. Uncorrelated, non-genetic random class effects are
assumed to have an identity relationship matrix between levels of the
effect. Each model line contains trait name, fixed effects, non-genetic
random effects and genetic random effects. Fixed effects may be class
effects, covariates or covariates nested within a class effect.
Similarly, random effects may be class effects or covariates nested
within a class effect. The residual random effect does not need to be
specified. The minimum statistical model contains one fixed effect and
one genetic random effect.

\paragraph{7.1.2. Syntax}\label{Stat04}

\begin{quote}
MODEL\\
\textless trait1\textgreater{} \textasciitilde{} \textless class
effects\textgreater{} {[}covariates{]} {[}class effect * covariates{]}
\&!RANDOM G(\textless direct genetic effect\textgreater)
{[}\textless non-genetic random effects\textgreater{]}\\
{[}\textless trait2\textgreater{} \ldots{]}\\
\textless\ldots\textgreater{}\\
{[} \ldots{]}
\end{quote}

Section:

\textbf{MODEL}\\
The MODEL section specifies the statistical model for the traits in the
analysis.

Qualifiers:

\textbf{\textasciitilde{} (tilde)}\\
The tilde separates the trait from the statistical model.

\textbf{* (star)}\\
The star is used for nesting a covariate within a class effect, to yield
a regression coefficient for each level of the class effect. This can be
used for both fixed and random nested covariates. The star must not be
used to model an interaction of class effects.

\textbf{!RANDOM}\\
The !RANDOM qualifier separates the fixed effects from the random
effects in the model.

\textbf{G(\ldots)}\\
The G(\ldots) function links a random effect to the inverse genetic or
genomic relationship matrix.

\paragraph{7.1.3. Associated output files}\label{Stat05}

{\def\LTcaptype{none} % do not increment counter
\begin{longtable}[]{@{}
  >{\raggedright\arraybackslash}p{(\linewidth - 2\tabcolsep) * \real{0.5000}}
  >{\raggedright\arraybackslash}p{(\linewidth - 2\tabcolsep) * \real{0.5000}}@{}}
\toprule\noalign{}
\begin{minipage}[b]{\linewidth}\raggedright
Output file
\end{minipage} & \begin{minipage}[b]{\linewidth}\raggedright
Description
\end{minipage} \\
\midrule\noalign{}
\endhead
\bottomrule\noalign{}
\endlastfoot
Solfix.txt & Solutions of fixed effects by trait (class effects,
covariates and nested covariates) \\
Solfix.out & As Solfix.txt, but for alphanumerical class labels \\
Solr00.txt & Solutions of non-genetic, uncorrelated random effect 00 by
trait, where 00 ranges \\
& from 1 to the number of non-genetic, uncorrelated random effects
across traits \\
& (class effects, covariates and nested covariates) \\
Solr00.out & As Solr00.txt, but for alphanumerical class labels \\
\end{longtable}
}

\subsubsection{7.2. Repeatability models}\label{Stat06}

\paragraph{7.2.1. General}\label{Stat07}

Certain traits are measured just once in the lifetime of an individual.
Other traits may be measured repeatedly. Two observations on a single
individual are more similar than expected from having the same genotype.
A permanent environmental effect is usually added to the model to
account for non-genetic similarity of records of the same individual.
Such a permanent environmental effect has the same label as the direct
genetic effect, but with an identity relationship matrix between levels.
This permanent environmental effect should have its own column in the
data file and must not be the column with the direct genetic effect
(although identical).

\paragraph{7.2.2. Syntax}\label{Stat08}

\begin{quote}
MODEL\\
\textless trait\textbackslash1\textgreater{} \textasciitilde{}
\textless fixed effects\textgreater{} !RANDOM G(\textless direct genetic
effect\textgreater) \textless permanent environmental
effect\textgreater{}\\
\textless\ldots\textgreater{}\\
{[}\textless traitN\textgreater{} \ldots{]}
\end{quote}

\paragraph{7.2.3. Associated output files}\label{Stat09}

{\def\LTcaptype{none} % do not increment counter
\begin{longtable}[]{@{}
  >{\raggedright\arraybackslash}p{(\linewidth - 2\tabcolsep) * \real{0.5000}}
  >{\raggedright\arraybackslash}p{(\linewidth - 2\tabcolsep) * \real{0.5000}}@{}}
\toprule\noalign{}
\begin{minipage}[b]{\linewidth}\raggedright
Output file
\end{minipage} & \begin{minipage}[b]{\linewidth}\raggedright
Description
\end{minipage} \\
\midrule\noalign{}
\endhead
\bottomrule\noalign{}
\endlastfoot
Solr00.txt & Solutions of non-genetic, uncorrelated random effect 00 by
trait, where 00 ranges \\
& from 1 to the number of non-genetic, uncorrelated random effects
across traits \\
& (among which the permanent environmental effect) \\
Solr00.out & As Solr00.txt, but for alphanumerical class labels \\
\end{longtable}
}

\subsubsection{7.3. Maternal genetic models}\label{Stat10}

\paragraph{7.3.1. General}\label{Stat11}

Some traits are affected by the genotype of the animal itself and the
genotype of its dam at the same time. An example is weaning weight in
beef cattle. For such traits, a maternal genetic model should be used.
The inverse numerator relationship matrix is applied both to the direct
genetic effect (animal) and the maternal genetic effect (dam). The
maternal genetic effect must be a separate field in the data file and
each individual in this field must exist as an individual in pedigree,
genotype file or other resource of genetic similarity. For biological
dams, this is self-evident, but for foster dams, this requires special
attention.

\paragraph{7.3.2. Syntax}\label{Stat12}

\begin{quote}
MODEL\\
\textless trait\textbackslash1\textgreater{} \textasciitilde{}
\textless fixed effects\textgreater{} !RANDOM
\textless\ldots\textgreater{} G\textless direct genetic
effect\textgreater, \textless maternal genetic effect\textgreater)\\
\textless\ldots\textgreater{}\\
{[}\textless traitN\textgreater{} \ldots{]}
\end{quote}

The maternal genetic effect is placed within the brackets of function G
to link it with the relationship matrix between individuals.

\paragraph{7.3.3. Associated output files}\label{Stat13}

{\def\LTcaptype{none} % do not increment counter
\begin{longtable}[]{@{}
  >{\raggedright\arraybackslash}p{(\linewidth - 2\tabcolsep) * \real{0.5000}}
  >{\raggedright\arraybackslash}p{(\linewidth - 2\tabcolsep) * \real{0.5000}}@{}}
\toprule\noalign{}
\begin{minipage}[b]{\linewidth}\raggedright
Output file
\end{minipage} & \begin{minipage}[b]{\linewidth}\raggedright
Description
\end{minipage} \\
\midrule\noalign{}
\endhead
\bottomrule\noalign{}
\endlastfoot
Solani.txt & The solutions of the maternal genetic effect are included
as additional columns in \\
& Solani.txt. The exact layout of Solani.txt is printed at the end of
solver.log. \\
Solani.out & As Solani.txt, but for alphanumerical ID labels \\
\end{longtable}
}

\subsubsection{7.4. Social interaction models}\label{Stat14}

\paragraph{7.4.1. General}\label{Stat15}

The social interaction model (or group selection) is used to estimate
the effects of an individual's genotype on its own performance and on
the performance of its pen mates simultaneously. It should be used for
groups of a single group size, but a slightly varying group size is also
supported. For a single group size, a genetic effect for social
interaction is fitted for each pen mate. This effect can be interpreted
as the genetic value for supporting or inhibiting the expression of the
direct genetic merit of pen mates. The genetic variance of this social
effect is dependent on the size of the group, so performance in small
and large groups by design should not be combined as one trait. If the
size of groups is the same by design, but it varies slightly due to
death or removal from the pen, it is still possible to fit a social
interaction model by adding a covariate of either 0 (not present) or 1
(present) and apply a nested regression of this covariate within pen
mate. The ID label used for not-present pen mates must appear in the
pedigree or genotype file, too, but no information is added to its
social genetic value when the covariate is zero (not present).

\paragraph{7.4.2. Syntax of the social interaction model with one group
size for all groups for the MiX99 solver}\label{Stat16}

\begin{quote}
DATAFILE\\
\textless ID individual\textgreater{} \textless field type I or
A\textgreater{}\\
\textless\ldots\textgreater{}\\
mate1 I\\
mate2 I\\
\textless\ldots\textgreater{}\\
mateN I\\
MODEL\\
\textless trait1\textgreater{} \textasciitilde{}
\textless\ldots\textgreater{} !RANDOM \textless\ldots\textgreater{}
G(\textless ID individual\textgreater,mate1 AND mate2
\textless\ldots\textgreater{} AND mateN)\\
{[}\textless trait2\textgreater{} \ldots{]}\\
\textless\ldots\textgreater{}\\
{[}\textless traitN\textgreater{} \ldots{]}
\end{quote}

The pen mates need to be defined in the data file. The number of
additional columns is equal to the number of pen mates (mate1, mate2,
\ldots, mateN).

Qualifier:

\textbf{AND}\\
The function AND combines the effects of different mates to one design
matrix. There are a few constraints when combining effects; all of them
are rare:

\begin{itemize}
\tightlist
\item
  When `AND' is used for group selection, it cannot be used for
  IBD-haplotypes. In other words, haplotype effects cannot be included
  in a social interaction model.
\item
  Combining of effects needs to be the same for any traits for which
  effects are combined. In other words, traits measured when the animal
  was in a different group need to be analyzed in a separate analysis.
\item
  The parser supports an evaluation in which some traits have social
  genetic effects and other traits do not.
\item
  There can be only one group of combined genetic effects, so no commas
  must be placed between the effects of mates. It means that the genetic
  effects can only be combined to one other effect and not more.
\item
  The order of genetic effects should be kept the same across traits
  with a social interaction model.
\end{itemize}

\paragraph{7.4.3. Syntax of the social interaction model with slightly
varying group sizes for the MiX99 solver}\label{Stat17}

\begin{quote}
DATAFILE\\
\textless ID individual\textgreater{} \textless field type I or
A\textgreater{}\\
\textless\ldots\textgreater{} mate1 I\\
present1 R\\
mate2 I\\
present2 R\\
\textless\ldots\textgreater{}\\
mateN I\\
presentN R\\
MODEL\\
\textless trait1\textgreater{} \textasciitilde{}
\textless\ldots\textgreater{} !RANDOM \textless\ldots\textgreater{}
G(\textless ID individual\textgreater,present1*mate1 AND present2*mate2
\textless\ldots\textgreater{} AND presentN*mateN)\\
{[}\textless trait2\textgreater{} \textless\ldots\textgreater{]}\\
\textless\ldots\textgreater{}\\
{[}\textless traitN\textgreater{} \textless\ldots\textgreater{]}
\end{quote}

Both the pen mates and their presence need to be defined in the data
file. The number of additional columns is therefore equal to two times
the number of pen mates (mate1, mate2, \ldots, mateN, present1,
present2, \ldots, presentN).

\paragraph{7.4.4. Syntax of the social interaction model for the hpblup
solver}\label{Stat18}

\begin{quote}
DATAFILE\\
Animal I\\
\textless\ldots\textgreater{}\\
mate1 I\\
mate2 I\\
\textless\ldots\textgreater{}\\
mateN I\\
MODEL\\
\textless trait1\textgreater{} \textasciitilde{}
\textless\ldots\textgreater{} \textless\ldots\textgreater{} !RANDOM
\textless\ldots\textgreater{} G(Animal,LINK(mate1))\\
{[}\textless traitN\textgreater{} \textless\ldots\textgreater{]}\\
LINKEDEFFECTS\\
mate1 \textasciitilde{} mate2 \ldots{} mateN
\end{quote}

The pen mates need to be defined in the data file. The number of
additional columns is equal to the number of pen mates (mate1, mate2,
\ldots, mateN). For slightly varying group sizes, just use a zero for a
missing pen mate.

Section:

\textbf{LINKEDEFFECTS}\\
The section LinkedEffects is used to specify which effects are linked to
the leading effect and should be combined with the leading effect. The
leading effect should be presented before the tilde (\textasciitilde);
linked effects should be presented after it. The number of lines in the
LinkedEffect section matches the number of unique occurrences of the
LINK function in the MODEL section. So if multiple traits contain
G(Animal,LINK(mate1)) in the model, only one line is needed in the
LinkedEffect section.

Qualifier:

\textbf{LINK}\\
The function LINK specifies the leading effect of a set of linked
effects in the model.

\paragraph{7.4.5. Associated output files}\label{Stat19}

{\def\LTcaptype{none} % do not increment counter
\begin{longtable}[]{@{}
  >{\raggedright\arraybackslash}p{(\linewidth - 2\tabcolsep) * \real{0.5000}}
  >{\raggedright\arraybackslash}p{(\linewidth - 2\tabcolsep) * \real{0.5000}}@{}}
\toprule\noalign{}
\begin{minipage}[b]{\linewidth}\raggedright
Output file
\end{minipage} & \begin{minipage}[b]{\linewidth}\raggedright
Description
\end{minipage} \\
\midrule\noalign{}
\endhead
\bottomrule\noalign{}
\endlastfoot
Solani.txt & The solutions of the social genetic effects are included as
additional columns in \\
& Solani.txt. The exact layout of Solani.txt is printed at the end of
solver.log. \\
Solani.out & As Solani.txt, but for alphanumerical ID labels \\
\end{longtable}
}

\subsubsection{7.5. Random regression models}\label{Stat20}

\paragraph{7.5.1. General}\label{Stat21}

There are two types of random regression models supported by MiXBLUP,
the non-genetic and genetic random regression model. Both original
covariates and polynomials derived from an independent variable may be
used in the model. In a non-genetic random regression model, the
regression of the observations on an independent covariate is fitted as
a random effect. Random regression in MiXBLUP has to be specified as
regression nested within a class variable. If no nested regression is
required, the user needs to add a column of ones to the data file, and
fit the covariate within this class effect of a single level. The
internal structure of MiXBLUP requires that any random effect be
associated with a class effect. This can be seen in the parameter files,
too, where variance-covariance matrices are all specified by class
effect. In a genetic random regression model, trait observations are
regressed on the covariate within animals, taking into account the
genetic relationships between animals. The estimated breeding values
from such an analysis concern the animal-specific parameters of the line
or curve fitted. The user needs to convert these estimates to estimated
breeding values at a given level of the covariate or for a function of
levels of the covariate. If the relationship between an observed trait
and an independent variable is non-linear, it may still be possible to
model the relationship with polynomials, as a special case of multiple
linear regression. Polynomial regression is a form of linear regression
in which the relationship between the independent variable x and the
observed trait is modelled as an nth degree polynomial in x by fitting
(n+1) covariates derived from x. Polynomials may be provided by the user
either in the data file or in a covariate table, or may calculated
during the preparations for the analysis and stored in a covariate
table. Polynomials calculated by MiXBLUP are Legendre polynomials.

\paragraph{7.5.2. Syntax of a simple non-genetic random regression model
for the MiX99 solver}\label{Stat22}

\begin{quote}
MODEL\\
\textless trait1\textgreater{} \textasciitilde{}
\textless\ldots\textgreater{} !RANDOM
\textless class\textgreater*\textless covariate\textgreater{}
{[}\textless class\textgreater*\textless covariate\textgreater{}
\textless\ldots\textgreater{]}\\
\textless\ldots\textgreater{}\\
{[}\textless traitN\textgreater{} \textless\ldots\textgreater{]}
\end{quote}

The random regression term consists of a class effect with field type
integer (I) or alphanumerical (A) and a covariate with field type real
(R). Each random regression term has to be present in the
variance-covariance matrix of the class effect in the parameter file
(see chapter \hyperref[Comp02]{6.1.}).

Qualifier:

\textbf{* (star)}\\
The star is used for nesting a covariate within a class effect, to yield
a regression coefficient for each level of the class effect. There is no
specific order of class effect and covariate in the term.

\paragraph{7.5.3. Syntax of a simple genetic random regression model for
the MiX99 solver}\label{Stat23}

\begin{quote}
MODEL\\
\textless trait1\textgreater{} \textasciitilde{}
\textless\ldots\textgreater{} !RANDOM
G(\textless ID\textgreater*\textless covariate\textgreater{[},\textless ID\textgreater*\textless covariate\textgreater{]})
\textless\ldots\textgreater{}\\
\textless\ldots\textgreater{}\\
{[}\textless traitN\textgreater{} \textless\ldots\textgreater{]}
\end{quote}

The regression terms nested within the individual's ID are placed within
the function G(\ldots) to indicate that the relationship matrix of
individuals should be used.

\paragraph{7.5.4. Syntax of a polynomial regression model using a
covariate table for the MiX99 solver}\label{Stat24}

\begin{quote}
CVRTABLE \textless name covariate\textgreater{}\\
MODEL\\
\textless trait1\textgreater{} \textasciitilde{}
\textless\ldots\textgreater{}
\textless class\textgreater*CVR(\textless n1\textgreater) !RANDOM
\textless class\textgreater*CVR(\textless n2\textgreater)
G(\textless ID\textgreater*CVR(\textless n3\textgreater))
\textless\ldots\textgreater{}\\
\textless\ldots\textgreater{}\\
{[}\textless traitN\textgreater{} \textless\ldots\textgreater{]}
\end{quote}

For the use of covariate table files with the MiX99 solver, see chapter
\hyperref[Obse10]{4.2.2.1.} and \hyperref[Obse12]{4.2.2.3.}.

Qualifier:

\textbf{CVR(\ldots)}\\
The CVR function is used in the MODEL section and is a shorthand for all
polynomial terms to be fitted and may be used in the same way as any
individual random regression term. The alternative way to specify
polynomial random regression is to use the individual columns of the
covariate table file. The names of the columns are cvr00, cvr01, cvr02,
\ldots, cvrnn. The label is lowercase and has exactly two digits ranging
from 00 to 99.

\paragraph{7.5.5. Syntax of a simple non-genetic random regression model
for the hpblup solver}\label{Stat25}

\begin{quote}
REGFILE\\
\textless field ID\textgreater{} \textless field type I or
A\textgreater{}\\
REG02 \textless file name REG02\textgreater{} !REGTYPE \textless R or
H\textgreater{}\\
REGPARFILE\\
REG02 \textless file name REG02\textgreater{}\\
MODEL\\
\textless trait1\textgreater{} \textasciitilde{}
\textless\ldots\textgreater{} !RANDOM
\textless class\textgreater*hpReg(\textless field ID\textgreater,2)
{[}\textless class\textgreater*hpReg(\textless field dam
ID\textgreater,2) \textless\ldots\textgreater{]}\\
\textless\ldots\textgreater{}\\
{[}\textless traitN\textgreater{} \textless\ldots\textgreater{]}
\end{quote}

For the hpblup solver, covariates in a random regression should be
provided in a covariate table or a general covariate file, but not as a
field in the data file. The random regression term consists of a class
effect with field type integer (I) or alphanumerical (A) and a covariate
with field type real (R). Each random regression term has to be present
in the variance-covariance matrix of the class effect in the parameter
file (see chapter 6.1).

Qualifier:

\textbf{* (star)}\\
The star is used for nesting a covariate within a class effect, to yield
a regression coefficient for each level of the class effect. There is no
specific order of class effect and covariate in the term.

\paragraph{7.5.6. Syntax of a simple genetic random regression model for
the hpblup solver}\label{Stat26}

\begin{quote}
REGFILE\\
\textless field ID\textgreater{} \textless field type I or
A\textgreater{}\\
REG02 \textless file name REG02\textgreater{} !REGTYPE \textless R or
H\textgreater{}\\
REGPARFILE\\
REG02 \textless file name REG02\textgreater{}\\
MODEL\\
\textless trait1\textgreater{} \textasciitilde{}
\textless\ldots\textgreater{} !RANDOM
G(\textless ID\textgreater*hpReg(\textless field ID\textgreater,2))
\textless\ldots\textgreater{}\\
\textless\ldots\textgreater{}\\
{[}\textless traitN\textgreater{} \textless\ldots\textgreater{]}
\end{quote}

For the hpblup solver, covariates in a random regression should be
provided in a covariate table or a general covariate file, but not as a
field in the data file. The regression terms nested within the
individual's ID are placed within the function G(\ldots) to indicate
that the relationship matrix of individuals should be used.

\paragraph{7.5.7. Syntax of a polynomial regression model using a
covariate table for the hpblup solver}\label{Stat27}

\begin{quote}
CVRTABLE !nCVRTABLES 2\\
TABLE01 \textless filename\textgreater{} !CVRNUM \textless nth
order\textgreater{} !CVRMIN \textless minimum value\textgreater{}
!CVRMAX \textless maximum value\textgreater{} !CVRSingleCov !CVRIndex
\textless index field name\textgreater{}\\
TABLE04 \textless filename\textgreater{} !CVRNUM \textless nth
order\textgreater{} !CVRMIN \textless minimum value\textgreater{}
!CVRMAX \textless maximum value\textgreater{} !CVRSingleCov !CVRIndex
\textless index field name\textgreater{}\\
MODEL\\
\textless trait\textgreater{} \textasciitilde{} \textless fixed
effects\textgreater{} \textless Class1\textgreater*TABLE01 !RANDOM
\textless Class2\textgreater*TABLE04 G(Animal*TABLE04)
\end{quote}

For the use of covariate table files with the hpblup solver, see chapter
\hyperref[Obse11]{4.2.2.2.} and \hyperref[Obse13]{4.2.2.4.}.

Qualifier:

\textbf{TABLEnn}\\
The TABLEnn label is a shorthand for a specific covariate table file. It
automatically fits all covariates in the file, unlike for the
CVR(\ldots) function for the MiX99 solver, which can be used to fit a
smaller number of covariates from a covariate table file. The names of
the covariates in the parameter file with trait variance-covariance
matrices are TABLEnn\_cc, where nn is the table number and cc the
covariate number starting with 00 (for example TABLE01\_00 for the first
covariate).

\paragraph{7.5.8. Associated output files}\label{Stat28}

{\def\LTcaptype{none} % do not increment counter
\begin{longtable}[]{@{}
  >{\raggedright\arraybackslash}p{(\linewidth - 2\tabcolsep) * \real{0.5000}}
  >{\raggedright\arraybackslash}p{(\linewidth - 2\tabcolsep) * \real{0.5000}}@{}}
\toprule\noalign{}
\begin{minipage}[b]{\linewidth}\raggedright
Output file
\end{minipage} & \begin{minipage}[b]{\linewidth}\raggedright
Description
\end{minipage} \\
\midrule\noalign{}
\endhead
\bottomrule\noalign{}
\endlastfoot
Solani.txt & The solutions of the genetic nested regression effects are
included as additional \\
& columns in Solani.txt. The exact layout of Solani.txt is printed at
the end of solver.log. \\
Solani.out & As Solani.txt, but for alphanumerical ID labels \\
Solr00.txt & The solutions of the non-genetic nested regression effects
are included as \\
& additional columns in Solr00.txt. The exact layout of Solr00.txt is
printed at the end \\
& of solver.log. \\
Solr00.out & As Solr00.txt, but for alphanumerical class labels \\
\end{longtable}
}

\subsubsection{7.6. Weighting residuals by record}\label{Stat29}

\paragraph{7.6.1. General}\label{Stat30}

If the common assumption of constant standard deviation of the residuals
(i.e.~homogeneous residual variance) is not met, it is possible to
weight individual records. Less precise measurements get less weight and
more weight is given to more precise measurements when estimating
breeding values. An example is the use of de-regressed breeding values
as a pseudo-phenotype. The standard deviation of the residual depends on
the reliability of the original estimated breeding value. A weighting
factor based on the reliability can be used to give more weight to
pseudo-phenotypes based on a relatively large amount of information.
Another example is variation in residual variances within contemporary
groups. Observations in contemporary groups with a large residual
variance can be given a proportionally lower weighting factor.

\paragraph{7.6.2. Syntax}\label{Stat31}

\begin{quote}
MODEL\\
\textless trait1\textgreater{} {[}!WEIGHT \textless weighting
factor\textgreater{]} \textasciitilde{} \textless fixed
effects\textgreater{} !RANDOM G(\textless ID\textgreater)
{[}\textless non-genetic random effects\textgreater{]}\\
\textless\ldots\textgreater{}\\
{[}\textless traitN\textgreater{} \textless\ldots\textgreater{]}
\end{quote}

Qualifier:

\textbf{!WEIGHT}\\
A field in the data file can be specified as a weighting factor for a
specific trait using the !WEIGHT qualifier.

\paragraph{7.6.3. Associated output files}\label{Stat32}

The standard output files are used for a weighted analysis.

\subsubsection{7.7. Combining effects across traits}\label{Stat33}

\paragraph{7.7.1. General}\label{Stat34}

If a trait measured in different cycles or parities or on individuals of
different strains and crosses is modelled as multiple traits, it may be
desirable to estimate fixed effects across these traits, in order to
increase the precision of the solutions of the model. Random effects can
easily be combined by specifying covariances between the traits that are
equivalent to a correlation close to unity. For fixed effects, it has to
be specified across which traits the effect should be estimated.

\paragraph{7.7.2. Syntax}\label{Stat35}

\begin{quote}
MODEL\\
\textless trait1\textgreater{} \textasciitilde{}
\textless fixed1\textgreater{} !RANDOM G(\textless ID\textgreater)
{[}\textless random1\textgreater{]}\\
\textless trait2\textgreater{} \textasciitilde{}
\textless fixed1\textgreater{} !RANDOM G(\textless ID\textgreater)
{[}\textless random1\textgreater{]}\\
\textless\ldots\textgreater{}\\
\textless traitN\textgreater{} \textasciitilde{}
\textless fixed\textgreater{} !RANDOM G(\textless ID\textgreater)
{[}\textless random\textgreater{]}\\
COMBINE\\
\textless fixed1\textgreater{} \textasciitilde{}
\textless trait1\textgreater{} \textless trait2\textgreater{}
\end{quote}

Section:

\textbf{COMBINE}\\
The section COMBINE allows to specify across which traits a fixed effect
should be estimated. It supports class effects, covariates and nested
covariates.

\paragraph{7.7.3. Associated output files}\label{Stat36}

The standard output files are used for an analysis with fixed and random
effects estimated across several traits.

\subsubsection{7.8. Correction of heterogeneous residual
variances}\label{Stat37}

\paragraph{7.8.1. General}\label{Stat38}

If residual variance within contemporary groups varies (heterogeneous
residual variance), the user may specify appropriate weighting factors
in the data file and weight records accordingly (see chapter
\hyperref[Stat29]{7.6.}). MiXBLUP also offers the possibility to
calculate appropriate weighting factors in a three-step approach. In the
first step, the traits are analysed with the assumption of homogeneous
residual variance. The residuals (ê) are read from the output of step 1
and the linearized squared residuals (z) for trait i and animal j are
calculated as

\pandocbounded{\includegraphics[keepaspectratio,alt={equation01}]{https://raw.githubusercontent.com/JantN1966/MasterManual/main/Images/StatMod01.jpg}}\\

Var(e\textsubscript{i}) is the residual variance of trait i used in the
first step and is obtained from the res(idual) matrix in the parameter
file or, if residual variance classes are used, the residual variance of
the corresponding class of the record. In the second step, these
linearised squared residuals are analysed using a suitable model. The
predicted phenotypes of this second model are used to calculate
weighting factors. The weighting factor for trait i and individual j is
calculated from the predicted value of the linearised squared residual
(Z\textsubscript{ij}) as

\pandocbounded{\includegraphics[keepaspectratio,alt={equation01}]{https://raw.githubusercontent.com/JantN1966/MasterManual/main/Images/StatMod02.jpg}}\\

where Z\textsubscript{i.} is the average predicted value of the
linearized squared residual for trait i across all individuals. In the
third step, the analysis of the first step is repeated, but with a
weighting factor added to account for heterogeneous residual variance.
MiXBLUP can run these three steps in a single process.

\paragraph{7.8.2. Syntax}\label{Stat39}

\begin{quote}
MODEL\\
\textless trait1\textgreater{} \textasciitilde{}
\textless fixed1\textgreater{} !RANDOM \textless random1\textgreater{}
G(\textless ID\textgreater)\\
\textless trait2\textgreater{} \textasciitilde{}
\textless fixed2\textgreater{} !RANDOM \textless random2\textgreater{}
G(\textless ID\textgreater)\\
\textless trait3\textgreater{} \textasciitilde{}
\textless fixed3\textgreater{} !RANDOM \textless random3\textgreater{}
G(\textless ID\textgreater)\\
VARMODEL\\
LSR1 \textasciitilde{} \textless fixed\textgreater{} !RANDOM
\textless random\textgreater{} G() !VARTRAIT
\textless trait1\textgreater{} LSR2 \textasciitilde{}
\textless fixed\textgreater{} !RANDOM \textless random\textgreater{}
G(\textless ID\textgreater) !VARTRAIT \textless trait2\textgreater{}\\
SOLVING\\
!DHGLM !HETCOR
\end{quote}

Section:

\textbf{VARMODEL}\\
The VARMODEL section specifies the statistical model for the second step
for the linearized squared residuals.

Qualifiers:

\textbf{!VARTRAIT}\\
The qualifier !VARTRAIT in the VARMODEL section is mandatory and links
the linearized squared residual to the original trait. Original traits
do not have to be all represented in the VARMODEL section.

\textbf{!DHGLM}\\
The option !DHGLM in the SOLVING section prepares MiXBLUP for multiple
calls of the kernel.

\textbf{!HETCOR}\\
The qualifier !HETCOR in the SOLVING section creates the data file and
instruction file for each step.

\paragraph{7.8.3. Associated output files}\label{Stat40}

The standard output files are used for an analysis with correction for
heterogeneous variances.

\subsubsection{7.9. Using a threshold model for a categorical trait
(MiX99 solver only)}\label{Stat41}

\paragraph{7.9.1. General}\label{Stat42}

The linear model used in MiXBLUP to estimate breeding values is based on
the assumptions that a trait has a continuous normal distribution, its
components of variance are homogeneous and residuals are uncorrelated
with genetic and non-genetic random effects. There are traits that are
recorded as categories. A binary trait has only two possible categories,
for example, present or absent, true or false, all or none. Traits with
more than two categories may be ordered, for example small-medium-large,
or unordered, such as red-yellow-blue. For categorical traits, the usual
assumptions of a linear model are violated.

It has been proposed to overcome this by assuming a continuous trait
underlying the categorical trait. Thresholds on the underlying scale
determine the recorded category on the observed scale. The threshold
model is more demanding than the linear model. There are methods
available for the threshold model: Newton-Raphson (NR) and
Expectation-Maximisation (EM). Both methods give the same results but
use a different route. Both methods have two levels of iterations. The
outer level iterates on the thresholds, given the set of estimated
solutions of the previous iteration. The inner level iterates on the
solutions given the current estimates of the threshold. Currently only
one categorical trait can be analysed with a threshold model in a
multi-trait evaluation of any number of traits analysed with a linear
model. Although theoretically incorrect, assuming a linear model for a
categorical trait often yield solutions that rank selection candidates
largely in correct order. This is especially the case for intermediate
prevalence of categories.

\paragraph{7.9.2. Input files}\label{Stat43}

Category labels must be numbered 1 to the number of categories for the
MiXBLUP kernel. MiXBLUP can rename category labels for this purpose from
a file with ordered labels by trait. The file is specified with
!CONVERTCAT. The first field is the trait name in the data file.
Subsequent fields contain the category labels. The position in the
sequence determines the new sequential integer code, 1..n. Although
MiXBLUP currently supports only a single trait with a threshold model in
combination with any number of traits with a linear model, multiple
traits may be specified for use across evaluations. In the example
below, the stature categories Small, Medium and Large are converted to
1, 2 and 3, according to their positions in the record. The diseased
categories 1 and 0 are converted to 1 and 2. Note that a binary trait
coded as 0/1 has to be converted to 1 and 2.

\begin{quote}
Stature Small Medium Tall\\
Diseased 1 0
\end{quote}

\emph{Example}. Category conversion table.

It is also possible to fix thresholds at a predefined value, using
!THRFIXED. Thresholds must be taken from a N(0,1) underlying
distribution. The file may contain any number of categorical traits, but
in any evaluation, only one categorical trait can be analysed with a
threshold model, currently. The number of thresholds to be specified is
the number of categories minus 1. Thresholds may be taken from a
previous analysis or calculated from the observed prevalence in a larger
set of data.

\begin{quote}
Stature -0.34 0.56\\
Diseased 0.0
\end{quote}

\emph{Example}. Table with fixed thresholds.

\paragraph{7.9.3. Syntax}\label{Stat44}

\begin{quote}
DATAFILE \textless filename\textgreater{} {[}!CONVERTCAT
\textless filename\textgreater{]}\\
\textless\ldots\textgreater{}\\
MODEL\\
\textless trait1\textgreater{} \textasciitilde{}
\textless fixed1\textgreater{} !RANDOM \textless random1\textgreater{}
G(\textless ID\textgreater)\\
\textless\ldots\textgreater{}\\
\textless trait3\textgreater{} \textasciitilde{}
\textless fixed3\textgreater{} !RANDOM \textless random3\textgreater{}
G(\textless ID\textgreater) !THRESHOLD \textless number of
thresholds\textgreater{}\\
SOLVING\\
{[}!THRMAXIT \textless maximum number of NR or EM iterations at outer
level\textgreater{]} {[}!THRMAXPCG \textless maximum number of
iterations at inner level\textgreater{]}\\
{[}!THRMETHOD \textless EM or NR\textgreater{]}\\
{[}!THRFIXED \textless filename\textgreater{]}
\end{quote}

Qualifiers:

\textbf{!CONVERTCAT}\\
The qualifier !CONVERTCAT in the DATAFILE section is optional. It can be
used to specify a file with category labels.

\textbf{!THRESHOLD}\\
The option !THRESHOLD in the MODEL section is mandatory and specifies
which trait is to be analysed with a threshold model.

\textbf{!THRMAXIT}\\
The qualifier !THRMAXIT in the SOLVING section can be used to specify
the maximum number of NR or EM iterations. The default number is 5,000.

\textbf{!THRMXPCG}\\
The qualifier !THRMAXPCG in the solving section is optional and
specifies the maximum number of iterations within each NR or EM
iteration. The default number is 100.

\textbf{!THRMETHOD}\\
The qualifier !THRMETHOD is optional and specifies the method to
implement the threshold model. The default method is NR. The alternative
method is EM.

\textbf{!THRFIXED}\\
The qualifier !THRFIXED is optional and can be used to specify a file
with fixed thresholds per trait.

\paragraph{7.9.4. Associated files}\label{Stat45}

The standard output files are used for an analysis with a threshold
model.

\hyperref[Tabl01]{Back to Table of Contents}

\newpage

\subsection{8. Control of analysis and output}\label{Cont01}

\textbf{This chapter describes the control part of the instruction file,
which can be used to control the analysis and the output generated from
the analysis.}

\subsubsection{8.1. Control of the analysis}\label{Cont02}

\paragraph{8.1.1. General}\label{Cont03}

Control of an iterative process like solving a linear system to estimate
breeding values involves setting the convergence criterion and the
maximum number of iterations, specifying whether the run is a
continuation or a new start and defining the starting values for a new
evaluation. An more advanced option is to specify the type of
preconditioner to be used for solving the system. Generally, the default
type of preconditioner is optimal. The default varies across models to
specify genetic similarity between individuals.

\paragraph{8.1.2. Syntax}\label{Cont04}

\subparagraph{8.1.2.1. Syntax when using MiX99 solver}\label{Cont05}

\begin{quote}
SOLVING\\
{[}!MAXIT \textless number of rounds\textgreater{]}\\
{[}!STOPCRIT \textless convergence criterion\textgreater{]}\\
{[}!NOPEEK{]}\\
{[}!PEEKFIRST \textless iteration number\textgreater{]}\\
{[}!PEEKEVERY \textless number of rounds\textgreater{]}\\
{[}!PEEKKEEP{]}\\
{[}!RESTART{]}\\
{[}!GFROMDISK{]}\\
PRECON \textless n, b, d\textgreater{}\\
{[}!WITHINBL \textless b, d\textgreater{]}\\
{[}!ACROSSBL \textless f, m, b, d\textgreater{]}
\end{quote}

Sections:

\textbf{SOLVING} The SOLVING section is used to control the process and
the output of the analysis.

\textbf{PRECON \textless option\textgreater{}} The PRECON section can be
used to change the type of preconditioner. Possible options are n (no
preconditioner), b (block-diagonal preconditioner) and d (diagonal
preconditioner).

Qualifiers:

\textbf{!MAXIT \textless number of iterations\textgreater{}} The
optional !MAXIT qualifier in the SOLVING section can be used to set the
maximum number of iterations to be used. If !MAXIT is not specified, the
default maximum number of iterations is 5,000.

\textbf{!STOPCRIT \textless convergence criterion\textgreater{}} If the
convergence criterion needs to be different from 1.0E-04, it can be set
with the optional !STOPCRIT qualifier in the SOLVING section.

\textbf{!NOPEEK} MiXBLUP stores intermediate results by default every
100th iteration. All solutions files are created and starting values for
a restart are stored as if solutions have converged. By default, only
the last set of preliminary results is kept. The name of each of the
file is the normal file name extended with \_PEEK, so for example
Solani\_PEEK.txt and solunf\_PEEK. The last set of preliminary results
will be removed when convergence has been attained or the maximum number
of iterations reached The process of storing preliminary results can be
avoided by specifying !NOPEEK.

\textbf{!PEEKFIRST \textless iteration number\textgreater{}} The
iteration number at which the preliminary results are stored for the
first time can be specified with !PEEKFIRST.

\textbf{!PEEKEVERY \textless number of iterations\textgreater{}} The
number of iterations between storing two subsequent sets of preliminary
results can be specified with !PEEKEVERY.

\textbf{!PEEKKEEP} Instead of only keeping the last set of preliminary
results, MiXBLUP can also retain each set of preliminary results. In
this case, the name of each file is the normal file name extended with
the iteration number, so for example Solani\_100.txt and
solunf\_100.txt. This option can be specified with !PEEKKEEP. This
option is useful for investigating the causes of unexpected convergence
behaviour.

\textbf{!RESTART} The optional qualifier !RESTART can be used to specify
that preliminary solutions of an interrupted analysis or old solutions
of the previous analysis are to be used as starting values for the new
evaluation. This option links old solutions to class effect levels using
the coded labels.

\textbf{!GFROMDISK} The !GFROMDISK qualifier instructs the solver to
read the inverse genomic relationship matrix from disk during solving.
This was the only option in previous versions of MiXBLUP. The new
default is to keep this matrix in memory, which is more demanding for
memory requirement, but it saves the time to read this matrix every
iteration.

\textbf{!WITHINBL \textless option\textgreater{}} The optional qualifier
!WITHINBL is used in the PRECON section and can be used to use a
different preconditioner for the within-block effects than the default
preconditioner type. Valid options are b (block-diagonal) and d
(diagonal).

\textbf{!ACROSSBL \textless option\textgreater{}} The optional qualifier
!ACROSSBL is used in the PRECON section and can be used to use a
different preconditioner for the across-block effects than the default
preconditioner type. Valid options are f (full), m (mixed), b
(block-diagonal) and d (diagonal).

\subparagraph{8.1.2.2. Syntax when using hpblup solver}\label{Cont06}

\begin{quote}
SOLVING\\
{[}!hpblup{]}\\
{[}!hpCriterion \textless type of convergence criterion\textgreater{]}\\
{[}!NumProc \textless number of cpu\textgreater{]}\\
{[}!MAXIT \textless number of rounds\textgreater{]}\\
{[}!STOPCRIT \textless convergence criterion\textgreater{]}\\
{[}!STARTVAL\_CHECK{]}\\
{[}!MIRACULIX{]}\\
{[}!MIRACULIX\_GPU{]}\\
{[}!NOPEEK{]}\\
{[}!PEEKFIRST \textless iteration number\textgreater{]}\\
{[}!PEEKEVERY \textless number of rounds\textgreater{]}\\
{[}!PEEKKEEP{]}
\end{quote}

Additional qualifiers:

\textbf{!hpblup} This qualifier is used to triggers MiXBLUP to call the
hpblup solver instead of the default MiX99 solver

\textbf{!hpCriterion} This qualifier is used to specify the convergence
criterion to be used, ck, cr or cd (default).

\textbf{!STARTVAL\_CHECK} The optional qualifier !STARTVAL\_CHECK only
applies to the hpblup solver and links old solutions to class effect
levels using the original class effect labels. Compared to the !RESTART
option, the !STARTVAL\_CHECK option will do an extra step of decoding
class effect levels of old solutions and re-coding them given the new
set of data. The !RESTART option can be used for the hpblup solver, but
old solutions will be linked to incorrect class effect levels if any
previously used class effect levels with a solution are no longer
present. In that case, solutions are still correct, provided that they
have converged, but it takes more iterations to obtain them than using
the correct old solutions as starting values.

\textbf{!NumProc } This qualifier is optional and can be used to specify
the number of cpus to be allocated to hpblup. This number has to be
equal to or lower than number available for the evaluation.

For the meaning of the other qualifiers, see previous chapter.

\subsubsection{8.2. Control of output}\label{Cont07}

\paragraph{8.2.1. General}\label{Cont08}

A successful analysis produces at least a log file and files with
solutions to all effects in the model. In some cases, additional results
may be required for development or evaluation purposes. Various options
are available to specify these additional files when required.

\paragraph{8.2.2. Syntax}\label{Cont09}

\subparagraph{8.2.2.1. Syntax when using MiX99 solver}\label{Cont10}

\begin{quote}
SOLVING\\
{[}!BASEANIMALSZERO \textless filename\textgreater{]}\\
{[}!YHAT{]}\\
{[}!EHAT{]}\\
{[}!YIELDDEV{]}\\
{[}!IDD{]}\\
{[}!DYD{]}\\
{[}!KEEPTMP{]}\\
{[}!SELINDEX \textless filename\textgreater{]}\\
{[}!FILTER \textless label\textgreater{]}\\
TMPDIR \textless work directory\textgreater{}
\end{quote}

Sections:

\textbf{TMPDIR} The TMPDIR section can be used to specify an existing
folder to store the temporary files of the kernel.

Qualifiers:

\textbf{!BASEANIMALSZERO \textless filename\textgreater{}} The estimated
breeding values of each individual in Solani.txt or Solani.out can be
presented as a deviation of a genetic base, which is the average of a
specified group of individuals, which are referred to as base animals.
The ID of the base animals should be given in the specified file, or in
BaseAnimals.dat. The file should contain one ID per row. The average per
trait of the group of base animals is included in the log file
MiXBLUP.log. If not all individuals are present in the data file, then a
warning is given in the log file MiXBLUP.log. The specified file
BaseAnimals.dat should contain the original IDs as they appear in the
pedigree file.

\textbf{!YHAT} The optional qualifier !YHAT creates a prediction for
each observed trait and each animal in the data file. The predictions
are stored in Yhat.txt, which is a text file that contains the animal ID
in the first column and predicted observations for each trait in the
model in subsequent columns. Missing observations in the data file get
the code -8192.0 in the file with predictions.

\textbf{!EHAT} The optional qualifier !EHAT stores the residual term of
each observed trait and each animal in the data file. The residuals are
stored in Ehat.txt, which is in text format and contains the animal ID
in the first column and residuals for each trait in the model in
subsequent columns. Missing observations in the data file get the code
-8192.0 in the file with residuals.

\textbf{!YIELDDEV} The optional qualifier !YIELDDEV stores the
observation corrected for fixed effects and non-genetic random effects
for each observed trait and animal in the data file. The yield
deviations are stored in YD.txt, which is in text format and contains
the animal ID in the first column and yield deviations for each trait in
the model in subsequent columns. Missing observations in the data file
get the code -8192.0 in the output file.

\textbf{!IDD} The optional qualifier !IDD (``individual daughter
deviation'') stores the yield deviation corrected for the genetic
contribution of the dam for each observed trait and animal in the data
file. The individual sire progeny deviations are stored in IDD.txt,
which is in text format and contains the animal ID in the first column
and yield deviations for each trait in the model in subsequent columns.
Missing observations in the data file get the code -8192.0 in the output
file.

\textbf{!DYD} The optional qualifier !DYD (``daughter yield deviation'')
stores the individual yield deviations averaged by sire. The daughter
yield deviations are stored in soldyd.txt, a text file that contains
sire in the first column, the number of progeny included in the progeny
yield deviation in the fourth column, the progeny yield deviations by
trait followed by the same number of fields to indicate whether the
progeny yield deviation of the corresponding trait is valid (value is 1)
or not (value is 0).

\textbf{!KEEPTMP} The optional qualifier !KEEPTMP can be used to stop
the removal of temporary files at the end of an analysis, for example to
check for possible errors. The default is that all large temporary files
are deleted as soon as they are no longer required.

\textbf{!SELINDEX \textless filename\textgreater{}} The qualifier
!SELINDEX can be used to automatically calculate a selection index value
as the sum of weighted genetic solutions (weighted EBV). The selection
index value is added as an additional column in the Solani output file.
The file specified after the qualifier contains the selection index
weighting factor for each combination of genetic effect and trait in the
model. The syntax is \textless trait\textgreater(\textless genetic
effect\textgreater) \textless selection index weighting
factor\textgreater, for example: phen1(animal) 1.0.

\textbf{!FILTER \textless label\textgreater{}} The qualifier !FILTER can
be used to automatically prune the Solani output file. It can be used to
include or exclude individuals from the solutions file. The pedigree
file should contain a field either with the text
\textless label\textgreater{} or something else. To exclude individuals
marked with `dead', use !FILTER \textless\textgreater dead. To include
individuals born in 2025 only, use !FILTER 2025.

\subparagraph{8.2.2.2. Syntax when using hpblup solver}\label{Cont11}

\begin{quote}
SOLVING\\
{[}!BASEANIMALSZERO \textless filename\textgreater{]}\\
{[}!YHAT{]}\\
{[}!EHAT{]}\\
{[}!YIELDDEV{]}\\
{[}!KEEPTMP{]}\\
{[}!SELINDEX {]}\\
{[}!FILTER \textless label\textgreater{]}\\
TMPDIR \textless work directory\textgreater{}
\end{quote}

\hyperref[Tabl01]{Back to Table of Contents}

\newpage

\subsection{9. Reliabilities}\label{Reli01}

\textbf{Besides estimating genetic effects (or breeding values), MiXBLUP
supports a second type of analysis to quantify the amount of information
available to estimate the genetic effect of each individual. This is
expressed as the reliability of the estimated (genomic) breeding value.
This chapter describes how reliabilities can be calculated with
MiXBLUP.}

\subsubsection{9.1. General}\label{Reli02}

A reliability is a measure of the information that is available for the
estimate of a breeding value. The reliability is dependent on the
heritability and the presence of observations for the individual itself.
The biggest impact on the reliability comes from the number of progeny
with observations. Calculation of reliability is not supported for
weighted residuals or marker haplotype models. Exact reliabilities can
be calculated only for a relatively small number of individuals in the
pedigree, say less than 100,000 individuals. Of these individuals, no
more than say 40,000 individuals can have a genotype record. If there
are more individuals in the evaluation, it is possible to calculate
approximate reliabilities. Approximate reliabilities are built up from
approximate pedigree reliabilities and the additional information
provided by genomic relationships as a deviation from pedigree
relationships.

MiXBLUP supports the calculation of an approximate or exact pedigree or
genomic reliability of the EBVs of individuals for most statistical
models.

{\def\LTcaptype{none} % do not increment counter
\begin{longtable}[]{@{}
  >{\raggedright\arraybackslash}p{(\linewidth - 4\tabcolsep) * \real{0.3333}}
  >{\raggedright\arraybackslash}p{(\linewidth - 4\tabcolsep) * \real{0.3333}}
  >{\raggedright\arraybackslash}p{(\linewidth - 4\tabcolsep) * \real{0.3333}}@{}}
\toprule\noalign{}
\begin{minipage}[b]{\linewidth}\raggedright
Type of evaluation
\end{minipage} & \begin{minipage}[b]{\linewidth}\raggedright
Reliabilities
\end{minipage} & \begin{minipage}[b]{\linewidth}\raggedright
Remarks
\end{minipage} \\
\midrule\noalign{}
\endhead
\bottomrule\noalign{}
\endlastfoot
Basic statistical model & direct & exact or approximate \\
Maternal genetic model & direct \& maternal & exact or approximate \\
Social interaction model & direct \& indirect & exact or approximate \\
Random regression model & direct & approximate; only for calculated
totals \\
Pedigree relationships & direct \& indirect & exact or approximate \\
Genomic relationships & pedigree + genomic & exact or approximate;
requires hpblup \\
SNP covariates & pedigree + genomic & exact or approximate; requires
hpblup \\
\end{longtable}
}

\subsubsection{9.2. Exact reliabilities}\label{Reli03}

\paragraph{1.1.1. General}\label{Reli04}

Exact reliabilities are calculated from the prediction error variance of
an estimate. Prediction error variance is obtained from the diagonal of
the inverse of the coefficient matrix. The need of a full inverse limits
the size of a dataset for which exact reliabilities can be calculated. A
prediction error variance is available for each genetic and non-genetic
solution of the evaluation. Reliabilities are available only for the
solutions of genetic effects of individuals. Exact reliabilities are
pedigree reliabilities if no genomic information is used and genomic
reliabilities if it is used. Exact reliabilities and prediction error
variances are only available for the hpblup solver.

\paragraph{9.2.1. Syntax}\label{Reli05}

\begin{quote}
SOLVING\\
!hpblup\\
!reliab\_exact
\end{quote}

Qualifier:

\textbf{!reliab\_exact} The !reliab\_exact is used to specify the
calculation of exact reliabilities for genetic effects and the
prediction error variance for each solution in the analysis. It has to
be used in tandem with !hpblup.

\paragraph{1.1.1. Associated output files}\label{Reli06}

{\def\LTcaptype{none} % do not increment counter
\begin{longtable}[]{@{}ll@{}}
\toprule\noalign{}
Output file & Description \\
\midrule\noalign{}
\endhead
\bottomrule\noalign{}
\endlastfoot
aggregate\_pev.dat & Prediction error variance of each solution \\
Relani.out & Reliabilities of genetic effects of individuals \\
\end{longtable}
}

\subsubsection{9.3. Approximate reliabilities}\label{Reli07}

\paragraph{9.3.1. General}\label{Reli08}

Approximate pedigree reliabilities are calculated for families within
blocks. It is a completely different process than estimating breeding
values. Approximate reliabilities are calculated in a separate analysis.

\paragraph{1.1.1. The concept of blocks in the reliability calculation
in MiXBLUP}\label{Reli09}

\subparagraph{9.3.1.1. Block variable}\label{Reli10}

The calculation of reliabilities in MiXBLUP requires the use of a family
block variable. The objective of using a block variable is to minimise
the use of memory during the calculation by ordering of equations in
equation family blocks (Lidauer and Strandén, 1999). Using a block
variable has no benefit for breeding value estimation with MiXBLUP. To
take advantage of this concept, it is important that each equation
family block contains as many closely connected equations as possible,
and that the number of equations connecting equation family blocks is as
low as possible. A family consists of an individual, its parents and its
progeny.

\subparagraph{9.3.1.2. Common-block variable}\label{Reli11}

Equations that connect all (or many) equations, such as individuals with
progeny in many different blocks, can be grouped into one or several
common blocks. Common blocks refer to blocks of equations, which will be
kept in memory for all families. Equations of common blocks must be
ordered to appear at the bottom of the MME, i.e.~animals of common
blocks must have largest block sorting variables. This ordering of mixed
model equations yields for many animal breeding problems a structured
coefficient matrix that has nearly doubly-bordered block diagonal form.
The number of common blocks can be specified by the user. The default is
that no common blocks are used. MiXBLUP will solve the mixed model
equations even if such a structure cannot be achieved. However,
pre-processing time may be longer than usual. As a consequence, it is
important to keep this concept in mind when editing the data for the
approximation of reliabilities.

\subparagraph{9.3.1.3. Sorting data and pedigree file on block
variable}\label{Reli12}

The concept of equation family blocks requires that data and pedigree
records be sorted on the block variable. MiXBLUP automatically checks
whether files are sorted and sorts them if necessary.

\subparagraph{9.3.1.4. Strategies for block definition}\label{Reli13}

For example in dairy cattle, equations for animals in the same herd
represent an equation family. Many models across species contain such
effects and are therefore suitable block variables to be used to group
equations into equation families. A good block variable orders the
records such that all (or almost all) records of the same animal and its
close relatives (parents and progeny) are in the same block. If the data
does not contain such a variable, it might be possible to generate a
suitable blocking variable. Again in dairy cattle, if a model contains a
herd-year-season effect (such like a herd-test-day) but not a herd
effect, it is advisable to include the herd code into the data and use
it as the blocking variable. MiXBLUP reads the data by blocks with one
or several blocks at a time. If there is only one block in a large data
file, all iteration files are read into the random access memory at the
same time, which might exhaust computer resources. If the data can be
read into the memory then this may be sensible. When benefits of using
equation family structure are desired, the block code of the animal has
to be given in the pedigree file. Each animals block code needs to be
the same as the one specified in the data file. Animals with records in
different data blocks (e.g.~in different herds) have to be coded with
the code of one of the different data blocks where it has observations,
e.g., the block with most of its observations. If an animal does not
have an observation, but it is a parent to an animal with observations
in the data file (e.g.~pedigree animal of a particular herd), then it
should receive the same block code as its offspring. This is most
suitable for a cow without observations. It should be assigned to a
block having most of its daughters. When an animal does not belong to
any equation family (no observations to give block code), or it is in
many different families through relationship information (e.g.~dairy
sires have progeny in many herds), a new block code should be given. It
is recommended to use a separate block code for animals with links to
many different equation family blocks. For instance, sires in a dairy
cattle population can be assigned to one block. These blocks should be
defined as common blocks and largest block code variables have to be
given to these blocks. Thus, sorting by the block code variable will
ensure that animals of common blocks will appear at the bottom of the
MME. Animals that cannot be included into any equation family can be
grouped into one or several own groups, depending on the number of such
animals. An equation family should always have a reasonable size. For
example, if the pedigree has equation families with 50 to 2000 animals
per block and a block with 300,000 animals, it is advisable to split the
largest block into several smaller blocks. The MiXBLUP kernel reads as
many animal blocks at a time as possible, and the largest animal block
dictates the memory requirements.

\paragraph{9.3.2. Differences between the syntax of the instruction file
for approximate reliability calculation and breeding value
estimation}\label{Reli14}

\subparagraph{9.3.2.1. Data file}\label{Reli15}

For a reliability calculation, every animal in the evaluation has to be
uniquely assigned to a level of the block variable. This block variable
must be present in the data file.

\subparagraph{9.3.2.2. Genetic similarity between
individuals}\label{Reli16}

The level of the block variable of an animal is also required in its
record in the pedigree file. Genetic groups are ignored in the
reliability calculation and replaced with unknown parents.

\subparagraph{9.3.2.3. Statistical model}\label{Reli17}

The statistical model for reliability calculation contains a single
fixed effect that is treated as nested within blocks, even if it is an
across-block effect. It should be the fixed effect with the largest
impact on the reliability, so the lowest average number of observations
within a class. The use of !WEIGHT is not supported.

\subparagraph{9.3.2.4. Control of analysis}\label{Reli18}

A reliability calculation is triggered with !RELIABILITY in the SOLVING
section.

\paragraph{1.1.1. Syntax}\label{Reli19}

\begin{quote}
DATAFILE \textless file name\textgreater{}\\
\textless ID\textgreater{} I/A\\
\textless\ldots\textgreater{}\\
\textless name block code\textgreater{} I !BLOCK\\
PEDFILE \textless file name\textgreater{}\\
\textless ID\textgreater{} I/A\\
\textless\ldots\textgreater{}\\
\textless block code\textgreater{} I !BLOCK\\
MODEL\\
\textless trait1\textgreater{} \textasciitilde{} BL(\textless largest
fixed class effect trait1\textgreater) !RANDOM \textless random
effects\textgreater{} G(\textless ID\textgreater)\\
\textless trait2\textgreater{} \textasciitilde{} BL(\textless largest
fixed class effect trait2\textgreater) !RANDOM \textless random
effects\textgreater{} G(\textless ID\textgreater)\\
SOLVING\\
!RELIABILITY\\
{[}!NCOMBLK \textless number\textgreater{]}\\
{[}!KEEPTMP{]}
\end{quote}

Qualifier:

\textbf{!NCOMBLK \textless number\textgreater{}} If common blocks are
used, the qualifier !NCOMBLK should be used to specify how many common
blocks there are. Common blocks should have a block identification that
positions the block at the end of the list of blocks. The specified
number of blocks at the end of the list of blocks are considered common
blocks.

\textbf{!KEEPTMP} The optional qualifier !KEEPTMP can be used to stop
the removal of temporary files at the end of an analysis, for example to
check for possible errors. The default is that all large temporary files
are deleted as soon as they are no longer required.

Qualifiers other than !NCOMBLK and !KEEPTMP have no meaning when
!RELIABILITY is specified, and will be ignored.

\paragraph{1.1.1. Associated output files}\label{Reli20}

{\def\LTcaptype{none} % do not increment counter
\begin{longtable}[]{@{}
  >{\raggedright\arraybackslash}p{(\linewidth - 2\tabcolsep) * \real{0.5000}}
  >{\raggedright\arraybackslash}p{(\linewidth - 2\tabcolsep) * \real{0.5000}}@{}}
\toprule\noalign{}
\begin{minipage}[b]{\linewidth}\raggedright
Output file
\end{minipage} & \begin{minipage}[b]{\linewidth}\raggedright
Description
\end{minipage} \\
\midrule\noalign{}
\endhead
\bottomrule\noalign{}
\endlastfoot
Relani.txt & The reliability of direct genetic effects of traits in the
model. The exact layout of Relani.txt is printed at the end of
reliabilities.log. \\
Relani.out & As Relani.txt, but for alphanumerical ID labels \\
Relani\_indirect.txt & The reliability of indirect genetic effects of
traits in the model. The exact layout of Relani\_indirect.txt is printed
at the end of reliabilities.log. \\
Relani\_indirect.out & As Relani\_indirect.txt, but for alphanumerical
class labels \\
\end{longtable}
}

\subsubsection{9.4. Approximate genomic reliabilities}\label{Reli21}

\paragraph{1.1.1. General}\label{Reli22}

MiXBLUP 3.0 contained two approximations of genomic reliabilities
proposed by Misztal et al.~(2013; J. Dairy Sci. 96 :647--654). They have
been removed from MiXBLUP 3.2, as they were not useful in practice.
Instead, the algorithm of Gao et al.~(2023; Genetics Selection Evolution
55:1) is now used. This algorithm first calculates approximate pedigree
reliabilities as in Chapter 9.3 and then calculates the approximate
genomic contribution to the reliability of a genomic estimated breeding
value.

\paragraph{1.1.1. Instruction file}\label{Reli23}

The instruction file for approximate genomic reliabilities differs from
a breeding value estimation in the same way as described above, but with
an ERMFILE section with a genotype file added.

\paragraph{1.1.1. Syntax}\label{Reli24}

\begin{quote}
SOLVING !greliability
\end{quote}

Qualifiers:

\textbf{!greliability} The qualifier !greliability is used to calculate
approximate genomic reliabilities.

\paragraph{1.1.1. Associated output files}\label{Reli25}

{\def\LTcaptype{none} % do not increment counter
\begin{longtable}[]{@{}ll@{}}
\toprule\noalign{}
Output file & Description \\
\midrule\noalign{}
\endhead
\bottomrule\noalign{}
\endlastfoot
Relani\_ssg.out & Approximate genomic reliabilities \\
Relani.out & Approximate pedigree reliabilities \\
\end{longtable}
}

\subsubsection{9.5. Command-line interface for calculating
reliabilities}\label{Reli26}

\paragraph{1.1.1. Syntax}\label{Reli27}

The recommended syntax for calling MiXBLUP to calculate reliabilities
is: \textgreater MiXBLUP.exe rel --type \textless reliability
type\textgreater{} \textless instruction file\textgreater{}

The old syntax is also still supported: \textgreater MiXBLUP.exe
\textless instruction file\textgreater{}

\textbf{\textless reliability type\textgreater{}} \textbar Type
\textbar{} Equivalent Flag \textbar{} Description\textbar{} \textbar{}
--- \textbar{} --- \textbar{} \textbar exact \textbar{} !reliab\_exact
\textbar{} Exact reliabilities\textbar{} \textbar pedigree \textbar{}
!reliability \textbar{} Approximate pedigree reliabilities\textbar{}
\textbar genomic \textbar{} !greliability \textbar{} Approximate genomic
reliabilities\textbar{}

\hyperref[Tabl01]{Back to Table of Contents}

\newpage

\subsection{10. Multi-trait genome-wide association studies
(GWAS)}\label{MTGW01}

\textbf{A Genome-Wide Association Study (GWAS) is an approach that
examines the entire genome to identify genetic variations associated
with a particular complex trait of interest. It involves scanning the
genomes of many individuals to find genetic differences, typically
single nucleotide polymorphisms (SNPs), that occur more frequently in
individuals with a specific complex trait of interest. By identifying
these associated genetic variants, breeders can gain insights into the
underlying biology of the complex trait of interest.}

\subsubsection{10.1. General}\label{MTGW02}

Genome-wide association study (GWAS) is a common tool in genetic
research for identifying loci associated with complex traits. However,
the increasing availability of single nucleotide polymorphism (SNP)
genotypes and whole-genome sequencing (WGS) data presents significant
computational challenges due to the large number of individuals and
SNPs. Traditional mixed linear model association (mlma) analyses, as
implemented in for example the software GCTA, while considered the gold
standard, are computationally intensive. For limited datasets, such as
those with a single trait of interest and up to 40,000 genotyped
individuals, MiXBLUP (Chapter 10.2) can be used to obtain frequentist
p-values needed for GWAS using a single-step genomic BLUP. The
implemented approach is only feasible for a single trait and up to
40,000 genotyped individuals. The computational limitations are removed
by using an approximate GWAS based on the solutions of single-step SNP
best linear unbiased prediction (ssSNPBLUP) or single-step genomic BLUP
using a component-wise Ta decomposition of the inverse of G
(ssGTacBLUP). The approximate GWAS approach implemented in MiXBLUP
(Chapter 10.3) includes two steps. First, SNP effects are estimated by
solving ssSNPBLUP with a preconditioned conjugate gradient method.
Second, prediction error variances for all SNP effects, needed for
computing frequentist p-values, are approximated by the diagonal
elements of the inverse of the coefficient matrix of a SNPBLUP model.
For efficiency, a sliding-window approach is used assuming that linkage
disequilibrium between SNPs that are more than 50,000 SNPs apart is
zero, resulting in multiple inversions of coefficient matrices of size
equal to 50,000, instead of one single inversion of a matrix
encompassing all SNPs.

\subsubsection{10.2. Computation of frequentist p-values for limited
datasets}\label{MTGW03}

\paragraph{10.2.1. General}\label{MTGW04}

For limited datasets, such as those with a single trait of interest and
up to 40,000 genotyped individuals, MiXBLUP can be used to obtain
frequentist p-values needed for GWAS using a single-step genomic BLUP.
The implemented approach requires computation of the inverse of the
coefficient matrix of the evaluation, which limits its application to a
single trait and up to 40,000 genotyped individuals.

\paragraph{10.2.2. Syntax for calculating frequentist
p-values}\label{MTGW05}

\begin{quote}
ERMFILE \textless genotype file\textgreater{} !Construct SSmat
!SingleStep !METHOD VanRaden !NoScale !NoReg\\
\textless\ldots\textgreater{}\\
SOLVING\\
!hpblup\\
!pvalue\_exact
\end{quote}

Qualifier:

\textbf{!pvalue\_exact} The qualifier !pvalue\_exact is used to specify
the calculation of frequentist p-values through full inversion of the
coefficient matrix of ssGBLUP. For calculating frequentist p-values with
a single-step genomic evaluation, it is recommended to use a genomic
relationship matrix computed following the first approach of VanRaden
(2028) (!METHOD VanRaden) together with the !NoScale and !NoReg options.
This option is only available with the solver hpblup.

\paragraph{10.2.3. Associated output files}\label{MTGW06}

Frequentist p-values are saved in a file called Pvalue.out. The format
of the file is the same as the one of Relani.out. The order of the
frequentist p-values follow the order of the SNPs in the genotype file.

\paragraph{10.2.4. Example (move to appendix later)}\label{MTGW07}

\begin{quote}
TITLE GWAS for Trait1\\
DATAFILE Data.txt !MISSING -9\\
fixeff1 I\\
ID1 I\\
trait1 T\\
ERMFILE Genotype.bed !Plink !CONSTRUCT SSmat\\
ID1 I\\
!SingleStep\\
!LAMBDA 1.0\\
!ALPHA 0.95\\
!BETA 0.05\\
!OMEGA 1.0\\
!MAF 0.005\\
!NOSCALE\\
!NOREG\\
!METHOD VanRaden\\
PEDFILE Pedigree.txt !CalcInbr\\
ID1 I\\
sire I\\
dam I\\
PARFILE ../para.var\\
MODEL\\
trait1 \textasciitilde{} fixeff1 !RANDOM G(ID1)\\
SOLVING\\
!hpblup\\
!pvalue\_exact\\
END
\end{quote}

\subsubsection{10.3. Approximation of frequentist p-values for
large-scale datasets}\label{MTGW08}

\paragraph{10.3.1. General}\label{MTGW09}

The approximate GWAS approach implemented in MiXBLUP includes two steps.
First, SNP effects are estimated by solving ssSNPBLUP or ssGTAcBLUP with
a preconditioned conjugate gradient method. See Chapter X for more
details on how to run a ssSNPBLUP or ssGTAcBLUP evaluation with MiXBLUP.
Second, prediction error variances for all SNP effects, needed for
computing frequentist p-values, are approximated by the diagonal
elements of the inverse of the coefficient matrix of a SNPBLUP model.
For efficiency, a sliding-window approach is used assuming linkage
disequilibrium between SNPs more than 50,000 SNPs apart is zero,
resulting in multiple inversions of coefficient matrices of size equal
to 50,000, instead of one single inversion of a matrix encompassing all
SNPs. This approach allows MiXBLUP to approximate frequentist p-values
for large-scale multi-trait single-step genomic evaluations.

\paragraph{10.3.2. Syntax for calculating frequentist
p-values}\label{MTGW10}

\begin{quote}
ERMFILE \textless genotype file\textgreater{} !Construct SSmat
!SingleStep !Tac\\
\textless\ldots\textgreater{}\\
SOLVING\\
!hpblup\\
!gwas \textless file with SNP effect solutions\textgreater{}
\textless SNP effect ID\textgreater{}
\end{quote}

Or

\begin{quote}
SNPFILE \textless genotype file\textgreater{} !NoPrune !NoCheck\\
\textless\ldots\textgreater{}\\
SOLVING\\
!hpblup\\
!gwas \textless file with SNP effect solutions\textgreater{}
\textless SNP effect ID\textgreater{}
\end{quote}

Qualifier:

\textbf{!gwas }

The qualifier !gwas is used to approximate frequentist p-values for
large-scale datasets. The first field corresponds to the file with
estimated SNP effects obtained from a previous ssSNPBLUP or ssGTAcBLUP
evaluation. This file is called Solreg\_mat.txt. The second field
corresponds to the first ID of the SNP effect. The approximation of
frequentist p-values for large-scale datasets relies on the
approximation of genomic reliabilities (Chapter 9). Therefore,
approximating frequentist p-values requires an instruction file for
approximating genomic reliabilities, with the qualifier !gwas (instead
of !greliability) in the SOLVING section.

SNP effect solutions are stored in Solreg\_mat.txt, both for ssGBLUP,
using a Ta or Tac decomposition of G, and ssSNPBLUP. / is the hpblup
EFFECT\_ID of the SNP effect, which can be found in the hpblup
instruction file hpInstr.txt.

\paragraph{10.3.3. Associated output files}\label{MTGW11}

Approximate p-values are saved in a file called Pvalue\_approx.out. The
format of the file is the same as the one of Relani.out. The order of
the approximate p-values follows the order of the SNPs in the genotype
file.

\paragraph{10.3.4. Example (move to appendix later)}\label{MTGW12}

\begin{quote}
TITLE Approximate GWAS for Trait1\\
DATAFILE data.txt !MISSING -9\\
fixeff1 I\\
blockid I !BLOCK\\
ID1 I\\
trait1 T\\
ERMFILE Genotype.bed !Plink !CONSTRUCT SSmat !SingleStep !NumProc 10\\
ID1 I\\
!LAMBDA 1.0\\
!ALPHA 0.95\\
PEDFILE pedigree.txt !CalcInbr\\
ID1 I\\
sire I\\
dam I\\
blockped I !BLOCK\\
PARFILE ../para.var\\
MODEL\\
trait1 \textasciitilde{} BL(fixeff1) !RANDOM G(ID1)\\
SOLVING\\
!gwas Solreg\_mat.txt 3\\
END
\end{quote}

\hyperref[Tabl01]{Back to Table of Contents}

\newpage

\subsection{11. Descriptive analyses}\label{Desc01}

\textbf{To be able to configure a genetic evaluation appropriately, and
interpret the results correctly, it is important to understand the data
structure in detail. MiXBLUP provides tools to analyze and describe the
data structure.}

\subsubsection{11.1. General}\label{Desc02}

Optimal genetic and genomic evaluations require sufficient valid data
records for each trait, sufficient genetic and genomic relationships in
the data and sufficient information to fit the model chosen (Chapter
\hyperref[Desc03]{11.2.}). A more specific detail is whether the
objective choice of base populations in the pedigree is sufficiently
supported by the data (Chapter \hyperref[Desc11]{11.3.}).

\subsubsection{11.2. Descriptive statistics of performance data, genomic
data and pedigree}\label{Desc03}

\paragraph{11.2.1. General}\label{Desc04}

Aspects of data structure that are important to consider for optimal
genetic evaluations are the amount of valid information per trait,
amount of genetic and genomic relationships in the data, and amount of
information available for the model fitted.

\paragraph{11.2.2. Syntax}\label{Desc05}

\begin{quote}
DATAFILE \textless file name\textgreater{} !Stats N {[}D H L{]}
\end{quote}

Qualifier:

\textbf{!Stats} The !Stats qualifier can be used to specify the
calculation of descriptive statistics of the data structure. It can be
used with either the default or the hpblup solver. There are four types
of statistics that can be produced: N for numbers of records in data,
pedigree and genotype file (Chapter \hyperref[Desc06]{11.2.3.}) ; D for
means and standard deviations of traits and covariates in the data
(Chapter \hyperref[Desc07]{11.2.4.}); H for grouping class effect levels
for each trait by the number of records per class (Chapter
\hyperref[Desc08]{11.2.5.}) and L for a table by trait with number of
records for each class effect level (Chapter
\hyperref[Desc09]{11.2.6.}). For large evaluations, it is recommended to
use !STATS NDH, as the option L might produce a very large output file.
Types may be specified in any order. If D, H or L are specified, N is
automatically included.

\paragraph{11.2.3. Counts (option N)}\label{Desc06}

The counts option (option N) produces five tables. The first table shows
the number of records in data, pedigree and genotype file. It provides
some basic information on the amount of available information in the
input files. The second table presents a breakdown of the number of
records per trait. It provides the number of data records, the number of
individuals with a data record, the number of individuals with data and
genotype, which is the size of the reference population of the trait,
and the number of invalidated records. Invalidation may occur if model
information is missing or if the observation is outside of the valid
range. The third table gives the number of levels of class effects in
the model across traits. This may be useful for software that estimate
variance components and require the number of level of each class effect
upfront. The fourth table shows pedigree completeness. The fifth table
provides some information on the use of base populations. See Chapter
\hyperref[Desc11]{11.3.} for more detailed information on the subjective
choice of base populations.

\paragraph{11.2.4. Means and standard deviations (option
D)}\label{Desc07}

Means, standard deviations, and minimum and maximum values (option D)
are presented in two tables, one for traits and one for covariates. If
minimum or maximum values fall outside the valid range, then the
qualifier !MinMax can be used to restrict the range in the data
evaluated.

\paragraph{11.2.5. Class effect levels grouped by available information
(option H)}\label{Desc08}

For the amount of information available for class effect levels (option
H), class effect levels are grouped by the number of records available
(0, 1, 2, 3-5, 6-10, 11-20, 21-50, 51-100, 101-200, 201-500, more than
500). The number of class effect levels in each category is presented
for each combination of class effect and trait. Many class effect levels
with only few records for a trait results in inaccurate estimates and
poor correction for systematic non-genetic factors. It may also increase
the likelihood of confounding between effect levels of different class
effects, which causes only the sum of the two effects to be estimable.
Small class effect levels are more of a problem for fixed effects than
random effects, because variance of level effects and a correlation or
genetic relationship structure between levels provide additional
information and constraints for the latter.

\paragraph{11.2.6. Valid data records for each combination of trait and
class effect level (option L)}\label{Desc09}

It is also possible to print a table with all class effect levels for
each combination of trait and class effect and the number of valid data
records available (option L). It can be used to identify in which part
of the data small class effect levels occur and which effect levels
should be combined to achieve a larger number of data records per level.
For large evaluations, the table and output file may become very large.

\paragraph{11.2.7. Associated output files}\label{Desc10}

{\def\LTcaptype{none} % do not increment counter
\begin{longtable}[]{@{}ll@{}}
\toprule\noalign{}
Output file & Description \\
\midrule\noalign{}
\endhead
\bottomrule\noalign{}
\endlastfoot
Statistics.log & Output file with tables for the options specified \\
\end{longtable}
}

\subsubsection{11.3. Diagnostics of use of base
populations}\label{Desc11}

\paragraph{11.3.1. General}\label{Desc12}

Any pedigree has individuals without known parents. In some cases it is
reasonable to assume that individuals without known originate from the
same large population, but in other cases there is prior knowledge that
these individuals come from different populations. Multiple base
populations are generally referred to as genetic groups or unknown
parent groups. The assumptions are that these groups are infinitely
large, so groups, and individuals within a group, are unrelated.
Discarded pedigree information and genomic information of descendants
may be used to take into account that base populations are finite in
reality. Genetic groups with relationships within and between groups are
referred to as metafounders (Legarra et al., 2015). Allocation of
individuals without known parents to a base population is a subjective
process. The aim is to stay as close as possible to the true base
population to minimize bias, but with sufficient individuals in each
base population to minimize the mean squared error. For metafounders, it
is important that there is sufficient information to estimate
relationships within and between base populations. MiXBLUP provides
statistics of quantity, quality, and proximity of genomic information
for base populations. Quantity of genomic information is expressed as
equivalent number of base animals genotyped (Chapter
\hyperref[Desc13]{11.3.2.}). Quality of genomic information is expressed
as auto-similarity with other base populations (Chapter
\hyperref[Desc14]{11.3.3.}). Proximity of genomic information is
expressed as number of generations between pedigree and genomic base
populations (Chapter \hyperref[Desc15]{11.3.4.}).

\paragraph{11.3.2. Equivalent number of base animals
genotyped}\label{Desc13}

The first statistic quantifies the amount of genotype information
available for each base population. For this purpose we define a genomic
base population which consists of all genotyped animals with a path of
non-genotyped ancestors to a base population. Genotyped animals that are
descendants of two genotyped parents are not included as they do not
provide genotype information to a base population (c equals 0.0).
Imputation of non-genotyped animals from genotyped descendants only
using pedigree relationships, causes loss of information, because the
genotype information has to be distributed to two parents in every
generation without genotypes. We therefore introduce equivalent number
of base animals genotyped for a base population (Neq\textsubscript{i}).
It is calculated as:

\pandocbounded{\includegraphics[keepaspectratio]{https://raw.githubusercontent.com/JantN1966/MasterManual/main/Images/DescStat01.jpg}}

where N is the number of animals in the genomic base population,
q\textsubscript{ij} is the fraction relating the contribution of base
population i to the total genetic value of the animal j,
gener\textsubscript{j} is the maximum number of generations between
animal j and base population i, and c\textsubscript{i}j is the
contribution of the genomic base animal j to base population i. So for
example, a genotyped individual has one non-genotyped parent, hence c
equals 0.50. The genotyped individual is linked to a specific base
population only through one grandparent, hence q is 0.25.

\paragraph{11.3.3. Auto-similarity to another base
population}\label{Desc14}

The second statistic is the extent to which the same genotype
information was used for a pair of base populations, which we call
auto-similarity of one base population to another. To illustrate the
concept, imagine 24 balls, 8 red, 8 blue and 8 orange. The balls are
placed in two bowls, so the first bowl contains 5 red and 4 blue balls.
The second contains the remaining 3 red, 4 blue and 8 orange. The
similarity of bowl one to bowl two is 3 red + 4 blue over 9 balls is
0.78. The similarity of the second bowl to the first one is 3 red + 4
blue over 15 balls is 0.47. The colors in the illustration are genotyped
animals and the bowls are base populations. So auto-similarity of base
population i to j (AS\textasciitilde i to j\textasciitilde) is
calculated as:

\pandocbounded{\includegraphics[keepaspectratio]{https://raw.githubusercontent.com/JantN1966/MasterManual/main/Images/DescStat02.jpg}}

where q\textsubscript{ik} and q\textsubscript{jk} and
c\textsubscript{ik} are defined as q\textsubscript{ij} and
c\textsubscript{ij} in Chapter \hyperref[Desc13]{11.3.2.}. An
AS\textasciitilde i to j\textasciitilde{} of 0 means that genotype
information available to two base populations is independent. A value of
1 means that genotype information available is identical.

\paragraph{11.3.4. Number of generations between pedigree and genomic
base populations}\label{Desc15}

The third statistic is the weighted average number of generations
between base animals in a base population and the genomic base
population. It is calculated as:

\pandocbounded{\includegraphics[keepaspectratio]{https://raw.githubusercontent.com/JantN1966/MasterManual/main/Images/DescStat03.jpg}}

where q\textsubscript{ij}, c\textsubscript{ij} and
gener\textsubscript{j} are defined as above. delta\textsubscript{Gener}
differentiates between many remote genotyped descendants and fewer
proximate ones for a given Neq.

\paragraph{11.3.5. Syntax}\label{Desc16}

\begin{quote}
PEDFILE \textless file name\textgreater{} !Groups 1.0 !DiagBasePop
\end{quote}

Qualifier: \textbf{!DiagBasePop} The qualifier !DiagBasePop can be used
to specify that additional statistics on the suitability of base
populations for use as genetic groups or metafounders. The !DiagBasePop
is currently only available for the hpblup solver.

\paragraph{11.3.6. Associated output files}\label{Desc17}

{\def\LTcaptype{none} % do not increment counter
\begin{longtable}[]{@{}
  >{\raggedright\arraybackslash}p{(\linewidth - 2\tabcolsep) * \real{0.5000}}
  >{\raggedright\arraybackslash}p{(\linewidth - 2\tabcolsep) * \real{0.5000}}@{}}
\toprule\noalign{}
\begin{minipage}[b]{\linewidth}\raggedright
Output file
\end{minipage} & \begin{minipage}[b]{\linewidth}\raggedright
Description
\end{minipage} \\
\midrule\noalign{}
\endhead
\bottomrule\noalign{}
\endlastfoot
DiagBasePop.log & Summary of diagnostics \\
DBP\_DescriptiveStats.log & Descriptive statistics of base
populations \\
DBP\_EquivalentNrBaseIndiv.log & Equivalent number of genotyped
individuals in base population \\
DBP\_AutoSimilarity.log & Auto-similarity of each base population with
other base populations \\
DBP\_WeightedNrGenerations.log & Number of generations between pedigree
and genomic base populations \\
\end{longtable}
}

\hyperref[Tabl01]{Back to Table of Contents}

\newpage

\subsection{12. Validation studies with MiXBLUP}\label{Vali01}

Estimated breeding values (EBV) are used to identify the best
individuals to become the parents of the next generation. The implicit
assumption is that EBV are a reasonable prediction of performance of
future progeny. BLUP evaluations to estimate breeding values often span
many generations. Statistical model and components of variance and
covariance among traits may have been suitable for earlier generations,
but to what extent do EBV of the current generation predict the
performance of progeny well? This is explored in validation studies.

\subsubsection{1.1 General}\label{Vali02}

It is good management practice to do validation studies regularly. In
practice, validation studies are quite cumbersome to do, involve a lot
of manual work and are prone to making mistakes. This makes it difficult
to compare results between validation studies. MiXValidate has been
designed to facilitate validation studies that are easy to do and
transparent to interpret. MiXValidate creates a partial dataset given
the validation individuals and the genetic effect to validate. It
verifies that the list of validation individuals does not contain
multiple generations. It estimates breeding values for the partial
dataset and for the full dataset. It calculates validation statistics
that are well-defined. Finally, it displays the amount of information
that is available for validating the EBV and it leaves out validation
individuals that do not have sufficient data available for validation.

\subsubsection{12.1. Validation individuals}\label{Vali03}

MiXValidate supports two types of validation studies. Parent validation
focuses on how well parent EBV predict progeny performance. Individual
validation checks how well EBV predict future performance of the
individual itself. Parent validation can be done for either sires or
dams. The user provides a file with validation individuals. This is a
text file with just a list of original IDs, one ID per line. Validation
individuals have to be from a single generation. MiXValidate gives a
warning if more than one generation of a family is present among the
validation individuals. The youngest individual from such a family will
have more data removed for validation than other validation individuals,
thus affecting the validation statistics.

\subsubsection{12.2. Validation effect}\label{Vali04}

There are traits that are affected by the genotype of other individuals,
in addition to the genotype of the individual itself. For example, daily
gain of piglets in the first weeks of their life is also determined by
the genotype of the dam. Social behavior of pen mates affects the growth
of individuals, in addition to their own genotype. These are referred to
as indirect genetic effects and direct genetic effects, respectively.
The process of creating the partial dataset is completely dependent on
which genetic effect needs to be validated. Data records are omitted
based on the field in the data file with the validation effect. A trait
with a statistical model with direct and indirect genetic effects
therefore needs two separate MiXValidate analyses for validation: one
for the direct genetic effect as validation effect and one for the
indirect genetic effect as validation effect.

\subsubsection{12.3. Creating partial dataset}\label{Vali05}

The idea behind creating the partial dataset is to create a dataset that
reflects the situation at the point of selection for breeding or further
testing. Any performance records on relatives collected after selection
of the individual should be removed. The common practice to use a
specific date after which all data records are omitted, only works well
in the case of non-overlapping generations. With individual validation,
MiXValidate removes data records of validation individuals and their
siblings, and of descendants of these two groups. With parent
validation, MiXValidate removes data records of progeny of validation
individuals, descendants of this progeny and descendants of siblings of
validation individuals. If the user specifies to remove data records of
validation individuals, then these will be removed, as well as data
records of siblings of validation individuals. The above can be applied
to any genetic effect as validation effect (Table 1).

\textbf{Table 1.} Type of data records and type of removal in partial
dataset

{\def\LTcaptype{none} % do not increment counter
\begin{longtable}[]{@{}
  >{\raggedright\arraybackslash}p{(\linewidth - 12\tabcolsep) * \real{0.1429}}
  >{\raggedright\arraybackslash}p{(\linewidth - 12\tabcolsep) * \real{0.1429}}
  >{\raggedright\arraybackslash}p{(\linewidth - 12\tabcolsep) * \real{0.1429}}
  >{\raggedright\arraybackslash}p{(\linewidth - 12\tabcolsep) * \real{0.1429}}
  >{\raggedright\arraybackslash}p{(\linewidth - 12\tabcolsep) * \real{0.1429}}
  >{\raggedright\arraybackslash}p{(\linewidth - 12\tabcolsep) * \real{0.1429}}
  >{\raggedright\arraybackslash}p{(\linewidth - 12\tabcolsep) * \real{0.1429}}@{}}
\toprule\noalign{}
\begin{minipage}[b]{\linewidth}\raggedright
Validation Effect
\end{minipage} & \begin{minipage}[b]{\linewidth}\raggedright
Direct
\end{minipage} & \begin{minipage}[b]{\linewidth}\raggedright
Direct
\end{minipage} & \begin{minipage}[b]{\linewidth}\raggedright
Maternal
\end{minipage} & \begin{minipage}[b]{\linewidth}\raggedright
Maternal
\end{minipage} & \begin{minipage}[b]{\linewidth}\raggedright
Maternal
\end{minipage} & \begin{minipage}[b]{\linewidth}\raggedright
Maternal
\end{minipage} \\
\midrule\noalign{}
\endhead
\bottomrule\noalign{}
\endlastfoot
\textbf{Validation Type} & \textbf{Parent} & \textbf{Individual} &
\textbf{Parent} & \textbf{Parent} & \textbf{Individual} &
\textbf{Individual} \\
\textbf{Data field} & \textbf{Direct} & \textbf{Direct} &
\textbf{Maternal} & \textbf{Direct} & \textbf{Maternal} &
\textbf{Direct} \\
Validation Individual & X1 & X2 & X1 & - & X2 & - \\
Siblings\textsuperscript{a} & X1 & X3 & X1 & - & X3 & - \\
Progeny\textsuperscript{b} & X2 & X3 & X2 & X3 & X3 & X3 \\
Descendants\textsuperscript{c} & X3 & X3 & X3 & X3 & X3 & X3 \\
\end{longtable}
}

\textsuperscript{a} Siblings of validation individuals\\
\textsuperscript{b} Progeny of validation individuals\\
\textsuperscript{c} Descendants of progeny or siblings of validation
individuals\\
X1 data records are removed if !ValNoOwnPerf is specified\\
X2 data records are removed and removed records are used for validation
statistics\\
X3 data records are removed and removed records are not used for
validation statistics

Data records of descendants are removed to allow the user to do
validation studies retrospectively, so in any generation before the
current generation. Descendants are not included in the validation
statistics. The user can also specify a list of individuals for which
all data records should be removed. MiXValidate will not remove any
additional records in the partial data. It can be used for example to
use a fixed date by adding all individuals with a data record after the
date to the remove list.

\subsubsection{12.4. Types of validation}\label{Vali06}

MiXValidate presents validation statistics of two types of forward
validation: LR method (Legarra and Reverter, 2018) and
adjusted-phenotype validation (Mäntysaari et al., 2010). The LR method
compares solutions of validation individuals from evaluation of the
partial data with the corresponding solutions from evaluation of the
full data. Adjusted-phenotype validation compares solutions of
validation individuals from evaluation of the partial data with omitted
adjusted phenotypes. Adjusted-phenotype validation is implemented for
direct genetic effects only, for now.

\subsubsection{12.5. Command-line syntax}\label{Vali07}

The recommended way to call MiXBLUP for a validation study is:
\textgreater MiXBLUP.exe val -i \textless instruction file\textgreater{}

The old syntax is now deprecated but still supported:
\textgreater MiXValidate.exe \textless instruction file\textgreater{}

\subsubsection{12.6. Syntax}\label{Vali08}

Starting point for the instruction file for MiXValidate is the MiXBLUP
instruction file of the routine evaluation. A new section VALIDATION
needs to be added and the qualifiers !YieldDev and !HeritabFile need to
be added to the SOLVING section.

\begin{quote}
SOLVING\\
{[}\ldots{]}\\
!YieldDev\\
!HeritabFile \textless name of file with heritability per
trait\textgreater{}\\
VALIDATION\\
!ValList \textless name of file with validation
individuals\textgreater{}\\
{[}!ValEffect \textless field name genetic effect to
validate\textgreater{]}\\
{[}!ValParent \textless sire \textbar{} dam\textgreater{]}\\
{[}!ValNoOwnPerf{]}\\
{[}!ValRemoveList \textless name of file with individuals to remove in
partial data\textgreater{]}\\
{[}!ValMinRec \textless n\textgreater{]}\\
{[}!ValSolutions \textless solutions file\textgreater{]}
\end{quote}

Qualifiers:

\textbf{!ValList \textless name file\textgreater{}} The !ValList
qualifier is mandatory and specifies the name of the file with
validation individuals. The file contains a record for each validation
individual with only the ID of the individual.

\textbf{!ValEffect \textless validation effect name\textgreater{}} The
!ValEffect qualifier is optional and specifies the genetic effect to
validate. Default is the direct genetic effect.

\textbf{!ValParent \textless sire \textbar{} dam\textgreater{}} The
!ValParent qualifier is optional and specifies that parent validation
will be used and for which parent. If the !ValParent qualifier is
omitted, then individual validation will be used. If the !ValParent
qualifier is used without a setting, then sire validation will be used.

\textbf{!ValNoOwnPerf} The !ValNoOwnPerf qualifier is optional and only
has a meaning in the context of parent validation. It specifies that
data records of validation individuals and their siblings will also be
removed in the partial data.

\textbf{!ValRemoveList \textless file with individuals to remove in
partial data\textgreater{}} The !ValRemoveList qualifier is optional and
causes MiXValidate to remove ONLY data records of the individuals in the
specified file, so no additional data records are removed.

\textbf{!ValMinRec \textless n\textgreater{}} The qualifier !ValMinRec
is optional and specifies the minimum number of data records available
for a validation individual to be included in the validation statistics.
This minimum number applies to own performance in case of individual
validation and progeny records in case of parent validation. Records of
other descendants and siblings of validation individuals and their
descendants do not count towards this minimum number. The default
minimum number is 1.

\textbf{!ValSolutions \textless name file\textgreater{}} The qualifier
!ValSolutions is optional and specifies the solutions file to use for
validation. !ValSolutions is not needed in most cases, but it can be
used to validate an index of all EBV in the evaluation, for example. The
default is the default MiXBLUP file with genetic effect solutions
(Solani.out or Solani.txt).

\textbf{!YieldDev} The qualifier !YieldDev is mandatory to enable
adjusted-phenotype validation.

\textbf{!HeritabFile \textless file with direct heritability per
trait\textgreater{}} The qualifier !HeritabFile is required for
adjusted-phenotype validation to get realized prediction accuracies. The
file contains a line for each trait. A line consists of trait name
(case-sensitive) and direct heritability.

\subsubsection{12.7. Validation statistics}\label{Vali09}

Validation statistics can be found in MiXBLUP.log. The first table gives
an overview of the amount of data available for validation. The example
(Figure 1) is taken from a parent validation study, using dams, of a
maternal genetic effect that was only fitted for trait1.

\pandocbounded{\includegraphics[keepaspectratio]{https://raw.githubusercontent.com/JantN1966/MasterManual/main/Images/ValStud01.jpg}}

\textbf{Figure 1.} Overview of data available for validation of a
maternal genetic effect

{\def\LTcaptype{none} % do not increment counter
\begin{longtable}[]{@{}
  >{\raggedright\arraybackslash}p{(\linewidth - 2\tabcolsep) * \real{0.5000}}
  >{\raggedright\arraybackslash}p{(\linewidth - 2\tabcolsep) * \real{0.5000}}@{}}
\toprule\noalign{}
\begin{minipage}[b]{\linewidth}\raggedright
label
\end{minipage} & \begin{minipage}[b]{\linewidth}\raggedright
Field name of trait as specified in MiXValidate instruction file
\end{minipage} \\
\midrule\noalign{}
\endhead
\bottomrule\noalign{}
\endlastfoot
N & Number of validation individuals \\
nExcl & Number of validation individuals excluded from calculation of
validation statistics because of insufficient number of data records \\
nMin & Lowest number of data records for a validation individual \\
nMax & Highest number of data records for a validation individual \\
Avg & Average number of data records per validation individual \\
\end{longtable}
}

For individual validation, the number of data records is always on the
individual itself only. For parent validation, it is the number of data
records across progeny. In the example, validation individuals consisted
of all dams in a generation. Of the 800 dams, 800 -- 286 = 514
individuals had progeny appearing in the maternal genetic effect field
of a data record, so progeny that were a dam, too. The number of data
record per validation individual varied from 0 to 72 with an average of
12. A validation study using the same validation individuals and the
same data records, but for the direct genetic effect gives a different
table (Figure 2).

\pandocbounded{\includegraphics[keepaspectratio]{https://raw.githubusercontent.com/JantN1966/MasterManual/main/Images/ValStud02.jpg}}

\textbf{Figure 2.} Overview of data available for validation of a direct
genetic effect Validation statistics are presented in two or three
tables, depending on the validation effect. The first table shows LR
validation statistics only for included validation individuals. The last
table show LR validation statistics for all validation individuals
(Figure 3 and Figure 4).

\pandocbounded{\includegraphics[keepaspectratio]{https://raw.githubusercontent.com/JantN1966/MasterManual/main/Images/ValStud03.jpg}}

\textbf{Figure 3.} Validation statistics for a maternal genetic effect
The explanation of columns in these tables is as follows.

{\def\LTcaptype{none} % do not increment counter
\begin{longtable}[]{@{}
  >{\raggedright\arraybackslash}p{(\linewidth - 2\tabcolsep) * \real{0.5000}}
  >{\raggedright\arraybackslash}p{(\linewidth - 2\tabcolsep) * \real{0.5000}}@{}}
\toprule\noalign{}
\begin{minipage}[b]{\linewidth}\raggedright
Solution
\end{minipage} & \begin{minipage}[b]{\linewidth}\raggedright
Combination of validation effect and trait name
\end{minipage} \\
\midrule\noalign{}
\endhead
\bottomrule\noalign{}
\endlastfoot
n & Number of validation individuals included in calculations \\
intercept & Intercept of regressing solutions from full data (W) on
solutions of partial data (P) \\
slope W on P & Slope of regressing solutions from full data (W) on
solutions of partial data (P) \\
correlation & Pearson correlation between solutions from full data and
solutions of partial data \\
level bias & Average of solutions from full data minus average of
solutions from partial data \\
pred accuracy & Covariance between solutions from full data and
solutions of partial data divided by genetic standard deviation among
validation individuals. \\
\end{longtable}
}

The genetic standard deviation among validation individuals is estimated
as the square root of (1-F\textsubscript{VI})*var\textsubscript{A}, with
F\textsubscript{VI} the average inbreeding coefficient among validation
animals and var\textsubscript{A} the additive genetic variance that
applies to the combination of validation effect and trait.

For validation studies of a direct genetic effect, the second table
shows the adjusted-phenotype validation statistics (Figure 4).

\pandocbounded{\includegraphics[keepaspectratio]{https://raw.githubusercontent.com/JantN1966/MasterManual/main/Images/ValStud04.jpg}}

\textbf{Figure 4.} Validation statistics for a direct genetic effect.

The explanation of columns in the second table is as follows.

{\def\LTcaptype{none} % do not increment counter
\begin{longtable}[]{@{}
  >{\raggedright\arraybackslash}p{(\linewidth - 2\tabcolsep) * \real{0.5000}}
  >{\raggedright\arraybackslash}p{(\linewidth - 2\tabcolsep) * \real{0.5000}}@{}}
\toprule\noalign{}
\begin{minipage}[b]{\linewidth}\raggedright
Solution
\end{minipage} & \begin{minipage}[b]{\linewidth}\raggedright
Combination of validation effect and trait name
\end{minipage} \\
\midrule\noalign{}
\endhead
\bottomrule\noalign{}
\endlastfoot
n & Number of validation individuals included in calculations \\
intercept & Intercept of weighted regression of adjusted phenotypes from
full data (W) on solutions of partial data (P) \\
slope W on P & Slope of weighted regression of adjusted phenotypes from
full data (W) on solutions of partial data (P) \\
correlation & Weighted Pearson correlation between adjusted phenotypes
from full data and solutions of partial data \\
level bias & Average of adjusted phenotypes from full data minus average
of solutions from partial data \\
pred accuracy & Weighted correlation between adjusted phenotypes from
full data and solutions of partial data divided by the expected
prediction accuracy. \\
\end{longtable}
}

The expected prediction accuracy is calculated as the square root of
N\textsubscript{prog} * h\textsuperscript{2} /(N\textsubscript{prog} *
h\textsuperscript{2} + 4 -- h\textsuperscript{2}) for parent validation
or the square root of h\textsuperscript{2} for individual validation,
where N\textsubscript{prog} is the average number of progeny records and
h\textsubscript{2} is the heritability of the combination of validation
effect and trait, read from the file specified with !HeritabFile.

Adjusted phenotypes are obtained from the evaluation of the full data
set and are the sum of genetic effects and the residual effect. The
expectation of intercept and level bias is zero for LR validation if a
suitable genetic base is used. A suitable genetic base consists of a
number of high-reliability males at least two generation before the
validation individuals. This can be done with !BaseAnimalsZero
\textless file\textgreater{} in the SOLVING section. The expectation for
slope is 1.0 for LR validation. For adjusted-phenotype validation, it is
0.5 for parent validation and 1.0 for individual validation. The
realized prediction accuracy is an estimate of the average accuracy of
estimated breeding values of validation individuals from the partial
data. Note that this different from the calculated accuracy of estimated
breeding values from the inverse coefficient matrix or approximations
thereof (!Reliability). The latter is a relative measure of the amount
of information available to estimate the breeding value assuming that
the model and variance components are appropriate and that there is no
confounding of effects. Note that this evaluation is based on the
assumption that the accuracy (or reliability) of the EBV of validation
individuals from the partial dataset is non-zero and more or less the
same across validation animals. The realized prediction accuracy is an
estimate of the average of this accuracy across validation individuals.
It is advised not to include individuals with a much lower accuracy of
the partial EBV in the group of validation individuals. Avoid
parent-offspring pairs among the validation individuals for this reason.

\subsubsection{12.8. Associated output files}\label{Vali10}

{\def\LTcaptype{none} % do not increment counter
\begin{longtable}[]{@{}
  >{\raggedright\arraybackslash}p{(\linewidth - 2\tabcolsep) * \real{0.5000}}
  >{\raggedright\arraybackslash}p{(\linewidth - 2\tabcolsep) * \real{0.5000}}@{}}
\toprule\noalign{}
\begin{minipage}[b]{\linewidth}\raggedright
Output file
\end{minipage} & \begin{minipage}[b]{\linewidth}\raggedright
Description
\end{minipage} \\
\midrule\noalign{}
\endhead
\bottomrule\noalign{}
\endlastfoot
MiXBLUP.log & Contains the summary tables of validation statistics \\
LogValidationData.txt & For each validation individual by trait: number
of data records, EBV partial \& EBV whole \\
\end{longtable}
}

\hyperref[Tabl01]{Back to Table of Contents}

\newpage

\subsection{13. Indirect prediction}\label{Indi01}

For an individual that is genotyped, but does not have a phenotype or
progeny with a phenotype yet, each genomic estimated breeding value
(GEBV) is just a function of SNP marker genotypes, SNP marker effects
and the average polygenic estimated breeding value of its parents. There
is limited benefit of solving equations for these animals iteratively.
These GEBV can be predicted indirectly as a separate step after the
routine genetic evaluation.

\subsubsection{13.1. General}\label{Indi02}

In some breeding programs, there is a constant influx of newly genotyped
young individuals, but there are only a limited number of routine
genetic evaluations in a year. Indirect prediction of their GEBV avoids
having to run the full genetic evaluation multiple times. In other
breeding programs, there may be a large number of genotyped individuals
that will never have a phenotype or progeny with a phenotype. Using
indirect prediction for these individuals reduces the size of the
genetic evaluation and saves computing time and resources. Indirect
prediction is a feature of the hpblup solver.

\subsubsection{13.2. Indirectly predicting genomic estimated breeding
values}\label{Indi03}

\paragraph{13.2.1. General}\label{Indi04}

If indirect prediction is used to predict GEBV of individuals without
phenotypes or progeny with phenotypes, then these individuals are not
included in the routine genetic evaluation. After the routine genetic
evaluation is completed, the user provides a pedigree file and a
genotype file that contain individuals in the evaluation and individuals
excluded from the routine genetic evaluation. MiXBLUP identifies the
latter group of individuals in the pedigree and genotype file and
predicts GEBV for these individuals indirectly.

\paragraph{13.2.2. Supported evaluations}\label{Indi05}

Indirectly predicting GEBV is only supported when using the hpblup
solver and either ssSNPBLUP or ssGBLUP applying the Ta or Tac
decomposition of G. MiXBLUP will verify that the routine genetic
evaluation is suitable for indirect prediction. If the routine genetic
evaluation is unsuitable, it will result in a fatal error.

\paragraph{13.2.3. Syntax}\label{Indi06}

The recommended syntax for calling MiXBLUP for indirect prediction is:

\begin{quote}
MiXBLUP pred -p \textless pedigree file\textgreater{} -g
\textless genotype file\textgreater{}
\end{quote}

The old syntax is still supported:

\begin{quote}
MiXPred.exe -p \textless pedigree file\textgreater{} -g
\textless genotype file\textgreater{}
\end{quote}

\paragraph{13.2.4. Output files}\label{Indi07}

The file with the solutions of the routine genetic evaluation is renamed
to Solani\_old.out. The solutions of routine genetic evaluation and
those of indirect prediction are written to Solani.out.

\hyperref[Tabl01]{Back to Table of Contents}

\newpage

\subsection{14. Running MiXBLUP}\label{Runn01}

\textbf{This chapter describes practical issues when analyzing data with
MiXBLUP.}

\subsubsection{14.1. Starting a MiXBLUP evaluation}\label{Runn02}

Solving mixed model equations using MiXBLUP involves execution of
several programs. The main executable is MiXBLUP.exe. This is the parser
and it either calls the MiX99 executables dataprocessor.exe, solver.exe
or reliabilities.exe, or it calls one of the hpblup executables
hpblup.exe or hpblup\_gpu.exe. For calculation of a genomic relationship
matrix, MiXBLUP calls calc\_grm.exe.

The recommended way to call MiXBLUP for breeding value estimation is:

\begin{quote}
MiXBLUP.exe blup -i \textless name instruction file\textgreater{}
\end{quote}

Calling MiXBLUP without any arguments prints a help message (also
available with the ``--help'' flag). Each flag under the new syntax has
a long and a short version, long versions are always preceded by two
dashes and short versions by a single dash. The user may choose the long
versions for readability and documentation or the short ones for brevity
and reduced typing.

Table with flag options:

{\def\LTcaptype{none} % do not increment counter
\begin{longtable}[]{@{}lllll@{}}
\toprule\noalign{}
Purpose & Command & Long flag & Short flag & Usage \\
\midrule\noalign{}
\endhead
\bottomrule\noalign{}
\endlastfoot
EBV calculation & blup & --instruction & -i & Instruction file \\
Reliability & rel & & & \\
Validation & val & & & \\
Indirect Prediction & pred & --pedigree & -p & Pedigree file \\
& & --genotype & -g & Genotype file \\
& any & --debug & -D & Debug options (s t m l) \\
\end{longtable}
}

The old syntax to call MiXBLUP is deprecated, but still available:

\begin{quote}
MiXBLUP.exe \textless name instruction file\textgreater{}
\end{quote}

\subsubsection{14.2. Choosing a breeding value evaluation or a
reliability calculation}\label{Runn03}

MiXBLUP either estimates breeding values, using either the default or
the hpblup solver, or calculates approximate reliabilities, using
reliabilities.exe. The type of analysis is controlled with the
!RELIABILITY qualifier in the SOLVING section in the instruction file.
If it is specified, a reliabilities calculation is started. See Chapter
\hyperref[Reli01]{9.} for the additional changes in the instruction file
when a reliabilities analysis is required. By default, a breeding value
analysis is started. The hpblup solver can be called with !hpblup in the
SOLVING section (see Chapter \hyperref[Cont01]{8.}).

\subsubsection{14.3. A breeding value analysis with previous solutions
as starting values}\label{Runn04}

In case of large evaluations of breeding values, there may be a
substantial saving in time to convergence of 10-30\% by using the
previous solutions as starting values for the current evaluation. This
is activated by specifying the !RESTART qualifier or the
!STARTVAL\_CHECK qualifier in the SOLVING section. It does not have an
effect on a reliabilities analysis. For the MiX99 solver, the only
additional file necessary for using previous solutions is the Solunf
file. If !RESTART is specified, MiXBLUP renames the file Solunf to
Solold. If the file Solold is present, the preprocessor dataprocessor
will create a file Solvec. This file will be used to initialise the
solution vector of the mixed-model equations before the start of the
iteration process (in solver). If any of the effects in the statistical
model has been defined with field type A (alphanumerical labels), then
the file Code.inp of the previous analysis must be present, too. The
file Code\_index.inp of a previous analysis is not used for a restart.
For the hpblup solver, it is better to use !STARTVAL\_CHECK instead of
!RESTART. Additional files needed are startval\_mixblup.new (in case the
pedigree contains multiple base populations) or solutions\_mixblup.dat
(pedigree contains a single base population). MiXBLUP renames these
files to startval\_mixblup.dat. This file is read by the hpblup solver
for initializing the solutions vector. The file hpCodes.bin is also
needed as it contains the key between original and coded labels.

\subsubsection{14.4. Monitoring and checking the process}\label{Runn05}

When developing new analyses, it may be useful to monitor the progress
of the analysis. This can be specified with ``-D s'' on the command
line, for example

\begin{quote}
MiXBLUP blup -i TestRun.inp -D s
\end{quote}

The output of the --debug or -D option is written to the screen log.
Useful debug options:

\textbf{-D s} The debug option s provides more detail on the stage of
the evaluation.

\textbf{-D t} The debug option t provides more detail on the run time of
the various stages of the evaluation.

\textbf{-D m} The debug option m provides more detail on the memory use
of the various stages of the evaluation.

\textbf{-D l} The debug option l (lowercase L) provides more detail on
the current host, the license host, the type of license, the path to the
license, and the end date of the license.

After the run is finished, it is worth to look through the various
log-files. For the MiX99 solver, check ERMcalc\_grm.log,
dataprocessor.log and solver.log. For the hpblup solver, check
ERMcalc\_grm.log, hpblup.log and log\_hpblup.dat. In MiXBLUP.log some
information is given about pre- and post-processing of the data. It also
lists error messages, if any. If all programs have run successfully, it
is worth to check solver.log for the MiX99 solver or hpblup.log and
convergence.dat for the hpblup solver, to see how the convergence was
reached. In cases with poor convergence, it will give a warning and some
model checking may be appropriate. When reliabilities are calculated,
one can check the reliabilities.log, or reliabilities\_direct.log and
reliabilities\_indirect.log when reliabilities are calculated for both
direct and indirect genetic effects, such as maternal genetic effects.

\subsubsection{14.5. Interrupting a process of the kernel}\label{Runn06}

The MiX99 solver can be interrupted by placing an empty file with file
name \textbf{STOP} in the folder of the analysis. The hpblup solver can
be interrupted by placing an empty file with file name \textbf{stopiter}
or \textbf{STOP} in the folder of the analysis. After every iteration,
both solvers check whether the stop file is present. If so, it will
start producing the output files as if convergence had been attained,
instead of the next iteration, and it will stop afterwards.

\hyperref[Tabl01]{Back to Table of Contents}

\newpage

\subsection{15. Decoding any file with coded class effect
labels}\label{Deco01}

\textbf{MiXBLUP produces quite a range of internal files with
information that might be useful for the user in specific cases.
Internal files contain coded labels of class effects and need to be
decoded to be useful.}

\subsubsection{15.1. General}\label{Deco02}

The solver requires that labels of class effect levels are coded 1 to N.
In order to make coding and decoding as efficiently as possible, the key
to code and decode labels is stored in a binary file. Therefore MiXBLUP
has tools to decode files instead of to decode files manually.

\subsubsection{15.2. Decoding coded labels}\label{Deco03}

\paragraph{15.2.1. General}\label{Deco04}

The decoding tool can be used for either the MiX99 or the hpblup solver.
It will detect which solver has been used. If both types of key exist in
the folder, the user is asked which key to use. If neither key exists,
an error is given. The tool can be used to decode individual coded
labels, a file with coded labels created by the user or an internal file
created by MiXBLUP.

\paragraph{15.2.2. Syntax}\label{Deco05}

The tool to decode coded labels is called from the command line as:

\begin{quote}
Coded2Original.exe
\end{quote}

If the tool identifies hpblup as the solver used, it will ask for the
field name in the genetic evaluation of the coded labels to decode.

\begin{quote}
hpblup: the file hpCodes.bin wil be used for decoding What is the field
name of IDs to be decoded (not case-sensitive)?
\end{quote}

The solver hpblup has a separate key for each class effect. Genetic
class effect levels (i.e.~IDs) of indirect genetic effects are coded
using the key of the direct genetic effect. The MiX99 solver just has a
single key for all alphanumerical class effects.

\begin{quote}
MiX99: the file code.inp wil be used for decoding
\end{quote}

\paragraph{15.2.3. Decoding individual coded class effect
labels}\label{Deco06}

For decoding (one or a few) individual class effect labels, answer `1'
to the question

\begin{quote}
Do you want to decode (1) individual IDs or (2) a file with coded IDs?
\end{quote}

The user then specifies the codes to decode one by one and closes by
entering code `0', which closes the tool.

\paragraph{15.2.4. Decoding a file that contains coded class effect
labels}\label{Deco07}

For decoding (one or a few) individual class effect labels, answer `2'
to the question

\begin{quote}
Do you want to decode (1) individual IDs or (2) a file with coded IDs?
\end{quote}

The user is then asked three more questions:

\begin{quote}
Enter name of file with IDs to convert:\\
Enter field number with coded IDs to decode:\\
Enter name of new file with decoded IDs:\\
\end{quote}

The new file is the same as the existing file, except for the coded
label replaced with the original label.

\hyperref[Tabl01]{Back to Table of Contents}

\newpage
